\chapter{Status of the Open Problems -- Detailed Derivations}

\label{app:status-detailed}



This file contains the detailed derivations, conditional results,

no-gos, and bounded results that were previously interleaved in

appH-microscopic-proofs.tex. Each block is labeled for cross-reference.





\section{Problem A: matrix-to-modular map}

\label{sec:detailed-Problem A: matrix-to-modular map}



\paragraph{Ledger label.} \label{detail:sec-A-paragraphledgerlabel-0}

Status: OPEN. Priority: CRITICAL PATH for the matrix-derived label.

Completion: PROGRAMME-CONDITIONAL. The matrix-to-modular map remains

unproved. A scope-limited Presentation-B model may retain

Assumption~\ref{ass:dictionary} and claims no derivation from the

matrix model.



For a constant-\(\Theta\) Moyal realization, the commutator of a

coordinate symbol \(x_\mu\) with a field is exact, with no

higher-derivative correction; see

Remark~\ref{rem:ch3-moyal-check}. That result is conditional on this

realization and is not a general expansion for fluctuating

\(\hat X_\mu\). The remark also records a metric-matching obstruction

for a real constant \(\Theta\). Neither result supplies the map from

\((\hat X,\Phi)\) to the causal-diamond modular Hamiltonian \(H_D\)

or the mutual-information functional

\(U_{\mathrm{ent}}[I(A:B)]\) before the composite metric is defined.







\section{Problem B: replica Ward identity}

\label{sec:detailed-Problem B: replica Ward identity}



\paragraph{Ledger label.} \label{detail:sec-B-paragraphledgerlabel-1}

Status: OPEN. Priority: CRITICAL PATH. Completion:

PROGRAMME-CONDITIONAL. The first-law identity is established, but

the stationarity condition for the tensor block is not derived.



For a differentiable family of normalized reduced states on a fixed

diamond algebra, the first variation of the entropy obeys the standard

first-law identity

\[

\delta S_D

=

\operatorname{Tr}(\delta\rho_D H_D^{(0)})

=

\delta\langle H_D^{(0)}\rangle,

\qquad

H_D^{(0)}=-\ln\rho_D^{(0)},

\]

where the modular Hamiltonian is held fixed at the reference state.

Under the usual path-integral assumptions, the replica construction

expresses this variation as a response of the microscopic action

insertion on the branched cover. These identities do not imply

\(\delta S_D=0\).





\paragraph{Which first law is established (conditional; not closure).} \label{detail:sec-B-paragraphwhichfirstlawisestabl-2}

The identity above is a statement on a fixed diamond algebra: only

the state varies, and \(H_D^{(0)}\) is held fixed.  As stated in the

ledger, it does not imply \(\delta S_D=0\).  The first law used in

Chapter~\ref{ch:entanglement-equilibrium},

Equation~\eqref{eq:first-law}, also contains the area term

\(\delta A/(4G_{\mathrm{eff}})\), which arises when the diamond itself

varies with the emergent metric.  That is a variation of the algebra,

not of the state on a fixed algebra, so the fixed-algebra identity does

not establish it.  Obtaining it requires the map between variations of

the scaffold diamond and of the emergent diamond noted above.  In

Chapter~\ref{ch:entanglement-equilibrium} the coefficient of

\(\delta A\) is only named \(1/(4G_{\mathrm{eff}})\); its identification

with the heat-kernel \(G_{\mathrm{eff}}\) is deferred to

Chapter~\ref{ch:heat-kernel} and depends on the open Problem~F.  The

form of Assumption~\ref{ass:EE} stated in

Chapter~\ref{ch:entanglement-equilibrium},

Equation~\eqref{eq:first-law-zero}, is an assumption that goes beyond

the fixed-algebra identity.  This paragraph records that distinction;

it does not derive the area term or the stationarity condition, and

Problem~B remains open.





\paragraph{Replica cone coefficient of the Poisson-limit operators and the area term of the first law (conditional; not closure).} \label{detail:sec-B-paragraphreplicaconecoefficien-3}

Keep the setting of the Presentation~A Poisson-limit paragraphs of Section~\ref{sec:status-F}: four matrix coordinates, the Poisson limit, the Euclidean section \(\eta=\delta\), the scalar operator \(\mathcal O_s\), the connection operator \(\mathcal O_v\) acting on the internal index of \(a^\mu\) in background gauge, the ghost operator of weight \(-2\), \(G=\rho\,g_\rho\) with \(\sqrt{|G|}=\rho^2\), Assumption~\ref{ass:emergent-op}, and one proper-time limit \(s_{\min,s}=\Lambda_{\mathrm{UV},s}^{-2}\) per sector, shared by the effective action and the replica entropy. Only the \(\Lambda_{\mathrm{UV},s}^2\) term is compared, for constant \(\Lambda_{\mathrm{UV},s}\); neither a value of \(G_{\mathrm{eff}}\) nor the relation \(S=A/4G\) is assumed.



\emph{Replica on the matrix chart.} Let \(Z_n\) be the partition function on the \(n\)-sheeted cover of the chart \(x\), branched over a two-surface \(\Sigma\), with every background field (\(\Theta\), hence \(G\), \(g_\rho\), \(\rho\), the connections and the endomorphisms) pulled back. Then \(G\) and \(g_\rho\) are conical with the same angle \(2\pi n\) and \(\rho\) is a regular density of the chart: the cone is on \(G\) and \(g_\rho\) together. For a Laplace-type sector whose connection and endomorphism stay bounded on the cover, the leading surface term of the heat trace is the flat-cone term~\cite{Solodukhin2011},

\[

  \operatorname{Tr}e^{-s\mathcal O}\big|_n-n\operatorname{Tr}e^{-s\mathcal O}\big|_1

  =\operatorname{tr}\mathbb I\;\frac1{12}\Bigl(\frac1n-n\Bigr)\frac{A_G(\Sigma)}{4\pi s}+O(s^0),

\]

independent of the connection and of the endomorphism. The fibre index of \(\mathcal O_v\) is internal, so the fibre is not rotated around \(\Sigma\) by a Levi-Civita holonomy. With \(S=(1-n\partial_n)(\ln Z_n-n\ln Z_1)|_{n=1}\) and \(\ln Z_n=\tfrac12\sum_sw_s\int\mathrm{d}s\,s^{-1}\operatorname{Tr}e^{-s\mathcal O_s}\),

\[

  S=\frac{A_G(\Sigma)}{48\pi}\sum_sw_s\operatorname{tr}\mathbb I_s\,\Lambda_{\mathrm{UV},s}^2

\]

(the \(\mu_s^2\) terms of Theorem~\ref{thm:Geff} subtracted). The same limits in Equation~\eqref{eq:universal-R} give, for the coefficient of \(\int\sqrt{G}\,R_G\) when the endomorphisms carry no explicit curvature, \(1/(16\pi G)=\sum_sw_s\operatorname{tr}\mathbb I_s\,\Lambda_{\mathrm{UV},s}^2/(192\pi^2)\). Hence \(S=A_G(\Sigma)/(4G)\) with that \(G\), in the variables \((G,\rho,\partial\Theta)\) of the contact-term paragraph below: for these operators the coefficient of \(\delta A\) in Equation~\eqref{eq:first-law} is the \(\LUV^2\) coefficient of Theorem~\ref{thm:Geff}, namely the coefficient \(A=1/(12\pi)\) of Equation~\eqref{eq:Geff-structure} per unit weighted fibre trace, as in the sector-sum paragraph of Section~\ref{sec:status-F}, and it is determined by the operators and the cutoff. This holds for the replica and variables just stated; the scope of the whole block is given at the end of the second paragraph below.



\emph{Contact term.} If an endomorphism is \(E_s=E_{0,s}+Q_sR_G\) and \(R_G\) is continued onto the cone with its conical delta, the heat trace gains \(-(1-n)\operatorname{tr}Q_s\,A_G/(4\pi s)\) and the entropy changes by \(-\tfrac{A_G}{8\pi}\sum_sw_s\operatorname{tr}Q_s\Lambda_{\mathrm{UV},s}^2=-2\pi A_G\sum_sw_s\operatorname{tr}Q_s\langle\phi^2\rangle_s\), \(\langle\phi^2\rangle_s=\Lambda_{\mathrm{UV},s}^2/(16\pi^2)\). This is the contact term of Susskind and Uglum~\cite{SusskindUglum1994}, \(S_{\mathrm{BH}}=S_{\mathrm{ent}}-Q\) in Equations~(322)--(323) of~\cite{Solodukhin2011}, with \(S_{\mathrm{BH}}\) built from the coefficient \(\operatorname{tr}(\tfrac16-Q_s)\) of Theorem~\ref{thm:Geff}. With \((G,\rho,\partial\Theta)\) as independent variables, the scalar endomorphism \(E_s=-\tfrac12\Box_G\ln\rho+\tfrac14|\partial\ln\rho|_G^2\), the connection endomorphism with \(\operatorname{tr}E_v=4E_s+\operatorname{tr}(B_\sigma B^\sigma)\), \(\operatorname{tr}B_\sigma=0\), \(\operatorname{tr}(B_\sigma B^\sigma)=-X\), and the ghost operator \(-\mathcal O_s\) all have \(Q=0\) and are bounded on the pulled-back cover, so the contact term vanishes sector by sector.



\emph{Dependence on the independent variables.} The Weitzenb\"ock identity quoted in the paragraph on the torsion invariant gives \(E_s=\tfrac14R_G+\tfrac1{16}T_{\lambda\mu\nu}T^{\lambda\mu\nu}+\tfrac18T_{\lambda\mu\nu}T^{\mu\lambda\nu}\) and \(\operatorname{tr}(B_\sigma B^\sigma)=-T_{\lambda\mu\nu}T^{\lambda\mu\nu}\) (the first checked at random two-jets of a closed \(\omega\), for both signatures). With the coframe torsion in place of \(\partial\Theta\) as independent variable, the scalar operator has \(Q_s=\tfrac14\), and its coefficient \(\mathcal A_{1,R,s}\) is \(\tfrac16-\tfrac14=-\tfrac1{12}\) per component instead of \(\tfrac16\), of opposite sign. The coefficient \(\mathcal A_{1,R,s}\) of Theorem~\ref{thm:Geff} is therefore not a function of these operators alone. With \(N_{\mathrm{tot}}=N_s+2\) as in the Poisson-limit paragraph, the weighted coefficient \(c\) of \(\sqrt G\,R_G\) in \(\operatorname{tr}(R_G/6-E)\) is:

\begin{center}

\begin{tabular}{@{}lcc@{}}

\toprule

sector (weight) & \(c\) in \((G,\rho,\partial\Theta)\) & \(c\) in \((G,\text{torsion})\)\\

\midrule

\(N_s\) scalar components & \(N_s/6\) & \(-N_s/12\)\\

connection, 4 components & \(2/3\) & \(-1/3\)\\

ghost, weight \(-2\) & \(-1/3\) & \(1/6\)\\

total & \(N_{\mathrm{tot}}/6\) & \(-N_{\mathrm{tot}}/12\)\\

\bottomrule

\end{tabular}

\end{center}



\emph{Which geometry carries the cone.} The first column is the replica entropy above, \(S=(\LUV^2A_G/8\pi)\,c\), with \(\rho\) regular. The cone on \(G\) can instead be carried by the Poisson structure: for \(\Theta^{12}=1/f(x^1,x^2)\), \(\Theta^{34}=1\), one has \(G_\perp=f^2\delta\) and \(\rho=f\), and \(f=|z|^{n-1}\) makes \(G\) conical of angle \(2\pi n\) with \(\rho\propto r^{n-1}\) singular, \(g_\rho\) singular, and torsion unbounded, \(|T|^2\propto(n-1)^2/u^2\) at proper distance \(u\) from \(\Sigma\). Then \(E_s\) contains \(-\tfrac12\Box_G\ln\rho=\pi(1-n)\delta_\Sigma\), a contact term with \(Q=\tfrac14\), and \(\mathcal O_s\) reduces to a Bessel problem of index \(\sqrt{m^2+((n-1)/2)^2}/n\). With the branch regular at \(m=0\), the first-order coefficient is that of the second column, for the scalar and ghost and for the components of \(a^\mu\) in this block class (a statement conditional on that branch and not used below; the factor \(-\tfrac12\) is also what the local \(a_1\) coefficient with the delta function, \(Q=\tfrac14\), gives). The two choices differ by the factor \(-\tfrac12\) in every sector (for the branch stated). The first is the replica of the same state, since \(\operatorname{Tr}\rho_D^n\) is built on the cover of the same background; the second changes the background with \(n\) and corresponds to a smooth-coframe reading of the induced action. Only the first is used in the identification above.





\paragraph{Gauge parameter and gauge function in the comparison of \(G_T\) with \(G_A\) (conditional; not closure).} \label{detail:sec-B-paragraphgaugeparameterandgaug-4}

Keep the setting, the class and the normalisation of the two preceding paragraphs. Write \(e_i\) for the contribution of a sector of weight \(w_i\) to the bracket \(1-N_{\mathrm{tot}}/24\) of the preceding paragraph, so that \(\sum_iw_ie_i\) is the net eigenvalue on the two modes with \(\delta\rho=0\) and \(c_T=-4\sum_iw_ie_i\); a scalar component has \(e=-\tfrac1{24}\) there and \(-\tfrac18\) on the mode with \(\delta\rho\neq0\). Let \(D_\mu=\{x_\mu,\cdot\}\), anti-self-adjoint on \(L^2(\rho\,\mathrm d^4x)\), \((\mathrm{Gr}\varphi)_\nu=D_\nu\varphi\), and \(\mathfrak J_\nu=D^\mu\Theta_{\mu\nu}\), the Poisson limit of the tree-level equation of motion; for the perturbation of the quadratic-form paragraph \(\mathfrak J_\nu=-\epsilon\cos(p\cdot x)\,(Mu)_\nu\), with \(M=q^2\eta-qq\) the kernel of Remark~\ref{rem:ch3-trace-quadratic}. The background gauge \(\tfrac1{2\alpha}\rho\,(D_\mu a^\mu)^2\), \(\alpha>0\), \(\kappa=1-1/\alpha\), leaves the ghost operator unchanged and replaces \(\mathcal O_v\) by the non-minimal operator \(\mathcal O_\kappa=\mathcal O_v-\kappa\,\mathrm{Gr}\mathrm{Gr}^\dagger\), with principal symbol \(q^2\mathbb I-\kappa\,qq^{T}\); one proper-time cut-off on \(\mathcal O_\kappa\) serves the effective action and the replica, and the normalisation of the gauge-fixing weight is local and carries no two-derivative term.



For every Poisson bivector \(\mathrm{Gr}^\dagger\mathrm{Gr}=\mathcal O_s\) and \(\mathcal O_v\mathrm{Gr}=\mathrm{Gr}\mathcal O_s-\{\mathfrak J_\nu,\cdot\}\): the commutator terms are fixed by \([D_\mu,D_\nu]=\{\Theta_{\mu\nu},\cdot\}\) and the Leibniz property of \(D_\mu\) on the bracket, and the \(\{\Theta_{\mu\nu},\cdot\}\) coupling of \(\mathcal O_v\) cancels them up to \(\mathfrak J\) (checked in exact arithmetic for a nonlinear bivector). Hence for \(\mathfrak J=0\), \(\mathcal O_\kappa\mathrm{Gr}=(1-\kappa)\mathrm{Gr}\mathcal O_s\), the longitudinal sector is the scalar operator with proper time rescaled by \(1-\kappa\), and up to zero modes \(\operatorname{Tr}e^{-s\mathcal O_\kappa}=\operatorname{Tr}e^{-s\mathcal O_v}+\operatorname{Tr}e^{-s(1-\kappa)\mathcal O_s}-\operatorname{Tr}e^{-s\mathcal O_s}\). The integrated \(\LUV^2\) density of the connection sector is then \(\tfrac{3+\alpha}6\,\rho R[g_\rho]+\rho^2X\) on such backgrounds, with the coefficient \(1\) of \(\rho^2X\) independent of \(\alpha\). On the replica cover with constant \(\bar\Theta\) the same intertwining, with the domain at \(\Sigma\) for which the longitudinal sector is the image of the regular scalar domain, gives the area coefficient \(\tfrac{3+\alpha}6\) for the connection sector in the units \(c\) of the table above, equal to the coefficient of \(\rho R[g_\rho]\). This depends on that choice of domain; other choices are not examined.



About constant \(\bar\Theta\), with \(\mathfrak J\neq0\), a numerical second-order evaluation of the heat trace (momentum-space Duhamel expansion, Gauss--Hermite quadrature, contour extraction of the \(p^4\) coefficient; it reproduces the scalar value above, the \(\alpha=1\) values \(\tfrac56\) and \(\tfrac12\) of the connection sector, and the weight \(3+\alpha^2\) of the \(\LUV^4\) term) gives for the connection sector the values, matched by the rational forms, \(\tfrac56-\tfrac{\alpha-1}{24}\) on the two modes with \(\delta\rho=0\) and \(\tfrac12-\tfrac{\alpha-1}8+\tfrac{\alpha^2-1}4\) on the mode with \(\delta\rho\neq0\) (\(\alpha=\tfrac13,\tfrac12,\tfrac23,1,2,3\) in one background, \(\tfrac12,1,2\) in another; rational fits to \(2\times10^{-3}\), not derived; \(\alpha=\tfrac13,\dots,3\) only). The two modes with \(\delta\rho=0\) thus show the \(\alpha\)-dependence of the \(\mathfrak J=0\) result, and the remainder, \(\tfrac{\alpha^2-1}2\rho|\nabla\ln\rho|^2_{g_\rho}\), acts on the mode with \(\delta\rho\neq0\) only. With scalar and ghost sectors independent of \(\alpha\), \(c_T(\alpha)=\tfrac{N_s+\alpha-1}6-\tfrac{11}3\) and the replica coefficient is \(c_A(\alpha)=\tfrac{N_s+\alpha+1}6\), so \(1/G_A-1/G_T=2\LUV^2/\pi\) for the tested range \(\tfrac13\le\alpha\le3\); on backgrounds with \(\mathfrak J=0\) and under the conditions above the \(\alpha\)-dependence cancels in the difference for every \(\alpha>0\). The TT value agrees with the exact lemma on \(\mathfrak J=0\) backgrounds. For \(\alpha\neq1\) the form is not spanned by \(\rho R[g_\rho]\) and \(\rho^2X\); the third-mode to \(\delta\rho=0\) ratio of a term \(c\,\rho R[g_\rho]\) is \(3\), and the net form attains this ratio at \(\alpha=3\), where the gap is unchanged, so that ratio does not diagnose the coefficient of \(\rho^2X\) outside the Feynman gauge. A gauge function \(D_\mu a^\mu+c\,\mathfrak J_\mu a^\mu\) (\(\alpha=1\), ghost operator \(-D_\mu D_\mu-c\mathfrak J_\mu D_\mu\)) adds \(c^2\rho^2\mathfrak J_\nu\mathfrak J_\nu\) to the density of the connection and ghost sectors together at quadratic order about \(\bar\Theta\), that is \(c^2q^2/(2p_\eta^2)\) to the eigenvalue on all three modes (two backgrounds, \(c=\tfrac12,1,2\)); this lowers \(c_T\) at fixed \(c_A\) and depends on the direction of \(p\) unless \(\bar\Theta\) is isotropic. Gauge functions with other first-derivative vectors are not examined.





\paragraph{Fermion replica gluing for \(\operatorname{Tr}\rho^n\) and the supertrace character (conditional; not closure).} \label{detail:sec-B-paragraphfermionreplicagluingf-5}

Let the state of the Poisson-limit fermions on the wedge be quasi-free, and let its modular flow be the rotation by \(2\pi t\) of the transverse plane of \(\Sigma\), acting on every field through its representation of the rotation group (for integer spin, the identity at integer \(n\)). The ungraded trace \(\operatorname{Tr}\rho^n\) then glues a Grassmann-odd field with the rotation by \(2\pi n\) and one sign from the cyclicity of the trace: in the polar frame the spinor is antiperiodic under the angle \(2\pi n\), and the Cartesian components of transverse spin \(\sigma=\pm\tfrac12\) acquire the monodromy \(-e^{2\pi\mathrm in\sigma}\), which is \((-1)^{n+1}\) for integer \(n\)~\cite{CasiniHuerta2009,FursaevMiele1997}. Diagonalising the replica matrix gives the twist phases \(e^{2\pi\mathrm ik/n}\), \(k=-(n-1)/2,\dots,(n-1)/2\)~\cite{CasiniHuerta2009}. For a block of two components with twist \(a\), \(|a|<\tfrac12\), and the less singular self-adjoint extension of the first-order operator \(\gamma^\mu D_\mu\), whose square is \(\mathcal O_f\), the conical heat-trace constant is \(a^2+\tfrac16\): the componentwise-regular value \(B_2(|a|)\), \(B_2(x)=x^2-x+\tfrac16\), plus \(|a|\) from the reversed sign of the one Bessel order that the first-order operator fixes. This extension is the one giving an analytic continuation: the componentwise-regular domain gives \(\sum_k|a_k|\) equal to \(n/4\) or \((n^2-1)/(4n)\) for even or odd \(n\), a parity-dependent and non-analytic \(\operatorname{Tr}\rho^n\). With the less singular extension, \(\sum_ka_k^2=(n^2-1)/(12n)\), analytic in \(n\), equal to the zeta-function value of the polar-frame cone on the branch \(n<1\), with \((1-n\partial_n)\) equal to \(-\tfrac16\) per block at \(n=1\). The area coefficient of a Dirac field (two blocks, weight \(-1\)) is therefore \(+\tfrac13\), as used above for antiperiodic gluing.



The unrotated periodic gluing above has the phase \(+1\) for all \(n\): its twists are those of \(\operatorname{Tr}\rho^n\) for odd \(n\) and those of the graded trace \(\operatorname{Tr}[(-1)^F\rho^n]\) for even \(n\), and with the componentwise-regular domain it gives \(-2\) times the Dirac-domain value at odd \(n\). It is not \(\operatorname{Tr}\rho^n\), and for the stated extension its values \(-\tfrac23\) and \(4-2n_D\) do not apply to \(\operatorname{Tr}\rho^n\). The graded trace has the twist \(a=\tfrac12\) at \(n=1\) and is not normalised there. The rotation by \(2\pi n\) is the identity on the vector index for integer \(n\), so the integer-\(n\) traces of the connection sector are those of the unrotated gluing; if the continuation from integer \(n\) is unique, the rotated branches \(3-n_D\) and \(5-n_D\) above do not reproduce them and are not continuations of \(\operatorname{Tr}\rho^n\). With \(c_A=\tfrac{N_s+\alpha+1}6+\tfrac{n_D}3\) and \(c_T\) as above, \(c_A-c_T=4(1-\tfrac{n_D}2)+n_D\), the sum of the weight of \(\rho^2X\) and the excess \(+1\) per Dirac field of the cone value \(+\tfrac13\) over its coefficient \(-\tfrac23\) of \(\rho R[g_\rho]\). For \(n_D=2\), \(N_s=6\), \(\alpha=1\): \(c_T=0\), \(c_A=2\).



In the character \(\mathsf Q_{10}\) of~\cite{MantaSteinackerTran2024}, Eq.~(3.7), with the flux along the four Poisson directions, the vector, six inert scalars, the ghosts and \(n_D\) Dirac fields in place of the sixteen-component spinor, the coefficient of order \(\alpha^0\) is \(N_{\mathrm{tot}}=8-4n_D\) and that of order \(\alpha^2\) is \(-b\,(x^2+y^2)\) with \(b=1-n_D/2\), where \(\pm\mathrm ix,\pm\mathrm iy\) are the eigenvalues of the flux. These are the weighted component count and the weighted index of the preceding paragraphs, so that \(c_T=N_{\mathrm{tot}}/6-4b\) vanishes at \(n_D=2\), \(N_s=6\) because the orders \(\alpha^0\) and \(\alpha^2\) of the character vanish. At \(n_D=2\) the order \(\alpha^4\) is \((x^2-y^2)^2/16\), which vanishes for self-dual or anti-self-dual flux. This concerns the character only, not the \(a_2\) coefficient of the operators and not the area side. These statements are conditional on the quasi-free state, the rotational modular flow, the choice of self-adjoint extension, the uniqueness of the continuation from integer \(n\), the domain of the first-order operator, the \(\LUV^2\) order and the class of the preceding paragraphs; they do not derive the state, the flow, the content, or a value or sign of \(G_{\mathrm{eff}}\), and they do not compute the next order of either side. Problems~B, C, D, F and~G remain open.



\paragraph{Replica prescriptions against published contact-term treatments (conditional; not closure).} \label{detail:sec-B-paragraphreplicaprescriptionsa-6}

Keep the setting, the fixed-state replica and the first-order smooth replica of the preceding paragraph, and write \(c_A\) and \(c_A^{\mathrm{sm}}\) for their area coefficients and \(c_R=N_{\mathrm{tot}}/6\) for the coefficient of \(\rho R[g_\rho]\); \(c_A=2\) and \(c_A^{\mathrm{sm}}=0\) at \(n_D=2\), \(N_s=6\). Put \(V_\mu=\partial_\mu\ln\rho\).



\emph{Scalar and ghost sectors.} For a scalar with a coupling \(\xi R\phi^2\), Herzog and Nishioka add the codimension-two term \(\delta I_E=-2\pi(1-n)\xi\int\phi^2\) (Eq.~(2.19) of~\cite{HerzogNishioka2016}), which cancels the delta of \(\xi R\) on the cone and gives an entropy independent of \(\xi\). With the heat-trace bookkeeping of the contact-term paragraph above, this term shifts the entropy by \(+2\pi\xi A\langle\phi^2\rangle\), equal and opposite to the contact \(-2\pi\xi A\langle\phi^2\rangle\), which is the expectation value of Wald's formula for \(\xi R\phi^2\)~\cite{DonnellyWall2016}, Eq.~(4); Herzog and Nishioka do not discuss Wald's formula, and their term is related to it only through this bookkeeping. The fixed-state replica, whose scalar and ghost operators carry no delta on the cone, therefore has the contact-free coefficient. The smooth replica, in which \(E_s\) contains \(-\tfrac12\Box_G\ln\rho\) with \(Q=\tfrac14\), has the form of the delta-potential entropy of Eqs.~(143)--(144) of~\cite{Solodukhin2011} with \(\xi\) replaced by \(Q\); that review does not accept this entropy as an entanglement entropy, and Donnelly and Wall assign the corresponding term to the generalised entropy as a Wald contact. Neither value is asserted here to be the entanglement entropy of the matrix state: the identifications are by the form of the coefficient, the cited analysis is for a scalar on a flat background, and the operators here have no independent \(\xi\). On the singular-\(\rho\) cone \(\ln\rho=(n-1)\ln|z|\), the operator \(\mathcal O_s\) is unitarily equivalent to \(-\Box+E_s\), with s-wave Bessel index \(\nu_0=|n-1|/(2n)>0\). The regular branch vanishes at \(\Sigma\), so the delta multiplies a function that vanishes there, and the constant \(-\nu_0/2\) gives the entropy coefficient \(\tfrac16-\tfrac14\) per component for \(n<1\) and \(\tfrac16+\tfrac14\) for \(n>1\), non-analytic at \(n=1\); the less singular branch exchanges the two. The smooth value \(-\tfrac1{12}\) is thus the \(n<1\) value of the regular branch, and removing the delta alone does not carry this cone to the contact-free value \(\tfrac16\). The terms of order \(((n-1)/2)^2\) in the indices \(\nu_k\) with \(k\neq0\) depend logarithmically on the tip scale and drop out at first order in \(n-1\).



\emph{Connection sector.} For the Maxwell field in Feynman gauge, Herzog and Nishioka remove the Ricci delta of the vector operator with \(\delta I=-\pi(1-n)\int\delta^{ij}A_iA_j\) (Eq.~(7.2) of~\cite{HerzogNishioka2016}), and note that this term is not gauge invariant. Kabat's contact is that delta: it is present in the black-hole entropy and absent from the bulk-mode entanglement entropy~\cite{Kabat1995}, and Donnelly and Wall identify it with edge modes~\cite{DonnellyWall2016}. Herzog and Nishioka state that their boundary term contributes oppositely to the contact term (Eq.~(5.20) of~\cite{HerzogNishioka2016}). The vector contact \(-1\) that gives \(c_A=1\) is therefore the term that the Herzog--Nishioka term removes. Its use presupposes a \(U(1)\) structure and a Ricci coupling that the operator \(\mathcal O_v\) does not have in the variables \((G,\rho,\partial\Theta)\), where the gauge group is symplectic and \(Q=0\); it remains an input and is not a prescription of this class.



\emph{Teleparallel form.} The identity for the Weitzenb\"ock connection of the paragraph on the torsion invariant gives \(E_s=-\tfrac12\nabla_\mu V^\mu+\tfrac14V_\mu V^\mu=\tfrac14\bigl(R_G+\mathcal T\bigr)+\tfrac14V_\mu V^\mu\) with \(\mathcal T=\tfrac14T_{\lambda\mu\nu}T^{\lambda\mu\nu}+\tfrac12T_{\lambda\mu\nu}T^{\mu\lambda\nu}-V_\mu V^\mu\), and the scalar density is \(\tfrac16R_G-E_s=-\tfrac1{12}R_G-\tfrac14\mathcal T-\tfrac14V_\mu V^\mu\). With \(\mathcal T=-R_G-2\nabla_\mu V^\mu\) the coefficient of \(R_G\) is \(\tfrac16\), which is the value in the quadratic form on the two modes with \(\delta\rho=0\), where \(V_\mu V^\mu\) is of higher order. For scalar content alone \(c_T=c_A\), as stated above, and the two replicas differ by the integrated boundary term \(R_G+\mathcal T=-2\nabla_\mu V^\mu\), the flux of \(\nabla\ln\rho\) through \(\Sigma\).



\emph{Dynamical metric.} In the replica construction of Lewkowycz and Maldacena the entropy is \(-\partial_n\) of the action of a regularised cone minus that of the cone, localised at the tip; the difference between the action of the true \(n\)-solution and that of the regularised cone is a first-order variation about a solution and vanishes by the equations of motion (Eqs.~(3.1)--(3.4) of~\cite{LewkowyczMaldacena2013}). That step requires the background to solve the equations of motion of the full action, of which the stationarity condition of the tensor block is the part not derived here. The regularised cone has the far field of the \(n=1\) solution with the angle identified modulo \(2\pi n\), so that every field other than the metric is regular there. With \(\Theta\) the only field, \(R_G^{(1)}=-2\Box_{\bar G}\,\delta\ln\rho\) excludes such a configuration at first order, as stated above. The only first-order regularisation in this class is the smooth family, whose far field has \(\rho\propto r^{n-1}\) and so differs from that of the cone; the construction of~\cite{LewkowyczMaldacena2013}, and the vanishing of its first bracket, are therefore not reproduced for it. For the smooth replica the area coefficient is \(c_A^{\mathrm{sm}}=-c_R/2\). This is the coefficient of \(\sqrt G\,R_G\) at fixed coframe torsion, so the smooth replica satisfies \(S=A/4G_{\mathrm{ind}}\) in those variables for every content, whereas \(c_A^{\mathrm{sm}}-c_T=(14-N_s-4n_D)/4\), which vanishes only for \(N_s+4n_D=14\). The generalised-entropy structure therefore does not imply \(G_A=G_T\) in this class. At \(n_D=2\), \(N_s=6\), the conditions \(c_R=0\) (\(N_s=4n_D-2\)) and \(c_T=0\) (\(N_s=22-8n_D\)) are independent and meet only there; \(c_A^{\mathrm{sm}}=0=c_T\) is then the absence of the \(\LUV^2\) term in the induced action and carries no information on \(G\), while \(c_A=2\) is the contact-free entanglement coefficient, in the relation, in form, \(S_{\mathrm{BH}}=S_{\mathrm{ent}}-Q\) of Eq.~(322) of~\cite{Solodukhin2011}, where \(Q\) is here the bosonic contact \((N_s+2)/4\) of the smooth reading and \(S_{\mathrm{BH}}\) is built from the induced action.



\emph{Intermediate flux.} For first-order smooth perturbations of constant \(\bar\Theta\) the integral of \(R_G^{(1)}\) over the transverse plane is \(-2\) times the flux of \(\nabla\ln\rho\), so no family of such perturbations joins the two replicas. If the first-order local densities of the pulled-back cone and of the smooth family are superposed with weight \(f\) on the smooth family, then \(c_A(f)=c_A-f(N_s+2)/4\), which is \(2-2f\) at \(n_D=2\), \(N_s=6\); the superposition is an assumption. The indices \(\nu_k=\sqrt{k^2+((n-1)/2)^2}/n=|k+\mathrm i(n-1)/2|/n\) are those of an imaginary twist, so the periodic quadratic dependence of the area coefficient on a real Abelian holonomy is not used.



These statements are for the Poisson-limit operators, the \(\LUV^2\) term, one loop, \(\alpha=1\), and the two replicas stated; they do not derive the replica or the state selected by the matrix model, the stationarity condition, the covariance of the induced action in \(\Theta\), the vector contact, or the tip domain, they treat neither edge modes of the symplectic gauge sector nor fermion edge modes, and they give no value or sign of \(G_{\mathrm{eff}}\). Problems~B, C, D, F and~G remain open.





\paragraph{Cone logarithmic coefficient at \(N_s=6\), \(n_D=2\) on the fixed-state replica (conditional; not closure).} \label{detail:sec-B-paragraphconelogarithmiccoeffi-7}

Keep the setting of the preceding paragraph and the fixed-state replica with antiperiodic Dirac gluing and the less singular extension. Let \(\Sigma\) be the fixed plane of a rotation under which \(\Theta\) is invariant, so that \(\Sigma\) is totally geodesic for \(G\). With unit normals \(n_i\) to \(\Sigma\), write \(R_{ii}=R_{\mu\nu}n^\mu_in^\nu_i\) and \(R_{ijij}=R_{\mu\nu\rho\sigma}n^\mu_in^\nu_jn^\rho_in^\sigma_j\) (sums over \(i,j\)), and \(R_\Sigma\) for the intrinsic scalar curvature of \(\Sigma\), so that \(R_\Sigma=R-2R_{ii}+R_{ijij}\) for totally geodesic \(\Sigma\). The pulled-back background is the same on every sheet, so \(\operatorname{Tr}\rho^n\) is the sum over the twist phases \(e^{2\pi\mathrm ik/n}\) of the single-sheet problem with a flat flavour vortex of flux \(a_k=k/n\) at \(\Sigma\) (integer \(k\) for the bosonic sectors, \(k=-\tfrac{n-1}2,\dots,\tfrac{n-1}2\) for the Dirac fields). For a Laplace-type sector \(-\nabla^2+E\) the \(s^0\) surface term is linear in the dimension-two data. For a component with a periodic phase it is the scalar term of Equation~(70) of~\cite{Solodukhin2011}, \(c_B(n)\operatorname{tr}(\tfrac R6-E)+d_B(n)(R_{ii}-2R_{ijij})\) with \(c_B=\tfrac\pi3\tfrac{1-n^2}n\) and \(d_B=-\tfrac\pi{180}\tfrac{(1-n^2)(1+n^2)}{n^3}\), and it has no term in the bundle curvature, because every periodic twist acts as a common phase and the sum of the Bernoulli polynomials \(B_3(k/n)\) in the linear coefficient vanishes. For a Dirac field the common phase is accompanied by the spin term \(\gamma^{12}\Omega_{12}\) of \((\gamma^\mu D_\mu)^2\), with \(\Omega_{12}=[D_1,D_2]\), whose part supported on \(\Sigma\) is attractive in the component where it is positive; the less singular extension then gives the \(E\)-linear weight \(-\tfrac12c_B\) per component and a term \(-\tfrac12\sum_ka_k^2\operatorname{tr}(\gamma^{12}\Omega_{12})\) with \(\sum_ka_k^2=(n^2-1)/(12n)\) and \(\Omega_{12}\) the bundle curvature in the normal plane (from the exact Landau--Aharonov--Bohm spectrum; the terms odd in \(a\) cancel in the symmetric sum over \(k\)). The remaining geometric part of the Dirac term is fixed at first order in \(n-1\) by the value \(a=11/360\) of the Dirac field (Equation~(275) of~\cite{Solodukhin2011}, after Duff; not re-derived here), which gives \(-\tfrac{11}{90}\) for a sphere. With one proper-time limit, the coefficient of \(\ln\LUV\) in the replica entropy is \(\int_\Sigma\sigma_{\log}\), where, with \(X'=-\operatorname{tr}(b_\sigma b^\sigma)\) in the vector sector, \(\tau_\Omega=\operatorname{tr}(\gamma^{12}\Omega^{f}_{12})\) the normal-plane curvature of the Dirac connection, with the same \(\Omega_{12}=[D_1,D_2]\), \(e\) the scalar endomorphism of the preceding paragraph and \(W=W_{ijij}\) the Weyl projection,

\[

  24\pi\,\sigma_{\log}=\tfrac43R_G+2R_{ii}-2R_{ijij}+2W-12\,e+\tfrac32X'-\tau_\Omega .

\]

On the self-dual and anti-self-dual backgrounds, \(W=0\) and \(e=R_G/6\), so that \(24\pi\sigma_{\log}=\tfrac32X'-\tau_\Omega\). At three sampled points of \(\Sigma\) for the rotation-invariant self-dual example above this equals \(2R_\Sigma\), twice the intrinsic curvature of \(\Sigma\) (exact rational arithmetic; not proved), so that \(\sigma_{\log}=R_\Sigma/12\pi\). The cone logarithmic coefficient is then an Euler density of \(\Sigma\) with nonvanishing local value. The coefficient \(\propto\int\rho^2\mathcal S_4\) of the induced action vanishes on these backgrounds, and so does the first variation of \(\int\rho^2\mathcal S_4\) along the smooth replica, because \(\mathcal S_4\ge0\) vanishes there. The two sides therefore differ at logarithmic order, as coefficients, for the fixed-state replica on these backgrounds; this is a statement about the fixed-state replica only and not about \(G_{\mathrm{eff}}\). The first-order smooth replica of the preceding paragraphs gives zero on both sides, and which replica the matrix state selects is not derived here. Contact and edge-mode terms of the connection and Dirac sectors, which are not included, shift the coefficient of \(R_\Sigma\) and could cancel the difference.



These statements are for the Poisson-limit operators, the Euclidean section, one loop, \(\alpha=1\), \(N_s=6\), \(n_D=2\), the totally geodesic \(\Sigma\), the fixed-state replica and one common proper-time limit. They use Equation~(70) of~\cite{Solodukhin2011} and the Dirac coefficient \(a=11/360\) as inputs, do not treat finite terms, edge modes, other extensions, higher orders in \(n-1\), or the non-self-dual \(\mathfrak J=0\) backgrounds beyond the general formula, and give no value or sign of \(G_{\mathrm{eff}}\). Problem~B remains open and PROGRAMME-CONDITIONAL.





\paragraph{Content dependence of the logarithmic comparison off \((N_s,n_D)=(6,2)\) (conditional; not closure).} \label{detail:sec-B-paragraphcontentdependenceofth-8}

Keep the setting of the two preceding paragraphs with \(N_s\) scalars and \(n_D\) Dirac fields in place of \((6,2)\), the vector with weight \(1\), the ghosts with weight \(-2\), and write \(N_{\mathrm{tot}}=N_s+2-4n_D\), \(\kappa_X=1-\tfrac12n_D\) and \(\bar a_1=\sum_iw_i\operatorname{tr}(\tfrac R6-E_i)=N_{\mathrm{tot}}(\tfrac R6-e)+\kappa_XX'\) for the weighted \(a_1\) density, whose integral is the \(\LUV^2\) density of the induced action. Off \((6,2)\) the logarithmic coefficients are not constrained to vanish, so this is a comparison with nonzero values, for contents that are not the maximally supersymmetric one. In units \(1/24\pi\) on a totally geodesic \(\Sigma\), the replica with all sectors periodic and no deltas, the fixed-state replica of the preceding paragraphs and the replica completed by the uniform \(E_\delta\) and by \(\Omega_\delta=2\pi(1-n)\mathcal J\delta_\Sigma\) of the vector bundle give

\[

  24\pi\sigma_{\mathrm{per}}=\bar a_1-\tfrac{N_{\mathrm{tot}}}{30}(R_\Sigma-3W),\qquad

  24\pi\sigma_{\mathrm{cmp}}=-\tfrac12\bar a_1+\kappa_X\tau_\Omega-\tfrac{N_{\mathrm{tot}}}{30}(R_\Sigma-3W),\qquad

  24\pi\sigma_{\mathrm{fs}}=24\pi\sigma_{\mathrm{cmp}}+\tfrac32\bar a_1^{\mathrm{bos}}-\tau_\Omega ,

\]

with \(\bar a_1^{\mathrm{bos}}\) the bosonic part of \(\bar a_1\) and \(\tau_\Omega\) as in the preceding paragraphs. For one Dirac field the antiperiodic value of the preceding paragraphs equals the periodic value plus the cross terms of a uniform \(E_\delta\) and of \(\Omega_\delta\) with \(\mathcal J\to\tfrac12\gamma^{12}\), an exact identity in \(R\), \(R_{ii}\), \(R_{ijij}\), \(e\), \(X'\) and \(\tau_\Omega\) given the Dirac weights \((-\tfrac12,1,\tfrac12)\); the contacts therefore equal \(-\tfrac32\bar a_1+\kappa_X\tau_\Omega\) for every content, proportional to the weighted \(a_1\) density and to the weighted Dynkin index, the data of the \(\LUV^2\) order. They vanish at \((6,2)\) for this reason, and the three coefficients differ off \((6,2)\) only by terms in \(\bar a_1\) and \(\tau_\Omega\). These relations were checked in exact arithmetic at twenty sampled points (the rotation-invariant self-dual family, the non-self-dual \(\mathfrak J=0\) family, \(\mathfrak J\neq0\) radial profiles and generic closed \(\omega\)); at the \(\mathfrak J\neq0\) and generic points they are pointwise algebraic identities, with \(\Sigma\) not necessarily totally geodesic, and the rotation-invariant closed form rests on eleven points. On the rotation-invariant self-dual family \(24\pi\sigma_{\mathrm{cmp}}=(\kappa_X-N_{\mathrm{tot}}/30)R_\Sigma=\tfrac{1}{30}(28-N_s-11n_D)R_\Sigma\) for every content, without the term in \(g\) of the preceding paragraph, and \(24\pi\sigma_{\mathrm{fs}}=24\pi\sigma_{\mathrm{cmp}}-R_\Sigma+\tfrac32g\). On the self-dual and anti-self-dual points and on the sampled \(\mathfrak J=0\) points the weighted \(a_2\) density equals \(N_{\mathrm{tot}}s+\kappa_X(v-4s)\), with \(s\) and \(v\) the single-sector scalar and vector values, since the Dirac value minus \(4s\) is half of the vector value minus \(4s\); it vanishes on all these backgrounds simultaneously only for \(N_{\mathrm{tot}}=\kappa_X=0\), that is \((N_s,n_D)=(6,2)\).



\emph{Weyl coefficient on plane waves.} A comparison that is not affected by the choice between the variables \((G,\rho,\partial\Theta)\) and \((G,\text{torsion})\) for the endomorphism contacts is the coefficient of \(C^2\), because \(R\), \(e\), \(X'\) and \(\bar a_1\) vanish to first order on the modes with \(\delta\rho=0\) about a constant self-dual or anti-self-dual \(\bar\Theta\). It is read from the quadratic form of the weighted \(a_2\) on those modes, against the coefficient of the Weyl projection \(W\) in the first-order area-side density, where a metric functional \(\kappa_CC^2+\kappa_E\mathcal E_4+\kappa_RR^2\) is assigned the entropy density \(12\kappa_CW+12\kappa_ER_\Sigma+12\kappa_RR\) by the cross-term rule. The rule does not fix the entropy of a part of the quadratic form carried by non-metric invariants of \(\Theta\); that part is the gap below. At second order the single-sector values satisfy \(v-4s=\tfrac1{12}\operatorname{tr}\Omega^2\) and \(d-4s=\tfrac12(v-4s)\), the vector connection \(b\) contributes \(X^{(1)}=0\), and on the modes considered \(\tfrac1{12}\operatorname{tr}\Omega^2=-\tfrac13C^2\) while \(\tau_\Omega^{(1)}=-R_{ijij}=-2W+R_\Sigma\) (first-order formulas checked against the jet code at \(\epsilon=10^{-5}\); six modes for each of the two chiralities). The coefficient of \(C^2\) is then \(\kappa_C^T=N_{\mathrm{tot}}/120-\kappa_X/3\), and the coefficients of \(W\) of the three replicas minus \(12\kappa_C^T\) are

\[

  \text{periodic: }4-2n_D,\qquad \text{fixed-state: }4-n_D,\qquad \text{completed: }2-n_D,

\]

independent of \(N_s\). The first two equal the differences \(c_A-c_T\) at order \(\LUV^2\) of the paragraphs on the quadratic form and on fermionic content; the completed gap has no \(\LUV^2\) counterpart of the same form. The periodic and completed gaps vanish on the line \(n_D=2\) for every \(N_s\), where both coefficients equal \(N_{\mathrm{tot}}\) times the conformal-scalar value and are nonzero for \(N_s\neq6\); the fixed-state gap vanishes on \(n_D=4\). The completed gap is \(2\kappa_X\). On these modes \(\operatorname{tr}\Omega^2=-4C^2\) is twice its Levi-Civita value \(-2C^2\), whereas \(\operatorname{tr}(\mathcal J\Omega_{12})=-R_{ijij}\) has its Levi-Civita value, so that the cross term with \(\Omega_\delta\) accounts for half of the contribution of the \(\Omega^2\) term to the coefficient of \(W\). The bundle curvature \(\Omega^v\) is not determined by the Riemann tensor: at fixed first jet its second jet has 33 independent components against the 20 of the Riemann tensor (45 random second jets, exact rank), and the Riemann second jet is a linear image of that of \(\Omega^v\). The matching on \(n_D=2\) is the universal \(N_{\mathrm{tot}}\) counting and does not test the index-weighted sector.



The contents \((N_s,n_D)\neq(6,2)\) are parameters of the computation and are not proposed field contents of the matrix model. The replicas are the hybrid fixed-state replica, whose Dirac fields have the antiperiodic gluing and are not the trace of the \(n\)-th power of one state, the periodic replica, whose Dirac fields have the unrotated periodic gluing (the graded trace), and the replica completed by the uniform \(E_\delta\) and by \(\Omega_\delta\) on every component, with one common proper-time limit for all sectors.



These statements are for the Poisson-limit operators, one loop, \(\alpha=1\), the contents stated with the Yukawa couplings dropped, a totally geodesic \(\Sigma\), cross terms at first order in \(1-n\) on the \(n<1\) branch, the imported contact prescription and the delta content of the preceding paragraph, with the hypotheses stated there, and, for the plane-wave comparison, a constant self-dual or anti-self-dual \(\bar\Theta\), the modes with \(\delta\rho=0\) and the Weyl coefficient only; the coefficient of the Euler density is invisible to the quadratic form. They use the Dirac weights and Equation~(70) of~\cite{Solodukhin2011} as before, treat neither finite terms, nor the mode with \(\delta\rho\neq0\), nor the Hodge distribution of the Ricci-type delta, and give no value or sign of \(G_{\mathrm{eff}}\). Problem~B stays open and PROGRAMME-CONDITIONAL.





\paragraph{Edge and contact terms of the vector and connection sector at the Weyl coefficient (conditional; not closure).} \label{detail:sec-B-paragraphedgeandcontacttermsof-9}

Keep the setting of the preceding paragraph: the weighted \(a_2\) density, the modes with \(\delta\rho=0\) about a constant self-dual or anti-self-dual \(\bar\Theta\), the completed replica, a totally geodesic \(\Sigma\), cross terms at first order in \(1-n\), and the completed gap \(2\kappa_X\) in the coefficient of \(W\), read in the basis \((W,R_\Sigma,R,e,X')\). On the twelve modes \(W\) and \(R_\Sigma\) are independent (rank 2), with \(R_{ii}=W-R_\Sigma\) and \(R_{ijij}=2W-R_\Sigma\).



\emph{Endomorphism contacts.} To first order in the amplitude about a constant \(\bar\Theta\) the endomorphisms \(E_V\) of the vector and \(E_f\) of the Dirac operator equal \(e\,\mathbb I\): the antisymmetric part is a divergence of \(\Theta^{\rho\sigma}\partial_\rho\Theta^{\mu\nu}\), which vanishes at first order about constant \(\bar\Theta\), and the symmetric traceless part is quadratic (exact arithmetic, cubic jets of both chiralities, two amplitudes). Since \(e\) vanishes for \(\delta\rho=0\), every cross term of \(a_2\) with a delta of an endomorphism, uniform or confined to the normal block, vanishes at this order for every content: a contact of Kabat type has no coefficient of \(W\) in this presentation. For the Hodge operator of the metric \(G\) (\(E=-\mathrm{Ric}\), \(\Omega=\mathrm{Riem}\)), the same cross-term rule on the same modes gives \(6/5\) and \(-31/15\) for the Maxwell field with its ghosts, equal to its \(12\kappa_C\) in the coefficient of \(W\), of which the endomorphism contact carries \(3\) (the value is \(-9/5\) without it). The matrix-derived vector operator is not the Hodge operator: on these modes \(\mathrm{Riem}^2=2C^2\) and \(\mathrm{Ric}^2=C^2/2\), and \(v-4s\) is \(-C^2/3\) against \(C^2/12\). That mechanism is therefore not available for it.



\emph{Edge modes.} In the classes treated in~\cite{DonnellyWall2016}, products of a two-dimensional base with a bifurcate horizon and a compact fibre, and \(S^1\times H^3\), the edge sector of the Maxwell field is a scalar of wrong-sign statistics on the entangling surface, and its logarithmic density is the two-dimensional trace anomaly of the Laplacian of \(\Sigma\), \(-R_\Sigma/24\pi\); that paper does not treat the dependence on the Weyl tensor, and leaves that on the extrinsic geometry open. In the basis above \(-R_\Sigma\) has coefficient \(0\) of \(W\). A weight \(w\) built from \(\operatorname{Pf}\Theta\) and \(\operatorname{tr}\Theta^2\) in \(-w^{-1}\operatorname{div}(w\nabla)\) adds nothing at first order, because the Hessians of both vanish on the modes. A potential \(E_\Sigma=\alpha R_{ii}+\beta R_{ijij}\) would change the density to \(-(R_\Sigma+6E_\Sigma)\) and would need \(\alpha+2\beta=\kappa_X/3\) for the coefficient \(-2\kappa_X\) of \(W\). This depends on \(n_D\) through the Dirac weight, whereas the vector carries one edge scalar, if it carries the edge sector of a Maxwell field, and the Dirac fields none. No edge term independent of \(n_D\) equals \(2\kappa_X\) for all \(n_D\).



\emph{Holonomy contact.} For the diagonal profile with \(\lvert\Theta^{34}\rvert=r^{-\epsilon c_b}\) the holonomy is \(H=c_a\mathcal J_{12}+\chi c_b\mathcal J_{34}\), with \(\chi=+1\) about a self-dual and \(\chi=-1\) about an anti-self-dual \(\bar\Theta\) (symbolic computation), and the coefficient of \(W\) in \(\operatorname{tr}(\mathcal J_{34}\Omega_{12})\) is \(2\chi\). The gap, in the extended cross-term rule of the preceding paragraph, is therefore the same for both chiralities, \(N_{\mathrm{tot}}(3c_a-2c_b-3)/30+2\kappa_X(2-c_a+c_b)\). With \(s=c_a-c_b\), the contraction of the first-order Weyl tensor with the Weyl tensor of the deltas is \(2sW\), and the Euler combination \(\mathrm{Riem}\cdot\mathrm{Riem}'-4\,\mathrm{Ric}\cdot\mathrm{Ric}'\) is \(2c_aR_\Sigma-2c_bW\), which contains \(W\) for \(c_b\neq0\). If that \(W\) part is dropped as topological, the gap is \((s-1)N_{\mathrm{tot}}/10-2(s-2)\kappa_X\): the metric part requires \(s=1\) and the connection part \(s=2\). The profile with \(\delta\rho=0\), \(c_b=-c_a\), has \(s=2\) and removes the \(\kappa_X\) part, leaving \(N_{\mathrm{tot}}/15\) (\(N_{\mathrm{tot}}/10\) after dropping the Euler part); its induced metric on \(\Sigma\) is singular at first order and the scalar of its Riemann delta, \(2c_a+4c_b\), has the sign opposite to that of the cone.



No edge, endomorphism-contact or holonomy term of this family therefore supplies \(-2\kappa_X W\) for all contents. The required term is proportional to the index-weighted sum that multiplies \(\operatorname{tr}\Omega^2\), the same for every fibre representation; of the candidates examined only a holonomy generator beyond \(\mathcal J_{12}\) attached to \(\Sigma\) has this property, and none of the profiles of the diagonal family provides it without a singular metric on \(\Sigma\). The gap is also relative to the comparator \(12\kappa_C^T\) assigned to \(\operatorname{tr}\Omega^2\) by the metric-universal rule; whether the entropy of this non-metric term is that value or the Noether value \(-2\kappa_X\) of the preceding paragraph, which equals the replica, is not decided here.



These statements are for the Poisson-limit operators, one loop, \(\alpha=1\), a totally geodesic \(\Sigma\), cross terms at first order in \(1-n\) on the \(n<1\) branch, the modes and replica of the preceding paragraphs, the cross-term rule extended to \(c_b\neq0\) (not derived), and an edge sector of the abelian type of the cited paper. They do not treat non-diagonal profiles, the non-delta singular first-order parts, the exact heat trace of the cone with holonomy, the symplectic gauge group, or the mode with \(\delta\rho\neq0\), and give no value or sign of \(G_{\mathrm{eff}}\). Problem~B stays open and PROGRAMME-CONDITIONAL.







\section{Problem C: conformal mode and Cassini}

\label{sec:detailed-Problem C: conformal mode and Cassini}



\paragraph{Ledger label.} \label{detail:sec-C-paragraphledgerlabel-10}

Status: OPEN. Priority: CRITICAL PATH. Completion:

PROGRAMME-CONDITIONAL. A scalar-tensor result is a reduced-scope

outcome only when its operator, coupling and phenomenological test

are supplied; it does not by itself meet the matrix-derived criterion.



The radial singlet has the tree-level mass

\[

m_\sigma^2=\lambda v^2.

\]

This does not by itself determine the mass or matter coupling of a

scalar mode of the composite metric. In particular, the stated

Repair~I prescription is intended to build a vierbein from the Lorentz-valued

polar factor \(\Lambda\), evaluated at \(\sigma=0\), rather than from

the radial field.





\paragraph{Status.} \label{detail:sec-C-paragraphstatus-11}

Under the stated Repair~I ingredients, the linearized composite

vierbein required for the later metric-scalar analysis is not

constructed. This is an obstruction to that frame prescription, not

a proof that a completed theory has or lacks a physical scalar pole.

Repair~II, an independent coframe, or additional symmetry-breaking

background tensors would change the construction and require a

separate analysis. The radial mass \(m_\sigma^2=\lambda v^2\) remains

established, but it is not a metric-scalar mass. The scalar spectrum

and matter coupling relevant to Cassini are therefore still

undetermined. The Repair~II pure-trace relation \(\psi=8\sigma/v\)

with \(m_\sigma^2=\lambda v^2\) is a potential-only statement; the

induced Einstein--Hilbert term adds to the kinetic coefficient of

\(\sigma\), so it does not fix a physical pole mass, mixing or

Cassini coupling.



In the Repair~II pure-trace truncation, adding the induced

Einstein--Hilbert term to the microscopic kinetic term gives, in the

healthy-sign normalization,

\(Z=2N-3/(4\pi G_{\mathrm{eff}}v^{2})\)

(Equation~\ref{eq:Z-sigma}), conditional on the sign of

\eqref{eq:Smicro} as corrected in Chapter~3 and on the declared

curvature convention. For \(Z>0\), the truncated pole has

\(m^{2}=V''/Z\); for \(Z<0\) it is ghost-like and, for \(\lambda>0\),

tachyonic in this truncation, which therefore cannot support

Assumption~\ref{ass:gap-psi} there. This signals an instability of the

chosen background and truncation, not a statement about the full scalar

spectrum. At \(Z=0\) the quadratic kinetic

term degenerates and establishes no canonically normalized

finite-mass pole. This result also assumes \(G_{\mathrm{eff}}>0\),

which is not derived (Problem~F), and does not determine mixing, the

non-minimal coupling, the matter coupling, or a Cassini bound.

If matter is assumed to couple only through \(g=\Omega^2\eta\), the

unmixed truncated scalar exchanges, for \(Z>0\), a Yukawa potential of

strength \(\beta=M_{\mathrm P}^{2}/(Nv^{2}-3M_{\mathrm P}^{2})\)

relative to the nominal Einstein-sector Newtonian potential

(Chapter~17); this is conditional on the assumptions listed there and

is not a Cassini bound.



\emph{Off-shell trace--tensor mixing from the induced

Einstein--Hilbert term (conditional; not closure).}

Assume Repair~II, which is assumed rather than derived, with

\(g_{\mu\nu}=\eta_{\mu\nu}+h_{\mu\nu}\) and

\(h_{\mu\nu}=T_{\mu\nu}+\tfrac14\eta_{\mu\nu}\psi\). Also assume the

induced Lorentzian term is

\(S_{\mathrm{EH}}=(16\pi G_{\mathrm{eff}})^{-1}\int\sqrt{-g}\,R\),

with \(G_{\mathrm{eff}}>0\). For a general symmetric \(h_{\mu\nu}\),

define \(\mathcal E_{\mu\nu}[h]:=-2G^{(1)}_{\mu\nu}[h]\); on traceless

tensors this agrees with \(\mathcal E_{\mu\nu}\) of

Equation~\ref{eq:einstein-op}. In these conventions, up to boundary

terms,

\[

  S_{\mathrm{EH}}^{(2)}

  =\frac{1}{16\pi G_{\mathrm{eff}}}

   \int\mathrm{d}^4x\,\frac14 h^{\mu\nu}\mathcal E_{\mu\nu}[h].

\]

For \(s_{\mu\nu}=\tfrac14\eta_{\mu\nu}\psi\),

\[

  \mathcal E_{\mu\nu}[s]

  =\frac12\bigl(\partial_\mu\partial_\nu\psi

                -\eta_{\mu\nu}\Box\psi\bigr).

\]

Expanding the quadratic action in this field basis gives

\[

  S_{\mathrm{EH}}^{(2)}

  =\frac{1}{16\pi G_{\mathrm{eff}}}

   \int\mathrm{d}^4x\left[

     \frac14T^{\mu\nu}\mathcal E_{\mu\nu}[T]

    +\frac14T^{\mu\nu}\partial_\mu\partial_\nu\psi

    +\frac{3}{32}(\partial\psi)^2

   \right].

\]

Thus, using \(\psi=8\sigma/v\) at linear order, the cross term is

\[

  \frac{1}{16\pi G_{\mathrm{eff}}}

  \int\mathrm{d}^4x\,\frac{2}{v}

  T^{\mu\nu}\partial_\mu\partial_\nu\sigma .

\]

The tensor-only term here is normalized relative to the assumed

\(S_{\mathrm{EH}}\); Equation~\ref{eq:EH-lin} states a structure and

is not used to set this coefficient. The cross term is

basis-dependent: it is the usual conformal-mode mixing in the

gauge-invariant quadratic Einstein--Hilbert action and can be

rearranged by a field redefinition. It does not by itself establish an

observable effect or scalar-exchange strength. Although it vanishes if

transversality is imposed, transversality cannot be assumed when

forming the coupled off-shell Hessian. The displayed microscopic action

on the fixed scaffold does not specify the full coupled, gauge-fixed

tensor operator or the remaining scalar couplings.

Problems~B, C, D, F and~G remain open.





\paragraph{Ledger label.} \label{detail:sec-C-paragraphledgerlabel-12}

Status: OPEN. Priority: CRITICAL PATH. Completion:

PROGRAMME-CONDITIONAL. Reduced-scope result: D2, imposed local

Lorentz symmetry, supports only a scope-limited Presentation-B

model; D remains OPEN and the matrix-derived criterion is unmet.



The counting of generators under the locked vacuum is given in

Chapter~10. The condensate term supplies a quadratic Higgs-type

contribution in the broken-direction combination \(\Xi\); it does not

establish a physical pole mass without kinetic terms and canonical

normalization. The diagonal connection's dynamics and torsion remain

open. Torsion-free behavior follows only under the conditional

Palatini--Holst hypotheses of Proposition~

\ref{prop:B-torsion-free}; NAMCG has not been shown to induce those

hypotheses.



For a Laplace-type operator,

\(a_1=\operatorname{tr}(R(g)/6-E)\); the explicit universal

curvature term is the Levi-Civita scalar \(R(g)\). Because \(E\) is

unspecified, no conclusion about additional \(B\)-dependence of

\(a_1\) follows, and a Palatini \(R(B)\) term is not established. If

\(B\) enters the bundle connection of the actual operator, \(a_2\)

may contain a \(\operatorname{tr}F(B)^2\) term; whether it survives

in the weighted trace depends on the unspecified operator

representation and weights. Deriving the \(B\)-dependence of \(E\)

and the bundle connection remains open.





\paragraph{Symmetries supplied by the stated Presentation~} \label{detail:sec-C-paragraphsymmetriessuppliedbyt-13}

(matrix-derived criterion of Problem~D; not closure).}

The matrix-coordinate action~\eqref{eq:Smatrix} is invariant under

unitary conjugation

\(\hat X^\mu\mapsto U\hat X^\mu U^\dagger\), provided the matter

variables transform so that their term is invariant: each matrix

commutator transforms by conjugation and the trace is cyclic. This

is a matrix gauge transformation

\(\delta\hat X^\mu=i[\epsilon,\hat X^\mu]\); the stated action and

infrared dictionary do not identify it with a transformation of an

independent tangent frame.



As an action-level statement, the commutator term is covariant under

constant Lorentz transformations of the spacetime index \(\mu\),

because its indices are contracted with \(\eta_{\mu\nu}\). The cubic

term is covariant if \(C_{\mu\nu\lambda}\) transforms as a tensor;

for fixed \(C_{\mu\nu\lambda}\), it is a symmetry only when \(C\) is

invariant. The matter term must obey the corresponding transformation

law. Likewise, the chosen background \(\Theta^{\mu\nu}\) in

Equation~\eqref{eq:NC} is invariant only under its stabilizer if held

fixed; transforming \(\Theta\) along with the fields expresses

covariance, not a symmetry of that fixed background. None of these

constant transformations establishes a spacetime-dependent local

Lorentz gauge symmetry.



On the \(C_{\mu\nu\lambda}=0\), central constant-\(\Theta\) branch of

Remark~\ref{rem:ch3-trace-quadratic}, evaluate the matrix-gauge variation

on the background \(\bar X\). By the Jacobi identity,

\[

  \delta_\epsilon[\bar X^\mu,\bar X^\nu]

  =i[\epsilon,[\bar X^\mu,\bar X^\nu]]=0,

\]

since \([\bar X^\mu,\bar X^\nu]=i\Theta^{\mu\nu}\mathbb I\). In the

irreducible Moyal-symbol realization, the corresponding linearized

fluctuation transforms as

\(\delta a^\mu=\Theta^{\mu\rho}\partial_\rho\lambda\), leaving

\(\mathcal F_\Theta^{\mu\nu}\) invariant as shown in that remark. This

local matrix-gauge statement is distinct from the global constant

Lorentz covariance described in the preceding paragraph, which requires

transforming \(\Theta\) as a tensor and is a symmetry of a fixed

\(\Theta\) only for transformations that stabilize it. On this branch,

matrix conjugation alone supplies neither a transformation of an

independent tangent-frame index nor six independent local Lorentz

parameters. This is a limitation of the stated matrix-gauge action on

the central background, not a no-go for an enlarged matrix model with a

derived frame map and additional fields.



More specifically, on the central, nondegenerate constant-\(\Theta\)

Moyal branch, define

\(D'_A f:=\Lambda_A{}^B\star D_B f\), where

\(D_B=i[\bar X_B,\,\cdot\,]\) are the inner derivations, whose

principal symbols \(\Theta_B{}^\mu\partial_\mu\) are linearly

independent. The Leibniz defect is

\[

  D'_A(f\star g)-(D'_Af)\star g-f\star(D'_Ag)

  =[\Lambda_A{}^B,f]_\star\star D_Bg .

\]

Independence of the principal symbols implies that this defect

vanishes for all \(f,g\) exactly when every \(\Lambda_A{}^B\) is

central. In the irreducible nondegenerate constant-\(\Theta\)

representation, the centre consists of constants, i.e.\ multiples of

the identity. This conclusion is limited to that branch and this

left-multiplication action; it does not rule out a connection-twisted

or module-valued frame, or an enlarged matrix model. A matrix-derived

frame module with a connection that compensates this defect is a

conjectured route, not a result established here. Problems~B, C, D, F

and~G remain open.



For real, nondegenerate constant \(\Theta\), consider

\(\Lambda\in\mathrm{SO}(3,1)\) that preserves the fixed tensor,

\(\Lambda\Theta\Lambda^T=\Theta\). In a Lorentz frame with

\(\Theta^{01}=E\), \(\Theta^{23}=B\), \(EB\ne0\), and all other

independent components zero, its identity-component stabilizer is

\(\mathrm{SO}^+(1,1)\times\mathrm{SO}(2)\): a boost in the

\(01\)-plane and a rotation in the \(23\)-plane. Equivalently, in

the self-dual representation its stabilizer is

\(\mathrm{SO}(2,\mathbb C)\), of real dimension two. This is a

constant-background stabilizer, not a local Lorentz gauge symmetry.



The corresponding linear automorphisms of the Heisenberg algebra

are implementable by inner automorphisms of its irreducible Weyl

representation. Infinitesimally, write \(\Lambda=\mathbb I+A\);

preservation of \(\Theta\) requires

\[

  A\Theta+\Theta A^T=0.

\]

Since \(\Theta\) is invertible,

\[

  H_{\mu\nu}:=(\Theta^{-1})_{\mu\rho}A^\rho{}_\nu

\]

is symmetric. On the common invariant domain of the Weyl

generators, the symmetrically ordered quadratic generator

\[

  G=\frac14 H_{\mu\nu}

    \{\bar X^\mu,\bar X^\nu\}

\]

satisfies

\[

  \delta\bar X^\mu=i[G,\bar X^\mu]

  =A^\mu{}_\nu\bar X^\nu.

\]

The finite implementation is the metaplectic unitary, unique up

to phase by the Stone--von Neumann theorem. This implements only

the constant stabilizer: it supplies neither six arbitrary local

Lorentz parameters nor an independent tangent-frame index or spin

connection. Invariance of the full action also requires the

remaining fields and couplings to transform compatibly. This

conditional representation-theoretic result does not establish

Problem~D1. Problems~B, C, D, F and~G remain open.



A conditional classical-symbol construction is possible if the

background Poisson tensor varies. Let

\(\Theta^{\mu\nu}(x)\) be a real, nondegenerate Poisson tensor,

so its bracket

\[

  \{f,g\}_\Theta

  :=\Theta^{\mu\nu}\partial_\mu f\,\partial_\nu g

\]

obeys the Jacobi identity. For local coordinate functions \(x_A\),

with the frame index \(A\) identified with the coordinate label,

define the Hamiltonian vector fields

\(V_A:=X_{x_A}\), where \(X_f(g):=\{f,g\}_\Theta\), and assume

they form an invertible frame. The Poisson Jacobi identity gives

\[

  [V_A,V_B]=X_{\Theta_{AB}},

  \qquad

  \Theta_{AB}:=\{x_A,x_B\}_\Theta .

\]

Define \(C^C{}_{AB}\) by

\([V_A,V_B]=C^C{}_{AB}V_C\). A provisional candidate inverse

metric is

\[

  g_{\mathrm{cand}}^{\mu\nu}

  :=\eta^{AB}V_A{}^\mu V_B{}^\nu .

\]

It is Lorentzian by construction when the frame is invertible, but

is not thereby identified with the emergent metric of

Chapter~\ref{ch:vierbein-emergence}.



Under a local frame change

\(V'_A=\Lambda_A{}^B(x)V_B\), with

\(\Lambda(x)\in\mathrm{SO}(3,1)\), the bracket, expanded in the

original frame \(V_F\), is

\[

  [V'_A,V'_B]

  =

  \left(

    \Lambda_A{}^D\Lambda_B{}^E C^F{}_{DE}

    +\Lambda_A{}^D V_D(\Lambda_B{}^F)

    -\Lambda_B{}^D V_D(\Lambda_A{}^F)

  \right)V_F .

\]

Accordingly, the structure functions of the new frame are the

displayed coefficients contracted with

\((\Lambda^{-1})_F{}^G\), since

\([V'_A,V'_B]=C'^G{}_{AB}V'_G\). For this provisional frame

and metric, the dual coframe \(\vartheta^A\) determines the

composite Levi-Civita spin connection through

\[

  \mathrm{d}\vartheta^A+\omega^A{}_B\wedge\vartheta^B=0,

  \qquad

  \omega_{AB}=-\omega_{BA}.

\]

It transforms inhomogeneously under the local frame change.



For constant \(\Theta\), the fields \(V_A{}^\mu\) are constant,

so \(C^C{}_{AB}=0\) and this composite connection and its

curvature vanish in that frame. A varying \(\Theta\) is therefore

needed for this construction to be nontrivial, although variation

alone does not guarantee nonzero curvature. This is a

classical Poisson-symbol construction: it does not address the

exact Moyal Leibniz defect for noncentral frame coefficients, derive

local frame rotations from the matrix action, or supply independent

connection dynamics. Problems~B, C, D, F and~G remain open.



For the hand-chosen Poisson background of

Remark~\ref{rem:ch3-varying-theta},

\[

  \Theta=a\,\partial_0\wedge\partial_1

         +b(y)\,\partial_2\wedge\partial_3,

  \qquad b(y)=b_0e^{\kappa y^2},

\]

the Hamiltonian frame \(V_A=X_{x_A}\), with

\(X_f(g)=\{f,g\}_\Theta\), is

\[

  (V_0,V_1,V_2,V_3)

  =(a\partial_1,-a\partial_0,b\partial_3,-b\partial_y).

\]

Its only nonzero frame bracket is

\[

  [V_2,V_3]=b'(y)V_2=2\kappa y\,b(y)\,V_2.

\]

The associated structure function is \(C^2{}_{23}=b'(y)\);

the composite connection contains the logarithmic derivative

\[

  \omega^2{}_3=\frac{b'(y)}{b(y)}\,\mathrm{d}x^3

  =2\kappa y\,\mathrm{d}x^3.

\]

With \(\eta_{AB}=\operatorname{diag}(-,+,+,+)\), the provisional

inverse metric constructed from this frame is

\[

  g_{\mathrm{cand}}^{\mu\nu}

  =\eta^{AB}V_A{}^\mu V_B{}^\nu

  =\operatorname{diag}(a^2,-a^2,b(y)^2,b(y)^2)

  =K_\Theta^{\mu\nu}.

\]

It therefore reproduces the tensor used in

Remark~\ref{rem:ch3-varying-theta}; with that remark's

normalization \(K_\Theta^{\mu\nu}=\ell_*^4\tilde g^{\mu\nu}\),

\(g_{\mathrm{cand}}=\ell_*^{-4}\tilde g\).



A dual coframe is

\[

  \vartheta^0=\frac{\mathrm{d}x^1}{a},\qquad

  \vartheta^1=-\frac{\mathrm{d}x^0}{a},\qquad

  \vartheta^2=\frac{\mathrm{d}x^3}{b(y)},\qquad

  \vartheta^3=-\frac{\mathrm{d}y}{b(y)}.

\]

Its torsion-free connection has

\(\omega^2{}_3=2\kappa y\,\mathrm{d}x^3\), and

\(\mathcal R^2{}_3=2\kappa\,\mathrm{d}y\wedge\mathrm{d}x^3\).

In the curvature convention of

Remark~\ref{rem:ch3-varying-theta}, this gives

\(R[g_{\mathrm{cand}}]=4\kappa b(y)^2\).

This identifies the candidate Poisson-frame metric with the

already chosen toy geometry, up to the stated constant

normalization; it is not an independent derivation of that

geometry.



The calculation is at classical Poisson-symbol order on a

hand-chosen background. It derives neither that background from

the matrix action nor local Lorentz gauge symmetry, a

noncommutative frame module, or independent connection dynamics.

Problems~B, C, D, F and~G remain open.



For the Poisson block above, restrict to the

\(C_{\mu\nu\lambda}=0\) branch of~\eqref{eq:Smatrix}, retain only

the quartic matrix term, and omit matter. The matrix equation is

\[

  [\hat X_\nu,[\hat X^\nu,\hat X^\mu]]=0,

\]

with indices contracted using \(\eta\). At leading

Poisson-symbol order, using

\([\hat f,\hat g]\sim i\{f,g\}_\Theta\), this requires

\[

  \eta_{\nu\lambda}\Theta^{\lambda\rho}

  \partial_\rho\Theta^{\nu\mu}=0.

\]

For \(\Theta^{01}=a\), \(\Theta^{23}=b(y)\), with \(y=x^2\),

the \(\mu=3\) component vanishes identically, while the

\(\mu=2\) component is \(b(y)b'(y)\). Nondegeneracy requires

\(b\ne0\), so this leading-order source-free equation forces

\(b'=0\) on each connected patch. Hence the hand-chosen curved

profile \(b=b_0e^{\kappa y^2}\), with \(\kappa\ne0\), is not a

source-free leading-order solution of this ansatz. Supporting that

profile would require a specified balancing contribution, such

as a cubic term, matter source, or corrections from a specified

deformation quantization. This calculation is at classical

Poisson-symbol order only; it does not define a beyond-leading-order

product for varying \(\Theta\), make a claim about the full matrix

model, or exclude other backgrounds. It is not a derivation or

closure of Problem~D1. Problems~B, C, D, F and~G remain open.



In particular, Presentation~A as specified contains no coframe

\(e^A{}_\mu\), independent tangent index \(A\), or spin connection

\(\omega_\mu{}^{AB}\). The infrared dictionary~\eqref{eq:dictionary}

is an assumption mapping a matrix commutator to a covariant derivative

with \(A_\mu\); it supplies neither a local Lorentz transformation

law for a derived frame nor the transformation and dynamics of

\(\omega_\mu{}^{AB}\). The boost obstruction recorded in Problem~F

concerns the displayed Presentation-B potential, not the

Presentation-A matrix action. These facts delimit what the stated

action and dictionary establish; they do not rule out an enlarged

matrix construction. Consistent with the preceding paragraph, the

matrix-derived criterion of Problem~D remains open: closure requires

an explicit matrix-to-frame map and a derivation of the local

transformation law, invariant action, and connection content from

that map. Problems~B, C, D, F and~G remain open.





\paragraph{Ghost-free versus boost-invariant contraction (conditional; not closure).} \label{detail:sec-C-paragraphghostfreeversusboosti-14}

On the real locked \(4\times4\) block write \(\Phi=\mathbb{I}+M\). Consider two

candidate quadratic kinetic contractions, both stated for the fixed scaffold

and neither derived from~\eqref{eq:Smicro}.

(i) The kinetic term in~\eqref{eq:Smicro}, with its positive-definite

matrix-space contraction, gives every component the same relative kinetic

sign, so there are no relative-sign differences; however, the displayed

contractions in~\eqref{eq:Smicro} and~\eqref{eq:V} are not invariant under

boosts of either index, as noted above.

(ii) A contraction with \(\eta\) on both indices, invariant under constant

independent \(\mathrm{SO}(3,1)\times\mathrm{SO}(3,1)\) transformations,

gives component \(M^I{}_A\) the relative kinetic sign

\(\eta_{II}\eta_{AA}\). Fix the overall sign so that \(M^0{}_0\) and

\(M^i{}_j\) have the healthy sign; then the six components

\(M^0{}_i,M^i{}_0\) have the opposite sign. This invariance is global: for

local transformations covariant derivatives and the corresponding connections

are needed, and the candidate expression does not specify them. If the six

broken generators of Proposition~\ref{prop:goldstone-count} are gauged and eat

the six \(\eta\)-antisymmetric components, three of the opposite-sign

components (the \(0i\) part of the antisymmetric block) are absorbed; three

opposite-sign \(\eta\)-symmetric \(0i\) components remain. Whether these are

removed by a constraint, a further gauge condition, or the diffeomorphism

structure of Chapter~\ref{ch:lorentz-locking} is not established. This is a

sign count for the quadratic form about \(\Phi=\mathbb{I}\) only; it is not a

spectrum or unitarity result and is not identified with the tensor sector of

Chapter~\ref{ch:lorentz-decomposition}. Boost invariance and absence of

opposite-sign components are therefore not both available from the

contractions considered here. Problem~D remains open.





\paragraph{Diffeomorphism symmetry of the displayed action (conditional; not closure).} \label{detail:sec-C-paragraphdiffeomorphismsymmetr-15}

In~\eqref{eq:Smicro} the field \(\Phi^I{}_A\) carries no coordinate index: it

is a set of scalar fields on the scaffold, and a world index appears only

through the identification of the tangent index with the coordinate index

in~\eqref{eq:vev-is-frame}. With \(g^{(0)}_{\mu\nu}\) held fixed

(Assumption~\ref{ass:aux-metric}), a diffeomorphism of the scaffold coordinates

acts on the argument \(x\) of~\(\Phi\), and \(S_{\mathrm{micro}}\) is

invariant under those diffeomorphisms that preserve \(g^{(0)}\), such as its

isometries; this does not exclude other symmetries. If \(g^{(0)}_{\mu\nu}\) is

transformed as a tensor together with \(\Phi\), the functional

\(S_{\mathrm{micro}}[\Phi,g^{(0)}]\) is covariant in the spurionic sense; this

does not by itself make \(g^{(0)}\) a dynamical or compensating field, nor

establish a gauge symmetry of the fixed-background theory. The composite

metric \(g=e^{T}\eta e\) built from the identified frame is a tensor only if

the identification \(A\leftrightarrow\mu\) is itself made covariant, which

\eqref{eq:Smicro} does not provide. The statement of

Section~\ref{sec:ch10-diffs} that diffeomorphisms are the gauge symmetry of

the composite metric, and the diffeomorphism-invariance hypotheses of

Theorem~\ref{thm:conservation} and Assumption~\ref{ass:Sent-density}, are

therefore conditional on such an extension, not consequences of

\(S_{\mathrm{micro}}\) as displayed. This is a scope statement; it does not

exclude a covariant completion, and Problems~B, D and~G remain open.





\paragraph{Related covariant quantum-space constructions, and what they d} \label{detail:sec-C-paragraphrelatedcovariantquant-16}

not supply (conditional; not closure).}

Steinacker and collaborators study the IKKT matrix model on covariant

quantum spaces in which the frame is read off from the adjoint action

\([X^a,\cdot\,]\) and averaged over an internal \(S^2\) fibre.  The Euclidean,

linearised analysis~\cite{Steinacker2016Gravity} is reported to give

linearised Einstein equations under conditions on the background, dimensional

reduction and an infrared cutoff; we have not checked these conditions.  A

Lorentzian space built on the minimal doubleton representation of

\(\mathfrak{so}(4,2)\), a quantised twistor space \(\mathbb{CP}^{1,2}\) over a

\(k=-1\) FLRW spacetime, is described in~\cite{MantaSteinacker2025}, and an

induced gravitational action with an explicit effective Newton constant is

reported at one loop for such backgrounds with fuzzy extra dimensions

in~\cite{MantaSteinackerTran2024}.  The absence of ghosts for the Lorentzian

higher-spin theory is reported in~\cite{Steinacker2019NoGhost}, at the expense

of manifest Lorentz invariance.  We have not reproduced any of these results.

Taking them as stated, they are analogous to, not the same as, the

construction tested above: the frame comes from a chosen representation

rather than from the vacuum value of a field \(\Phi^I{}_A\); the ghost problem

is addressed by giving up manifest Lorentz invariance, which the locking

assumption of Chapter~\ref{ch:lorentz-locking} presupposes; and gravity is

induced at one loop rather than obtained from entanglement stationarity as in

Problem~B.  They therefore do not close Problems~D, B, F, C or~G.  They report

that a Lorentzian, ghost-free construction exists in a matrix-model setting;

whether an analogue exists for the \(\Phi\)-field action of

Chapter~\ref{ch:matrix-action} with Lorentz locking retained is not

established here.  Adopting an induced-gravity route in place of Problem~B

would be a structural change to the programme and is not proposed here.





\paragraph{A compensating field alone does not change the spectru} \label{detail:sec-C-paragraphacompensatingfieldalo-17}

(conditional; not closure).}

Restore the shifts above by a Stueckelberg replacement

\(h_{\mu\nu}\to h_{\mu\nu}+\partial_{(\mu}\xi_{\nu)}\) in the \(\eta\)-contracted

candidate, with a new field \(\xi_\mu\) and gauge transformations

\(\delta h=\partial\epsilon\), \(\delta\xi=-\epsilon\).  The gauge \(\xi=0\) is

reachable and in it the action equals the original, so the physical spectrum

is unchanged: the signature \((7,3)\) of the self-adjoint block, the three

opposite-sign \(0i\) directions and the mass \(\lambda v^2\) persist.  The

compensating field is a reparametrisation of this candidate; removing the

ghosts would need structure beyond it, such as a kinetic structure for

\(\xi\) and a constraint, which is not examined here.  This concerns the

quadratic form of this candidate and does not exclude other completions.

Problems~B and~D remain open.





\paragraph{Diagonal invariance and the self-adjoint traceless bloc} \label{detail:sec-C-paragraphdiagonalinvarianceand-18}

(conditional; not closure).}

Write the real locked block as \(\Phi=v(\mathbb{I}+M)\), with

mixed-index components \(M^I{}_A\), and let the unbroken diagonal

\(\mathrm{SO}(3,1)\) of Proposition~\ref{prop:goldstone-count} act by

conjugation, \(\delta M=[\alpha,M]\) with \(\eta\alpha\) antisymmetric.

The sixteen components decompose as \(\mathbf 1\oplus\mathbf 6\oplus\mathbf 9\):

the trace, the six \(\eta\)-anti-self-adjoint components, and the nine

\(\eta\)-self-adjoint traceless components (six \(ij\) and three \(0i\)).

The nine components form the real irreducible representation \((1,1)\),

which has trivial commutant, so by Schur's lemma a quadratic form on them

that is invariant under the diagonal group is unique up to a scale.  The

form \(\mathrm{K}\) of the boost-invariant completed-action paragraph above,

restricted to this block, is such a form, with signature \((6,3)\); the

signature \((7,3)\) quoted there includes the trace direction.  Every

nonzero invariant form on the block is therefore \(c\,\mathrm{K}\), with

signature \((6,3)\) or \((3,6)\) according to the sign of \(c\).  A

numerical check (\texttt{numpy}, not a proof) found four independent

invariant quadratic forms on the sixteen components, whose restrictions to

the nine-component block span a one-dimensional space, consistent with

this count.



The statement assumes that the quadratic form is invariant under constant

diagonal \(\mathrm{SO}(3,1)\) acting by conjugation.  The displayed

potential does not have this property (see the paragraph on the assumed

group \(G_{\mathrm{int}}\supseteq\mathrm{SO}(3,1)\) in Problem~F); the

condition is a requirement on completions of the action, not a property of

the displayed action.  Under it, for a potential Hessian or for a kinetic

form with the spacetime index contracted through \(g^{(0)\mu\nu}\) only, as

in~\eqref{eq:Smicro}, the form on the nine components is indefinite unless

it vanishes, so a nonzero Hessian is not positive there, and a nonzero

kinetic form has at least three wrong-sign directions (three or six,

depending on the sign of \(c\)).  This holds for every diagonal-invariant

form without derivative--tangent mixing, which is a class of constructions

and not all of them.  Terms that tie the derivative index to the tangent

index through the vacuum identification of Equation~\eqref{eq:vev-is-frame}

add invariants and are not analysed here, so a Fierz--Pauli-type structure

is not excluded.  The trace and the anti-self-adjoint block are not

constrained by this argument.  It adds no field, symmetry or constraint,

does not identify the nine components with the tensor sector of

Chapter~\ref{ch:lorentz-decomposition}, and leaves Problem~D open.





\paragraph{Displayed action plus an assumed induced Fierz--Pauli term} \label{detail:sec-C-paragraphdisplayedactionplusan-19}

time-kinetic test (conditional; not closure).}

This paragraph keeps the displayed action as the tree-level term and

adds only the induced term assumed above; it introduces no new

tree-level structure.  On the \(\eta\)-self-adjoint block, with

\(h_{\mu\nu}\) proportional to \((\eta M)_{\mu\nu}\) as above, the

displayed kinetic term of Equation~\eqref{eq:Smicro} contracts the

components with the Euclidean form, \(\mathrm{Tr}\,M^{T}M=\sum_{\mu\nu}

h_{\mu\nu}^2\) for \(M=\eta h\), so take

\(Q=tk^2\sum_{\mu\nu}h_{\mu\nu}^2+g\,Q_{\mathrm{FP}}+Q_{\mathrm{mass}}\)

with \(t>0\), the same \(Q_{\mathrm{FP}}\) and the same sign convention

as above.  The displayed potential acts through the Hermitian part

\(H=(M+M^{T})/2\) of the Hessian paragraph in Problem~F; for \(M=\eta h\)

the \(0i\) entries of \(H\) cancel, so \(Q_{\mathrm{mass}}\propto

h_{00}^2+\sum_{ij}h_{ij}^2\) acts on seven components and the three

\(0i\) components are massless at this order.  A symbolic computation

(\texttt{sympy}, a check and not a proof) gives the

\(\omega^2\)-coefficients of the form at \(k=(\omega,0,0,0)\), which do not

depend on the mass, in the convention where healthy modes have negative

coefficient: \(-2t\) (once), \(-4t\) (three times), \(4g-2t\) (once, the

trace direction), \(-2(g+t)\) (twice) and \(-4(g+t)\) (three times).  All

ten are healthy exactly when \(t>0\) and \(-t<g<t/2\); for \(g>0\) only

\(g<t/2\) binds.  Without the mass term, the determinant at

\(k=(\omega,0,0,q)\) factors as

\(t^4(g+t)^2(q^2-\omega^2)^6\,

\bigl[(t-g)q^2-(t+g)\omega^2\bigr]^2\,

\bigl[t(t-g)(q^2-\omega^2)^2-2g^2(q^2+\omega^2)^2\bigr]\)

up to a constant.  The squared factor gives modes with

\(c^2=(t-g)/(t+g)\neq1\) for \(g\neq0\), and the final quartic factor is

also non-luminal for \(g\neq0\).  The dispersion is then

Lorentz-violating at order \(g/t\), in tension with the exact Lorentz

invariance that the locking construction of

Chapter~\ref{ch:lorentz-locking} is intended to deliver; this tension is

recorded, not resolved.



If the induced term is the linearised \((1/16\pi G_{\mathrm{eff}})\int

\sqrt{-g}\,R\), with \(\delta g_{\mu\nu}=2h_{\mu\nu}\) from

Equation~\eqref{eq:g-from-e} and the standard normalisation of the

linearised Einstein--Hilbert quadratic form (used, not re-derived), then

matching the coefficient of \(k^2h_{\mu\nu}h^{\mu\nu}\) gives

\(t=v^2\) from Equation~\eqref{eq:Smicro} and \(g=1/(16\pi G_{\mathrm{eff}})\).

For \(G_{\mathrm{eff}}>0\), the window above is then

\(1/G_{\mathrm{eff}}<8\pi v^2\), and \(1-c^2\approx 2g/t=1/(8\pi

G_{\mathrm{eff}}v^2)\) is small only if \(G_{\mathrm{eff}}v^2\gg1\).

Inserting the conditional parameterization of

Equation~\eqref{eq:Geff-structure}, the window reads

\(A\mathcal N_0\Lambda_{\mathrm{UV}}^2/v^2+B\ln(\Lambda_{\mathrm{UV}}/\mu)

+C\mu^2/v^2<8\pi\); if the logarithmic and \(\mu^2\) terms are

negligible, it reads \(A\mathcal N_0\Lambda_{\mathrm{UV}}^2<8\pi v^2\).

The coefficients \(A,B,C\), \(\mathcal N_0\), \(\mu\) and

\(\Lambda_{\mathrm{UV}}\) are undetermined (Theorem~\ref{thm:Geff} and

Problem~F), so this identifies the combination that Problem~F would have

to satisfy and does not decide whether it does.  The window has the same

structure as the condition \(Z>0\) with

\(Z=2N-3/(4\pi G_{\mathrm{eff}}v^2)\) of the Repair~II pure-trace

truncation above (a conformal-mode kinetic coefficient that turns

negative when \(G_{\mathrm{eff}}v^2\) is small), but the numerical factors

differ and no equivalence is claimed.  The result is for an assumed

linearised induced term; it is a statement about the time-kinetic matrix

and the dispersion at quadratic order, so it is not a spectrum or

unitarity result.  It does not treat the vector, ghost or \(\omega_\mu\)

sectors, does not identify the block with the tensor sector of

Chapter~\ref{ch:lorentz-decomposition}, and leaves Problems~B, C, D, F

and~G open.





\paragraph{Ledger label.} \label{detail:sec-C-paragraphledgerlabel-20}

Status: OPEN. Priority: CRITICAL PATH for the matrix-derived label.

Completion: PROGRAMME-CONDITIONAL. A scope retaining

Assumption~\ref{ass:lorentzian} gives up the claim of emergent or

matrix-derived Lorentzian signature.



Euclidean positivity of the matrix action and the formal Wick rotation of

one coordinate are recorded. They do not select Lorentzian signature without

an additional hypothesis.







\section{Problem F: coefficients (A,B,C)}

\label{sec:detailed-Problem F: coefficients (A,B,C)}



\paragraph{Ledger label.} \label{detail:sec-F-paragraphledgerlabel-21}

Status: OPEN. Priority: CRITICAL PATH. Completion:

PROGRAMME-CONDITIONAL. Reduced-scope result: treating

\(G_{\mathrm{eff}}\) as an input leaves F OPEN and does not satisfy

the matrix-derived criterion.



The conditional, truncated heat-kernel identity is

Theorem~\ref{thm:Geff}; it is not a computed \(A,B,C\)

parameterization or physical value of \(G_{\mathrm{eff}}\).





\paragraph{Quadratic Hessian of the displayed action (conditional; not closure).} \label{detail:sec-F-paragraphquadratichessianofthe-22}

Restrict to the square locked \(4\times4\) block of

Assumption~\ref{ass:locked-block}, expand~\eqref{eq:Smicro} with~\eqref{eq:V}

on the flat scaffold about \(\Phi=v\mathbb I+d\) at \(A_\mu=0\), and write

\(d=H+K\) with \(H=(d+d^\dagger)/2\) Hermitian and \(K=(d-d^\dagger)/2\)

anti-Hermitian. A direct expansion gives \(V^{(1)}=0\) and

\(V^{(2)}=\lambda v^2\,\mathrm{Tr}H^2\), independent of \(K\), and

\(\mathrm{Tr}(d^\dagger d)=\mathrm{Tr}H^2+\mathrm{Tr}(K^\dagger K)\) with no

cross term. The \(H\)--\(H\) block therefore has the Lorentzian field

equation \((\Box-\lambda v^2)H=0\) with the kinetic sign of~\eqref{eq:Smicro},

that is, the Euclidean operator \(-\nabla_E^2+\lambda v^2\) after Wick

rotation. This is the \(\Phi\)--\(\Phi\) block at fixed \(A_\mu=0\);

quadratic mixing of \(A_\mu\) with \(\Phi\) through the covariant

derivative is not evaluated, and no gauge-fixed coupled Hessian is given.

No identification of \(H\) with the \(\psi\) and \(T_{\mu\nu}\) of

Chapter~\ref{ch:lorentz-decomposition} is made here, since that split uses

\(\eta_{\mu\nu}\) and the vacuum identification of

Assumption~\ref{ass:frame-lock}. Equation~\eqref{eq:Smicro} displays no

field-strength term for \(A_\mu\), and \(\omega_\mu\) does not appear in it;

its operator is unspecified (Assumption~\ref{ass:local-lorentz}). The

displayed action also does not exhibit the Fierz--Pauli structure used as

hypothesis~5 of Theorem~\ref{thm:mT}, so any tensor-block terms beyond it

are an additional input. The gauge-fixed sector list, Goldstone and ghost

operators, measure, curvature couplings and weights are not specified by

this calculation, so it does not fix the sector sum of

Theorem~\ref{thm:Geff}, does not close F or D, and leaves the missing step

above unchanged.





\paragraph{Bookkeeping for one sector of the sum: the Hermitian scala} \label{detail:sec-F-paragraphbookkeepingforonesect-23}

block (conditional; not closure).}

This paragraph illustrates how Theorem~\ref{thm:Geff} would be applied

to one block.  It is not a derivation of that block's coefficients.

The Hessian paragraph above is evaluated on the fixed flat scaffold

\(g^{(0)}\) of Assumption~\ref{ass:aux-metric}, where \(R=0\), so the

displayed action supplies no curvature coupling for the emergent

metric.  The operator on the emergent background and the curvature

coupling \(Q\) below are the Presentation-B assumptions of

Chapter~\ref{ch:heat-kernel}, not outputs of the displayed action.



Take the Hermitian block \(H\) of the Hessian paragraph above, with sixteen real

components.  Assume (i) that this diagonal block at \(A_\mu=0\) is

treated as decoupled from the other sectors, so the coupled determinant

is not given; (ii) the determinant weight is \(w=1\) per real bosonic

component; and (iii) the operator of Equation~\eqref{eq:O-eta} on the

block is \(-\nabla^2+E\) with \(E=\lambda v^2+QR_E\), where \(Q\) is

proportional to the identity on the sixteen components (otherwise

\(1/6-Q\) below is replaced by its trace over the components).  With

\(a_1=\mathrm{tr}(R/6-E)\) of~\eqref{eq:a0a1} and the standard

curvature-linear part of

\(a_2\supset\tfrac12\,\mathrm{tr}(R/6-E)^2\), which has the same sign

convention as \(a_1\) and is used here without being re-derived, the

coefficients of \(R_E\) for this block are

\(\mathcal A_{1,R,H}=16(\tfrac16-Q)\) and

\(\mathcal A_{2,R,H}=16\lambda v^2(Q-\tfrac16)=-\lambda v^2\mathcal A_{1,R,H}\).

In the truncated, unresummed identity of Theorem~\ref{thm:Geff} with

common proper-time limits, the block would contribute

\(-\frac{\mathcal A_{1,R,H}}{32\pi^2}\bigl[(\LUV^2-\mu^2)

-\lambda v^2\ln(\LUV^2/\mu^2)\bigr]\) to \(C_R^E\); that is, in this

block's contribution the coefficient of \(\lambda v^2\ln(\LUV^2/\mu^2)\)

is minus the coefficient of \(\LUV^2-\mu^2\).  Both coefficients vanish

at \(Q=\tfrac16\); this concerns only this block's \(R_E\)-coefficient

in the truncated expansion.  Because the block has mass \(\lambda v^2\),

a resummed treatment replaces these terms by threshold functions, so

the expression is not the physical coefficient.  No comparison with

\(A,B,C\) of Equation~\eqref{eq:Geff-structure} is made.  The

anti-Hermitian block, the \(A_\mu\) and \(\omega_\mu\) sectors, the

tensor block, the ghosts, the continuation and the scheme are not

specified here, so this does not determine \(G_{\mathrm{eff}}\), does

not fix the sign of the Lorentzian term, and leaves Problem~F open.





\paragraph{The assumed group \(G_{\mathrm{int}}\supseteq\mathrm{SO}(3,1)\} \label{detail:sec-F-paragraphtheassumedgroupg_math-24}

and the displayed potential (conditional; not closure).}

Chapter~\ref{ch:lorentz-locking} takes \(G_{\mathrm{int}}\) to contain

\(\mathrm{SO}(3,1)\) acting on the internal index, with orbit

\(\delta d=v\alpha\) at the vacuum.  As mixed-index matrices

\(\alpha^I{}_J\) on the locked block, its three boost generators are

real symmetric matrices, for example \(B=E_{01}+E_{10}\), and so lie in

the Hermitian block \(H\) of the Hessian paragraph above, while the

rotations are antisymmetric.  For \(d=vB\) that paragraph gives

\(V^{(2)}=\lambda v^2\,\mathrm{Tr}H^2=2\lambda v^4\neq0\), so the

potential of~\eqref{eq:V}, which is built from the positive-definite

\(\mathrm{Tr}\,\Phi^\dagger\Phi\), is not invariant under these

generators.  Assumption~(ii) of the preceding paragraph, that the orbit

lies in the anti-Hermitian block, therefore does not hold for

\(G_{\mathrm{int}}=\mathrm{SO}(3,1)\); that count is for a different

model, and \(16-\dim G_{\mathrm{int}}\) does not apply here because the

orbit has components in \(H\).  For the displayed potential, three of

the six broken directions of Proposition~\ref{prop:goldstone-count}

carry mass \(\lambda v^2\), so that proposition's flat-Goldstone count

requires a boost-invariant potential, such as an \(\eta\)-contracted

form like the candidate discussed under Problem~D, whose quadratic form

is indefinite.  This dependence is recorded, not resolved.  Declaring

\(G_{\mathrm{int}}\) therefore does not by itself specify the

anti-Hermitian or \(A_\mu\) sectors of Problem~F: those depend on the

choice of potential, which is the open content of Problem~D.  The

statement is for the displayed action at quadratic order; it does not

rule out other constructions, does not evaluate any further sector of

Theorem~\ref{thm:Geff}, and leaves Problems~D and~F open.





\paragraph{Conditional gauge-fixed Hermitian block in the compac} \label{detail:sec-F-paragraphconditionalgaugefixed-25}

\(U(4)\) branch (not closure).}

This calculation concerns only the separate compact \(U(4)\) model

described above, not the \(\mathrm{SO}(3,1)\)-compatible construction.

Assume \(A_\mu\) is a Hermitian \(U(4)\) connection, with no

\(F_{\mu\nu}F^{\mu\nu}\) term, and write

\(\Phi=v\mathbb I+H+iY\), where \(H,Y\) are Hermitian. To quadratic

order the displayed action gives

\[

  D_\mu\Phi

  =\partial_\mu H+i(\partial_\mu Y-vA_\mu)+O(\text{quadratic}),

\]

\[

  S_{\Phi,A}^{(2)}

  =\int\mathrm{d}^4x\,\sqrt{-g^{(0)}}\,

   \operatorname{Tr}\!\left[

   -g^{(0)\mu\nu}\left(

    \partial_\mu H\,\partial_\nu H+

    (\partial_\mu Y-vA_\mu)(\partial_\nu Y-vA_\nu)

   \right)-\lambda v^2H^2\right].

\]

Thus the derivative mixing is between \(Y\) and \(A_\mu\), not \(H\)

and \(A_\mu\). Choose the linearized unitary gauge \(Y/v=0\); \(Y\)

is eliminated by this gauge condition and is not integrated as a

fluctuation in this gauge. For

\(\Phi\mapsto e^{i\alpha}\Phi\), one has \(\delta Y=v\alpha\) at the

vacuum, so the Faddeev--Popov operator for this condition is

\(\mathbb I_{16}\) there. A complex ghost pair has relative

determinant weight \(w=-2\) in the convention of

Equation~\eqref{eq:ch8-sector-sum}, but this operator is ultralocal.

In the same gauge \(A_\mu\) has only an algebraic quadratic term:

the action supplies no differential operator for it. Its determinant

and the ghost determinant are therefore not Laplace-type

heat-kernel sectors; their regulated local contributions are not

evaluated here.



For the differential Hermitian block \(H\) alone, further assume

that its operator on the emergent background is

\(-\nabla^2+\lambda v^2\), with no \(R_E\) term in \(E\), and assign

it common limits \(\LUV,\mu\). These are Presentation-B and cutoff

assumptions, not consequences of the displayed action on the fixed

scaffold. In the truncated, unresummed expansion,

\[

  \mathcal N_{0,H}=16,\qquad

  \mathcal A_{1,R,H}=\frac83,\qquad

  \mathcal A_{2,R,H}=-\frac83\lambda v^2 .

\]

The \(a_2\) coefficient uses, without re-derivation, the

curvature-linear term

\(\tfrac12\operatorname{tr}(R_E/6-E)^2\) used in the preceding

Hermitian-block bookkeeping paragraph; the expressions here are its

\(Q=0\) specialization. With the Euclidean matching convention

\(C_R^E=-1/(16\pi G_H)\), Theorem~\ref{thm:Geff} then gives the

single-block contribution

\[

  \frac1{G_H}

  =\frac{4}{3\pi}\left[

     \LUV^2-\mu^2

     -\lambda v^2\ln\!\left(\frac{\LUV^2}{\mu^2}\right)

   \right].

\]

Using \(\ln(\LUV^2/\mu^2)=2\ln(\LUV/\mu)\), this block alone

corresponds formally to

\(A_H=1/(12\pi)\), \(B_H=-8\lambda/(3\pi)\), and

\(C_H=-4/(3\pi)\) in Equation~\eqref{eq:Geff-structure}. These are

formal single-block values in a parameterization whose sectors need

not share \(\mu\); they are not full-theory coefficients. The

unresummed expansion requires \(\lambda v^2\ll\LUV^2\), so these

values are not threshold-correct. The algebraic determinants, other

fluctuation and ghost sectors, sector thresholds, continuation, and

finite curvature counterterm are not included. The displayed

nonnegative potential is not boost-invariant, so this compact branch

does not supply the missing \(\mathrm{SO}(3,1)\)-compatible sector

sum. Problems~B, C, D, F and~G remain open.





\paragraph{Which metric the scaffold one-loop action induces (conditional; not closure).} \label{detail:sec-F-paragraphwhichmetricthescaffol-26}

Consider the displayed action~\eqref{eq:Smicro} with

potential~\eqref{eq:V} on the square locked block, on the scaffold

\(g^{(0)}\) of Assumption~\ref{ass:aux-metric}, after Wick rotation,

with the internal connection held at a background value

\(\bar A_\mu\) and not integrated. Expand \(\Phi=\bar\Phi+\varphi\)

about a smooth background \(\bar\Phi\) and write \(\varphi\) in

thirty-two real components normalized so that the kinetic term is

\(\sum_a g^{(0)\mu\nu}\bar D_\mu\varphi_a\bar D_\nu\varphi_a\). The

kinetic term is bilinear in \(\Phi\), so its second variation does not

depend on \(\bar\Phi\), and the fluctuation operator is

\[

  \mathcal O_{\bar\Phi}

  =-g^{(0)\mu\nu}\bar D_\mu\bar D_\nu+E(\bar\Phi),

\]

where \(\bar D\) is built from \(g^{(0)}\) and \(\bar A_\mu\), and

\(E(\bar\Phi)\) is one half of the Hessian of \(V\) at \(\bar\Phi\);

\(E\) contains no derivative of \(\bar\Phi\). Use this operator

in~\eqref{eq:W-again} with the limits~\eqref{eq:s-limits}. For

\(I(\LUV)=\int_{\LUV^{-2}}^{\mu^{-2}}s^{-1}f(s)\,\mathrm{d}s\) one has

\(\partial_{\LUV}I=2\LUV^{-1}f(\LUV^{-2})\), so the

\(\LUV\)-dependence of \(W_E^{(1)}\) is fixed by the small-\(s\)

expansion of the heat trace through \(a_2\), up to terms that vanish as

\(\LUV\to\infty\). By~\eqref{eq:appA-a0}--\eqref{eq:appA-a2},

\[

  W_E^{(1)}\supset

  -\int\mathrm{d}^4x\,\sqrt{g^{(0)}}\left[

    \frac{32\,\LUV^4}{64\pi^2}

    +\frac{\LUV^2}{32\pi^2}

     \operatorname{tr}\Bigl(\tfrac16R[g^{(0)}]-E(\bar\Phi)\Bigr)

    +\frac{1}{32\pi^2}\ln\Bigl(\frac{\LUV^2}{\mu^2}\Bigr)a_2

  \right],

\]

where \(a_2\) contains \(E^2\), \(R[g^{(0)}]E\), curvature-squared

terms of \(g^{(0)}\), \(\bar F_{\mu\nu}\bar F^{\mu\nu}\) and total

derivatives. No \(\LUV^4\), \(\LUV^2\) or \(\ln\LUV\) term contains a derivative of

\(\bar\Phi\) outside a total derivative; terms with derivatives of

\(\bar\Phi\) have coefficients that are independent of \(\LUV\) or vanish

as \(\LUV\to\infty\). The \(\LUV^2\) curvature term

is the term~\eqref{eq:universal-R}, with \(\mathcal N_0=32\), for the

scaffold metric \(g^{(0)}\), which Assumption~\ref{ass:aux-metric}

holds fixed.



For the composite metric the corresponding coefficient vanishes. On

the flat scaffold take \(\bar A_\mu=0\) and

\(\bar\Phi=(v+\sigma)\mathbb I\), so that

\(g=(1+\sigma/v)^2\eta\) by~\eqref{eq:repairII-trace}; in the

Euclidean section \(\sqrt g\,R[g]=6(\partial\sigma)^2/v^2\) up to a

total derivative. The \(\LUV^2\) term above is

\(\tfrac{\LUV^2}{32\pi^2}\int\mathrm{d}^4x\operatorname{tr}E(\sigma)\),

with \(\operatorname{tr}E(\sigma)=16\lambda(v^2+4v\sigma+2\sigma^2)\),

and contains no \((\partial\sigma)^2\). Building the operator on \(g\)

instead, as in Assumption~\ref{ass:emergent-op}, gives

\(-\mathcal N_0\LUV^2(\partial\sigma)^2/(32\pi^2v^2)\) at the same

order, which is the induced term of~\eqref{eq:Z-sigma} with

\(1/(16\pi G_{\mathrm{eff}})=\mathcal N_0\LUV^2/(192\pi^2)\). The

replacement supplies the whole \(\LUV^2\) coefficient for the composite

metric; it does not correct a scaffold value.



The \((\partial\sigma)^2\) term of the scaffold one-loop action has no

\(\LUV^2\) or \(\ln\LUV\) piece. Take the limits~\eqref{eq:s-limits}

and \(p^2\ll\mu^2\ll\lambda v^2\). At \(\bar A_\mu=0\) the sixteen

Hermitian components have

\(E_H=\tfrac\lambda2\bigl(3(v+\sigma)^2-v^2\bigr)\) and the sixteen

anti-Hermitian components have

\(E_K=\tfrac\lambda2\bigl((v+\sigma)^2-v^2\bigr)\). At second order in

\(\sigma\) a block of \(n\) components with mass \(m\) gives

\[

  W_E^{(1)}\supset

  -\frac n4\,E'(0)^2\int\frac{\mathrm{d}^4p}{(2\pi)^4}\,

   \tilde\sigma(p)\tilde\sigma(-p)\,B(p),

\]

\[

  B(p)=\frac{1}{16\pi^2}\int_0^1\mathrm{d}x\left[

   E_1\!\left(\frac{m^2+x(1-x)p^2}{\LUV^2}\right)

   -E_1\!\left(\frac{m^2+x(1-x)p^2}{\mu^2}\right)\right],

\]

with \(E_1(x)=\int_x^\infty e^{-t}t^{-1}\,\mathrm dt\), the exponential

integral. For the Hermitian block, \(n=16\), \(E'(0)=3\lambda v\),

\(m^2=\lambda v^2\) and

\(\partial_{p^2}B|_{p=0}=-\bigl(e^{-\lambda v^2/\LUV^2}

-e^{-\lambda v^2/\mu^2}\bigr)/(96\pi^2\lambda v^2)\). For the

anti-Hermitian block, \(n=16\), \(E'(0)=\lambda v\), \(m=0\) and

\(\partial_{p^2}B|_{p=0}=(\LUV^{-2}-\mu^{-2})/(96\pi^2)\). The local

term \(\tfrac12\Delta Z\int\mathrm{d}^4x\,(\partial\sigma)^2\) therefore

adds to the microscopic value \(2N\) of \(Z\) in~\eqref{eq:Z-sigma}

\[

  \Delta Z=\frac{3\lambda}{4\pi^2}

   \Bigl(e^{-\lambda v^2/\LUV^2}-e^{-\lambda v^2/\mu^2}\Bigr)

   +\frac{\lambda^2v^2}{12\pi^2}\Bigl(\frac1{\mu^2}-\frac1{\LUV^2}\Bigr),

\]

that is, \(3\lambda/4\pi^2\) up to \(O(\lambda v^2/\LUV^2)\) and

\(O(e^{-\lambda v^2/\mu^2})\) from the Hermitian block, plus a positive

term from the anti-Hermitian block that depends on the infrared limit

\(\mu\) and is replaced by a term non-analytic in \(p^2\) as

\(\mu\to0\). Both terms are positive for \(\mu<\LUV\) and carry

no \(\LUV^2\) or \(\ln\LUV\) piece. In this one-loop calculation the

conformal-mode kinetic term keeps the sign of Lemma~\ref{lem:msigma}.

The opposite-sign term of~\eqref{eq:Z-sigma} is not produced in this

calculation; in the manuscript it arises from the induced

Einstein--Hilbert term, that is, from the operator built on the

emergent metric.





\paragraph{What the two preceding results imply for Problem~F (conditional; not closure).} \label{detail:sec-F-paragraphwhatthetwoprecedingre-27}

In the class of the preceding paragraph, neither the displayed action

nor its quantum corrections give the composite-metric fluctuation a

linearized diffeomorphism invariance, an induced Einstein--Hilbert

term, or the Fierz--Pauli block of hypothesis~5 of

Theorem~\ref{thm:mT}; the induced Fierz--Pauli term assumed in the

quadratic tests of Section~\ref{sec:status-D} is likewise not

generated in this way. These structures need an ingredient that ties

derivative indices to tangent indices. The

operator~\eqref{eq:O-eta} on the emergent background, used in

Theorem~\ref{thm:Geff}, is such an ingredient; it enters as

Assumption~\ref{ass:emergent-op}, which is assumed, not derived. The

one-loop \(\Phi\)--\(\Phi\) action of the displayed action on the

scaffold contributes to the \(\LUV^2\) and \(\ln\LUV\) coefficients

\(A\) and \(B\) of~\eqref{eq:Geff-structure} only through the curvature

of \(g^{(0)}\), not through that of the composite metric; the missing

step of Problem~F includes discharging

Assumption~\ref{ass:emergent-op}, or replacing it by another derived

ingredient. Not covered here: integration of the auxiliary \(A_\mu\) in

the first result, the hypotheses relaxed in the second, and

Presentation~A in general. No value or sign of \(G_{\mathrm{eff}}\)

follows. Problems~B, C, D and~F remain open.





\paragraph{Presentation~A in the Poisson limit: what four matrix coordinates can supply (conditional; not closure).} \label{detail:sec-F-paragraphpresentationainthepoi-28}

Keep the setting of the preceding paragraph: four matrix coordinates, the Poisson limit, matrix indices contracted with \(\eta=\mathrm{diag}(-1,1,1,1)\), a real nondegenerate \(\Theta\), \(K^{\mu\nu}=\eta_{\alpha\beta}\Theta^{\alpha\mu}\Theta^{\beta\nu}\), \(G=K^{-1}\), \(\rho=1/|\operatorname{Pf}\Theta|\), \(g_\rho^{\mu\nu}=\rho K^{\mu\nu}\). Let \(N=\Theta\eta\in\mathfrak{so}(3,1)\) with electric and magnetic parts \(\mathbf E,\mathbf B\), so \(\operatorname{Pf}\Theta=\pm\mathbf E\cdot\mathbf B\). For vectors \(a,b\) write \(a\circ b:=ab^T+ba^T\).



\emph{Image of \(\Theta\mapsto G\).} The characteristic polynomial of \(N\) is \(z^4+(\mathbf B^2-\mathbf E^2)z^2-(\mathbf E\cdot\mathbf B)^2\), so \(K\eta=-N^2\) has two doubly degenerate eigenvalues and \((K\eta)^2-\tfrac12\operatorname{tr}(K\eta)\,K\eta-(\operatorname{Pf}\Theta)^2\,\mathbb I=0\). Hence the six-dimensional image has codimension four in the ten-dimensional space of metrics and lies in the set cut out by this relation. Decomposing \(N=(n_+,n_-)\) in \(\mathfrak{sl}_2\oplus\mathfrak{sl}_2\), \(N^2=(\nu_++\nu_-)\mathbb I+2\,n_+\otimes n_-\), and a kernel element of the differential of \(\Theta\mapsto K\) must satisfy \(\delta n_\pm=\mp\lambda n_\pm\) and \(\lambda(\nu_+-\nu_-)=0\); since \(\operatorname{Im}\nu_+=0\) only if \(\mathbf E\cdot\mathbf B=0\), the map \(\Theta\mapsto G\) has rank six at every \(\Theta\) with \(\operatorname{Pf}\Theta\neq0\). The differential of \(\Theta\mapsto K\) has rank five when \(\Theta\neq0\) and \(\operatorname{Pf}\Theta=0\) (symbolic and exact-rational checks).



\emph{Linear response.} Let \(\Theta=\bar\Theta+\epsilon T\), \(T^{\mu\nu}=\bar\Theta^{\mu\rho}\partial_\rho a^\nu-\bar\Theta^{\nu\rho}\partial_\rho a^\mu\), \(a^\mu=u^\mu\cos(p\cdot x)\), \(q=\bar\Theta p\). The amplitude of \(\sin(p\cdot x)\) in \(\delta G\) is \((\bar Gu)\circ p+n\circ m\), with \(n=\bar G\bar\Theta p\), \(m=\bar G\bar\Theta\eta u\), and \(\delta\ln\rho=p\cdot u\). For \(u\propto q\) one has \(T=0\) (the matrix gauge). The first term is a linearised diffeomorphism. Since \(n\) and \(p\) are covectors, orthogonality is taken with the metric \(\bar K\): \(n^T\bar Kp=0\), while \(n\) is not parallel to \(p\) for generic \(p\), so the second term is not a diffeomorphism. The images of \(u\mapsto(\delta G,\delta\rho)\) and \(u\mapsto(\delta g_\rho,\delta\rho)\) are three-dimensional. For generic \(p\) the linearised Einstein operator about the constant background has rank three on this image, and its image lies in the six-dimensional space of symmetric tensors with \(p^\mu T_{\mu\nu}=0\) (``conserved''; exact-rational check for \(G\) and for \(g_\rho\)). Hence the tensor equation \(\mathcal E[h]=T\), and any linear local relation between \((\delta G,\delta\rho)\) and \(T\), can hold only for \(T\) in a subspace of dimension at most three of the six-dimensional space of conserved sources; varying the four functions \(a^\mu\) gives only the projection of the equation onto that image. This does not exclude an additional conserved stress that absorbs the codimension; such a stress would have to be derived. As a counting argument (not a derived result), four coordinate functions minus one matrix gauge function leave three gauge-invariant functions, against \(10-4+1=7\) for a diffeomorphism-invariant theory of a metric and one scalar. The tree-level quadratic action of Remark~\ref{rem:ch3-trace-quadratic} is proportional to \(\int\rho\,\eta_{\mu\alpha}\eta_{\nu\beta}T^{\mu\nu}T^{\alpha\beta}\), algebraic in \((\delta G,\delta\rho)\); every term with derivatives of the metric perturbation is a loop term. Whether the algebraic term is cancelled is the content of Assumption~\ref{ass:EE} and is not decided here.



\emph{Metric seen by the operators.} Suppose, as an assumed dictionary not contained in the four-matrix class, that there are \(D\ge4\) coordinate matrices \(X^I\) with a Lorentzian \(\eta_{IJ}\), and put \(\Phi^I{}_\mu:=\partial_\mu x^I\), \(g:=\Phi^T\eta\Phi\). Then \(K=\theta^{T}g\,\theta\) exactly; for \(D=4\), \(x^I=y^I\) and \(g=\eta\). The relation \(K=c\,g^{-1}\), \(c\neq0\), would give \(N^2=-c\) with \(N=\theta g\in\mathfrak{so}(g)\), whose characteristic polynomial has constant term \(-(\mathbf E\cdot\mathbf B)^2\le0<c^2\); \(c=0\) contradicts invertibility. So, pointwise for any real Lorentzian \(g\), the operator metric \(G=K^{-1}\) is not \(g\); this extends Remark~\ref{rem:ch3-moyal-check}, and Assumption~\ref{ass:emergent-op} read with the metric of that embedding reading is not discharged. The metric \(G\) itself has the form~\eqref{eq:g-from-e} for the coframe of the preceding paragraph. Let \(\ell\) be the fixed length scale in \(\Theta/\ell^2\), and \(e=\eta^{-1}\hat\omega\) with \(\hat\omega=(\Theta/\ell^2)^{-1}\). If \(e\) is identified with \(\langle\Phi\rangle/v\) as in~\eqref{eq:vev-is-frame}, then, with the internal index contracted with \(\delta\) and the locked block real, \(e^Te=\mathbb I\) on the vacuum manifold of~\eqref{eq:V} holds iff \(\hat\omega^2=-\mathbb I\), i.e. \(J:=\hat\omega\) is an orthogonal complex structure. Then \(G=\mathbb I-2nn^T\) with \(n=-Je_0\), where \(e_0\) is the first basis vector, so \(n\) is a spatial unit vector, and only two of the six directions of \(\delta\Theta\) are tangent to that manifold.



This is a statement for four matrix coordinates, the Poisson limit, a constant \(\eta\), linear order about constant \(\bar\Theta\) and generic \(p\); it does not apply to an independent field \(\Phi\) of Presentation~B, to the embedding statement above beyond four coordinates, or to Moyal corrections. It does not derive \(\langle\Phi\rangle/v\) as a coframe and gives no value or sign of \(G_{\mathrm{eff}}\). Problem~F, and Problems~A, B, C and~D, remain open.





\paragraph{Leaf dimension above four: fibre averages, a product factor and a conjugate-coordinate class (conditional; not closure).} \label{detail:sec-F-paragraphleafdimensionabovefou-29}

Work in the Poisson limit with matrix coordinates \(x^I\), \(I=0,\dots,D-1\), contracted with \(\eta=\mathrm{diag}(-1,1,\dots,1)\), on a symplectic leaf of dimension \(2m\ge4\), and four base functions \(y^\mu\) on the leaf. For a function \(\phi(y^\mu)\) the kinetic tensor is \(K^{\mu\nu}=\langle\eta_{IJ}V^{I\mu}V^{J\nu}\rangle\) with \(V^{I\mu}=\{x^I,y^\mu\}\) and \(\langle\cdot\rangle\) a positive average over the fibre; for a four-dimensional leaf it is \(\theta^{T}g\theta\).



\emph{Matrices as base coordinates.} Let \(y^\mu=x^\mu\), \(\mu=0,\dots,3\), so that \(V^{I\mu}=\Theta^{I\mu}\) is real and antisymmetric. Then \(K^{00}=\langle\sum_{I\ge1}(\Theta^{I0})^2\rangle\ge0\) and \(\sum_\mu K^{\mu\mu}=\langle\sum_{I\ge4}(\Theta^{I0})^2+\sum_{I\ge1,\,k\ge1}(\Theta^{Ik})^2\rangle\ge0\), whereas \(K=c\,\eta^{-1}\) has \(K^{00}=-c\) and \(\sum_\mu K^{\mu\mu}=2c\); hence \(c=0\). The case \(K=0\) forces \(\Theta^{I\mu}=0\) for \(\mu\le3\), so the \(x^\mu\) would be Casimir functions, contradicting nondegeneracy. Therefore \(K\ne c\,\eta^{-1}\) for every \(c\ne0\), for every \(D\), every compact fibre and every positive average: in this class, enlarging the leaf by a fibre does not remove the proportionality obstruction of the preceding paragraphs. For any nondegenerate real Lorentzian \(g_L\) on a leaf of dimension \(2m\ge4\), \(G_L=(\theta^{T}g_L\theta)^{-1}\) is not proportional to \(g_L\). Indeed \(K=c\,g_L^{-1}\) with \(c\neq0\) gives \(N^2=-c\,\mathbb I\) for \(N=\theta g_L\). For \(c>0\), \(\det N=c^m>0\), whereas \(\det N=(\operatorname{Pf}\theta)^2\det g_L<0\). For \(c<0\), \(N\) is \(g_L\)-skew with \(N^2=|c|\,\mathbb I\), so its \(\pm\sqrt{|c|}\) eigenspaces are totally null and have dimension at most one in Lorentzian signature, which forces \(2m\le2\); for \(2m=2\) the relation can hold (\(g_L=\mathrm{diag}(-1,1)\), \(\theta=2\epsilon\) gives \(K=-4g_L^{-1}\)). For a product leaf \(M\times\mathcal C\) with block-diagonal \(\theta\) and embedding, \(V^{i\mu}=0\) for the coordinates along \(\mathcal C\), so \(K\) is the tensor of the \(M\) factor, and to first order the response of \(K\) to \(\delta x^I=A^I(t)e^{ip\cdot y}\) is the four-dimensional one with amplitude \(\langle A^\mu\rangle\); transverse amplitudes and non-constant fibre parts do not enter. The rank statements of the preceding paragraph therefore hold with \(D\) the number of coordinates that vary over \(M\), and Kaluza--Klein modes enter only quadratically.



\emph{Conjugate coordinates.} The proportionality is not excluded when the base functions are not among the matrix coordinates. Let \(M^{AB}\), \(A,B=0,\dots,5\), generate \(\mathfrak{so}(4,2)\) with invariant form \(\mathrm{diag}(-,+,+,+,+,-)\) on the coadjoint orbit \(\mathbb{CP}^{1,2}\) (six-dimensional leaf), take the four matrix coordinates \(T^a=M^{a4}/R\), \(a=0,\dots,3\), and the base functions \(y^\mu=M^{\mu5}/R\). The functions \(y^\mu\) are not among the \(T^a\), and the leaf is not the four-dimensional leaf of Definition~\ref{def:matrix-coords-ext}. The Lie--Poisson relation gives \(\{T^a,y^\mu\}=\pm\eta^{a\mu}M^{45}/R^2\), because the form has \(\eta^{45}=\eta^{a5}=\eta^{4\mu}=0\), hence \(K^{\mu\nu}=(M^{45}/R^2)^2\eta^{\mu\nu}\) pointwise, with \(M^{45}\) a function of \(y^\mu\) on the orbit (checked on an explicit realisation). Here \(\eta^{\mu\nu}\) is the Cartesian form of the functions \(y^\mu\), not the pull-back of the matrix contraction \(\eta_{IJ}\), and the matrix-induced tensor on the six-dimensional leaf has rank four. This \(K\) is conformally flat, so the exact class carries one scalar function and no metric waves; whether fluctuations of \(T^a\) that are not fibre-constant supply a sufficient number of gauge-invariant functions is not computed here, and no argument is given that the \(y^\mu=M^{\mu5}/R\) are the base functions of the reduced theory.



\emph{Heat kernel over a fuzzy factor.} For the product leaf with a finite fuzzy factor, the scalar operator of the Poisson-limit paragraph is \(\mathcal O_s\otimes1+1\otimes\Box_{\mathcal C}\), so each four-dimensional sector is replicated over the \(\dim\mathrm{Mat}(\mathcal H_{\mathcal C})\) modes, the mode of mass \(m\) with the threshold factor \(F(m^2/\Lambda_{\mathrm{UV}}^2)=\int_1^\infty u^{-2}e^{-m^2u/\Lambda_{\mathrm{UV}}^2}\mathrm du\in(0,1]\) multiplying the coefficient of Theorem~\ref{thm:Geff}. In the truncated expansion the curvature coefficient is thus of order \(\Lambda_{\mathrm{UV}}^2\) times the weighted number of light modes; it is not set by the Kaluza--Klein scale. The supersymmetric one-loop mechanism of~\cite{Steinacker2023}, in which the corresponding terms cancel and the Newton constant is set by a finite mixed term, rests on field content that is not part of the present sector list.



This is a statement for the Poisson limit, the stated fibre classes and off-shell backgrounds; it does not identify \(G\) or \(\eta_{\mu\nu}\) with the emergent metric of Chapter~\ref{ch:vierbein-emergence}, gives no field equation, gap of the transverse scalars, \(G_{\mathrm{eff}}\) value or sign, and does not compute the rank of the linearised Einstein operator in the conjugate-coordinate class.





\paragraph{Conjugate-coordinate class: first-order response and linearised Einstein rank (conditional; not closure).} \label{detail:sec-F-paragraphconjugatecoordinatecl-30}

This paragraph supplements the preceding one, which left the rank of the linearised Einstein operator in this class uncomputed. Keep the conjugate-coordinate class of the preceding paragraphs in the Poisson limit: the leaf \(\mathbb{CP}^{1,2}\) with the Lie--Poisson bracket \(\{M^{AB},M^{CD}\}=s(\eta^{AC}M^{BD}-\eta^{AD}M^{BC}-\eta^{BC}M^{AD}+\eta^{BD}M^{AC})\), \(s=\pm1\) a convention, the matrices \(X^a=T^a+A^a\) with \(A^a\) arbitrary functions on the leaf, \(\eta=\mathrm{diag}(-1,1,1,1)\), and the base functions \(y^\mu\) held fixed. Put \(m=M^{45}/R^2\) and \(y^4=M^{45}/R\). Then \(\{T^a,y^\mu\}=s\,\eta^{a\mu}m\), \(\{y^\mu,y^\rho\}=-s\,M^{\mu\rho}/R^2\), and with \(A^\nu:=A^a\delta_a{}^{\nu}\) the first-order change of \(K^{\mu\nu}=\eta_{ab}\{X^a,y^\mu\}\{X^b,y^\nu\}\) is \(\delta K^{\mu\nu}=s\,m\,(\{A^\nu,y^\mu\}+\{A^\mu,y^\nu\})\). Additional matrices \(X^i\) with \(\{X^i,y^\mu\}=0\) in the background enter \(K\) only through \(\eta_{ij}\{X^i,y^\mu\}\{X^j,y^\nu\}\), which is quadratic in their fluctuations.



\emph{Fibre structure.} On an explicit realisation one finds \(M^{ab}y_b=0\) for \(a,b=0,\dots,4\) (indices lowered with \(\mathrm{diag}(-,+,+,+,+)\)), hence \(T^\mu=-M^{\mu\sigma}y_\sigma/y^4\) and \(y_\mu T^\mu=0\). Over a point of the base with \(y^4\neq0\) the fibre of \(y^a\) is a two-sphere on which the six functions \(M^{\mu\rho}\) have zero average, are first spherical harmonics and span a three-dimensional space \(W_y\), the projection to the indices \(0,\dots,3\) of the self-dual two-forms of the Euclidean space \(y^\perp\subset\mathbb R^{1,4}\); the four functions \(T^\mu\) span the same space of functions. Consequently the first-harmonic part of \(A^\nu\) consists of twelve amplitudes \(b^{\nu;\mu\rho}(y)\in W_y\) and coincides with the frame-type fluctuations \(A^\nu=\kappa^\nu{}_\sigma(y)T^\sigma\), defined modulo \(\kappa^\nu{}_\sigma\to\kappa^\nu{}_\sigma+k^\nu y_\sigma\).



\emph{Response.} Let \(\mathrm d\mu\) be the fibre measure induced by the Liouville measure on a sheet \(y^4>0\) (the other sheet is not included) and \(B^{\nu;\mu\rho}=\int\mathrm d\mu\,A^\nu M^{\mu\rho}\). Liouville invariance of the bracket gives \(\int\mathrm d\mu\,\delta K^{\mu\nu}=(m/R^2)\,\partial_\rho\bigl(B^{\nu;\mu\rho}+B^{\mu;\nu\rho}\bigr)\), so the fibre-integrated variation \(\delta\mathcal K^{\mu\nu}\) satisfies \(\partial_\mu\partial_\nu(\delta\mathcal K^{\mu\nu}/m)=0\) for every \(A\). Only the first-harmonic part of \(A\) contributes: fibre-constant fluctuations, the ordinary gauge fields, and harmonics of order two and higher give no linear response. Take a plane wave \(e^{iq\cdot y}\), frozen coefficients at a base point, \(|q|\ell\gg1\) with \(\ell^2=-y_ay^a\), and \(q^2\neq0\). To leading order \(\delta G_{\mu\nu}\propto\eta_{\mu\alpha}\eta_{\nu\beta}h^{\alpha\beta}\) with \(h^{\mu\nu}=q_\rho(b^{\nu;\mu\rho}+b^{\mu;\nu\rho})\). Antisymmetry in \((\mu,\rho)\) gives \(q_\mu q_\nu h^{\mu\nu}=0\), so the image has dimension at most nine. The rank is a consequence of the image being the whole hyperplane \(q_\mu q_\nu h^{\mu\nu}=0\) and does not depend on the particular three-dimensional space \(W_y\); this hyperplane is not invariant under longitudinal diffeomorphisms, so the gauge group of the class is that of divergence-free diffeomorphisms. In exact rational arithmetic, at base points \((y^0,y^4)=(5,3),(13,5),(5,4),(25,7)\) in units \(R=1\) and for generic rational \(q\), the image has dimension nine, the linearised Einstein operator of \(\eta\) has rank six on it, its image is the whole six-dimensional space of conserved tensors \(q^\mu T_{\mu\nu}=0\), and its kernel on the image is three-dimensional. Six is the largest possible rank, so it is the generic rank by lower semicontinuity. For null \(q\) the rank is three, and at \(y^4=0\) the image has dimension seven and the rank is five. For \(A^\nu=\{\Lambda,T^\nu\}\) the leading image is \(h=q\circ(\Lambda q)\) with \(\Lambda^{\mu\rho}\in W_y\), the three-dimensional space of transverse linearised diffeomorphisms \(\xi\cdot q=0\), which is the kernel found above; the longitudinal diffeomorphism is neither a matrix gauge transformation nor in the image. The twelve amplitudes thus split into three that leave \(G\) unchanged, three gauge directions and six gauge-invariant functions, against six conserved sources and against \(10-4+1=7\) for a metric and a density in a diffeomorphism-invariant theory; the density is fixed by the Liouville measure and the image obeys the one relation above. A change \(y^\mu\to y^\mu+\delta y^\mu\) of the base functions acts at leading order as a linearised diffeomorphism with parameter \(\int\mathrm d\mu\,\delta y^\mu\) on the fibre-integrated tensor and the fibre volume, and so adds no gauge-invariant function. The algebraic parts of the response and of the gauge variation are of relative order \(1/(|q|\ell)\) and are not computed, so the gauge invariance of \(\mathcal E[h]=T\) holds at the leading symbol only.



This is a kinematic statement for the Poisson limit, the stated class, linear order, frozen coefficients and the fibre-integrated tensor; fibre-non-constant parts of \(\delta K\) are not treated. It does not identify \(G\) with the emergent metric of Chapter~\ref{ch:vierbein-emergence}, derives no field equation, and does not show that the first-harmonic amplitudes are light, positive-norm or decoupled. A Yang--Mills-type tree-level quadratic action, analogous to that of Remark~\ref{rem:ch3-trace-quadratic} and built from first derivatives of \(A\), gives equations of second order in these amplitudes (no such action is computed in this class), whereas \(\mathcal E[h]\) is of third order; an equation \(\mathcal E[h]=T\) would therefore require an induced term or another structure, which is not decided here. Fibre harmonics of the fluctuation matrices form the higher-spin tower described in~\cite{SperlingSteinacker2019,MantaSteinacker2025}; the relation of the amplitudes \(b^{\nu;\mu\rho}\) to that tower is not examined here.





\paragraph{Conjugate-coordinate class: tree-level quadratic form of the level-one amplitudes (conditional; not closure).} \label{detail:sec-F-paragraphconjugatecoordinatecl-31}

This paragraph supplements the preceding ones, in which the tree-level action of the amplitudes \(b^{\nu;\mu\rho}\) was not computed. Keep the Poisson limit, the leaf \(\mathbb{CP}^{1,2}\), \(X^a=T^a+A^a\), the base functions \(y^\mu\) and the notation of the preceding paragraphs, in units \(R=1\), and write \(D^a=\{T^a,\cdot\}\), \(\mathcal J_{ab}=\{M_{ab},\cdot\}\), \(\Box=D_bD^b\). The quartic term of~\eqref{eq:Smatrix} with \(C_{\mu\nu\lambda}=0\), read with \([\cdot,\cdot]\to i\{\cdot,\cdot\}\) and \(\operatorname{Tr}\to\int\mathrm{vol}\), is \(S=(4g_{\mathrm{YM}}^2)^{-1}\int\mathrm{vol}\,\{X^a,X^b\}\{X_a,X_b\}\), with field equation \(\{X^b,\{X_b,X_a\}\}=0\).



\emph{Background.} The Lie--Poisson relations give \(\{T_a,\{T^a,T^b\}\}=-3T^b\), an identity of the polynomial Poisson algebra of \(\mathfrak{so}(4,2)\) that does not depend on the sign \(s\). Hence \(T^a\) is not stationary for the quartic term alone. It is stationary for \(S+(2g_{\mathrm{YM}}^2)^{-1}\mu^2\int\mathrm{vol}\,X_aX^a\), whose field equation is \(\{X^b,\{X_b,X_a\}\}=\mu^2X_a\), exactly when \(\mu^2=-3\). A constant \(C_{\mu\nu\lambda}\) cannot take the place of this term: it would add \(C_{abc}M^{bc}\) to the field equation, and these functions are linearly independent of \(T_a=M_{a4}\) on the orbit. The bilinear term is not part of~\eqref{eq:Smatrix}; it is an input of this class with a coefficient fixed by stationarity, and a mass term is reported for the covariant backgrounds of~\cite{SperlingSteinacker2019}.



\emph{Quadratic form.} With this term, \(S^{(2)}=(2g_{\mathrm{YM}}^2)^{-1}\int\mathrm{vol}\,A^aL(A)_a\) with \(L(A)_a=-\Box A_a+D_a(D_bA^b)-2s\,\mathcal J_{ab}A^b+\mu^2A_a\), an identity verified in the polynomial Poisson algebra for both signs \(s\). Gauge covariance of the field equation about a stationary point gives \(L(\{\Lambda,T\})=0\) for every function \(\Lambda\), and the same check fails for \(\mu^2\neq-3\). The quadratic form is therefore exactly invariant under \(\delta A^a=\{\Lambda,T^a\}\), not only at a leading symbol.



\emph{Leading symbol.} Since \(D^ay^\mu=s\,m\,\eta^{a\mu}\) with \(m=M^{45}\) a function of \(y\) alone, the only part of \(L\) with two derivatives along \(y\) is \(-\Box A_a+D_a(D_bA^b)\). The operator \(\mathcal J_{ab}\) acts on \(y^\mu\) as a Lorentz rotation with a coefficient linear in \(y\), and the fibre derivatives of \(A\) and the derivatives of \(m\) enter with coefficients of relative order \(1/(|q|\ell)\). For \(A^a=c^a(y)F\) with \(F\) a fixed fibre function and \(|q|\ell\gg1\), then, \(L(A)_a=-m^2(\partial^2c_a-\partial_a\partial_bc^b)F\) plus terms with at most one derivative of \(c\), so that

\[

S^{(2)}\simeq\frac{1}{4g_{\mathrm{YM}}^2}\int\mathrm{vol}\;m^2\,\mathcal F_{ab}\mathcal F^{ab},\qquad

\mathcal F_{ab}=\partial_aA_b-\partial_bA_a ,

\]

with \(\partial_a\) the derivative along \(y\) at fixed fibre argument. At this order each fibre function is an independent Maxwell field with the same coefficient, and there is no coupling between fibre harmonics. On the explicit realisation of the preceding paragraphs, evaluating \(L\) on \(A^a=c^a\mathrm e^{i\lambda q\cdot y}M^{12}\) reproduces \(\lambda^2m^2M^{12}\mathrm e^{i\lambda q\cdot y}(q^2c_a-q_aq\cdot c)\) with a relative remainder that falls as \(1/\lambda\).



\emph{Level-one sector.} Write \(M^{\mu\rho}=\sum_ie^iS_i^{\mu\rho}\) with \(S_i\) an orthogonal, equal-norm basis of \(W_y\) and \(e^i\) the three first spherical harmonics on the fibre; by the \(\mathrm{SU}(2)\) invariance of the fibre measure, \(\int\mathrm d\mu\,e^ie^j=\varkappa(y)\,\delta^{ij}\) with \(\varkappa>0\). For \(A^\nu=c^\nu_i(y)e^i\) one has \(b^{\nu;\mu\rho}=\varkappa\,c^\nu_iS_i^{\mu\rho}\), and the twelve amplitudes are three vector fields \(c_i\) with the action \((m^2\varkappa/4g_{\mathrm{YM}}^2)\sum_i\int\mathrm d^4y\,\mathcal F_{i,ab}\mathcal F_i^{ab}\). The kinetic matrix is \(\delta_{ij}(q^2\eta_{ab}-q_aq_b)\), definite in the index \(i\); the leading symbol has no mass term and no coupling to the other harmonics. The overall sign is that of the quartic term as displayed, common to all fibre harmonics including the fibre-constant gauge field, so no relative ghost arises at this order. It is opposite to the sign of the Lorentzian Maxwell action \(-\tfrac14F_{ab}F^{ab}\) in the same \(\eta\), and whether it is the healthy sign depends on the convention for the weight, which~\eqref{eq:Smatrix} does not fix. The Maxwell transformation \(\delta c_i^\nu=\partial^\nu\lambda_i\) is the leading form of \(\delta A^\nu=\{\Lambda,T^\nu\}\) with \(\Lambda=\lambda_ie^i\), under which \(h\) changes by a divergence-free linearised diffeomorphism. Masses, the connection induced by the \(y\)-dependence of \(S_i\) and \(m\), and the mixing of harmonics enter at relative orders \(1/(|q|\ell)\) and \(1/(|q|\ell)^2\), which frozen coefficients do not determine; they are not computed.



This is a statement for the Poisson limit, the stated class with the bilinear term, and the leading symbol. It does not determine masses, the signs of subleading terms or decoupling beyond leading order, and it does not identify the fibre-averaged tensor with the emergent metric of Chapter~\ref{ch:vierbein-emergence}.





\paragraph{Conjugate-coordinate class: matrix Noether identity and the divergence of the stress (conditional; not closure).} \label{detail:sec-F-paragraphconjugatecoordinatecl-32}

Work in the Poisson limit with \(F^{ab}=\{X^a,X^b\}\) and \(E_c=-\{X^b,\{X_b,X_c\}\}+\mu^2X_c\), the field equation of \((4g_{\mathrm{YM}}^2)^{-1}\int\mathrm{vol}\,(F_{ab}F^{ab}+2\mu^2X_aX^a)\), for arbitrary functions \(X^a\) on the leaf. The Jacobi identity and the Leibniz rule give \(\{X^a,E_a\}=0\) identically. The tensor \(\mathbb T^{ab}=F^{ac}F^{b}{}_{c}+\mu^2X^aX^b\), which is the variation of the action with respect to the contraction form, satisfies \(\{X_a,\mathbb T^{ab}\}=-F^{b}{}_{c}E^c+\tfrac14\{X^b,F_{cd}F^{cd}\}\); its divergence therefore vanishes on the field equation up to a gradient, and \(\mathbb T^{ab}-\tfrac14\eta^{ab}F_{cd}F^{cd}\) is divergence-free on the field equation exactly, the subtracted term being \(\tfrac14\eta^{ab}F_{cd}F^{cd}\) and not a trace of \(\mathbb T^{ab}\). About \(X^a=T^a\) the linear form obeys \(\{T^a,L(A)_a\}=-(3+\mu^2)\{A^a,T_a\}\), which vanishes exactly for \(\mu^2=-3\). For the massless model the matrix-level identity \([X_a,\mathbb T^{ab}]\propto\) equation of motion is standard; the Poisson-limit form \(\{X^a,E_a\}=0\) and the divergence identity with the \(\mu^2\) term are recorded here as exact identities in the polynomial Poisson algebra, and were verified there for both signs of the bracket. The tensor \(\mathbb T^{ab}\) carries matrix indices and \(\{X_a,\cdot\}\) is a frame derivative, not \(\nabla_\mu\mathbb T^{\mu\nu}\); it differs from the coordinate divergence of the fibre-averaged tensor by terms containing derivatives of the frame coefficients, of relative order \(1/(|q|\ell)\), which are not computed. The paragraph does not identify \(\mathbb T^{ab}\) with the source of an Einstein-type equation.





\paragraph{Conjugate-coordinate class: the \(\Theta\)-built heat kernel and the induced coefficient (conditional; bounded no-go).} \label{detail:sec-F-paragraphconjugatecoordinatecl-33}

Keep the Poisson limit, the leaf \(\mathbb{CP}^{1,2}\), \(X^a=T^a+A^a\), the bilinear term with \(\mu^2=-3\) and the notation of the preceding paragraphs; \(D^a=\{T^a,\cdot\}\), \(J^{AB}=\{M^{AB},\cdot\}\), \(\Box=\eta_{ab}D^aD^b\). The regulator is the proper-time cutoff \(s\ge\LUV^{-2}\) of the operator \(\Box_X\) itself, applied to every sector; it is built from Poisson brackets and the matrix contraction only. Background gauge fixing \(D_aA^a\) gives the ghost operator \(-\Box\) (weight \(-2\)), the vector operator \(H_V=-\Box-2sJ_{ab}\Sigma^{ab}\) (the same \(J_{ab}A^b\) as in the quadratic form above, with \(\mu^2\) added), and, for \(n_{\mathrm{ex}}\) further matrices with trivial background, \(-\Box+\mu_i^2\) with \(\mu_i^2\) an additional input. Fermions are absent. In exact arithmetic in the polynomial Poisson algebra, for both signs \(s\): \(\Box=C_{\mathfrak{so}(4,1)}-C_{\mathfrak{so}(3,1)}\) with the generators \(J^{AB}\), \(A,B\le4\), and \(J^{ab}\), \(a,b\le3\); \([\Box,J^{ab}]=0\) while \([\Box,D^a]\ne0\); and on functions of \(y^A=M^{A5}\), \(A\le4\), one has \(D^a=s\,(y^4\partial^a-y^a\partial_4)\), a rotation field of the hyperboloid \(y\cdot y=-\ell^2\). Fibre-constant functions therefore form an invariant subspace of \(\Box\), on which \(\Box=\rho^{-1}\partial_\mu(\rho\,m^2\eta^{\mu\nu}\partial_\nu\,\cdot)\) with \(m=y^4\), \(\rho=1/m\); for \(G^{\mu\nu}=m^2\eta^{\mu\nu}\) the Seeley--DeWitt density of this sector is \(\sqrt{|G|}\,(R_G/6-E)=(5m^2-3\ell^2)/(4m^6)\), which is not of definite sign (a formal evaluation for \(\eta=\mathrm{diag}(-1,1,1,1)\); no Euclidean heat-kernel statement is made for this sector). Fibre harmonics of order \(\ge1\) are not an invariant subspace of the perturbed operators, so this sector does not by itself give the response to the level-one amplitudes.



\emph{Hypoelliptic structure.} At every sampled point of the orbit (random points, and the sequence \(y^4\to0\)) the fields \(D^a\) have rank four and \(D^a\), \(J^{ab}\) together span the six-dimensional tangent space; the horizontal rank drops on \(y^4=0\) (one singular value proportional to \(y^4\)), which is excluded. Hence, for the Euclidean contraction \(\eta\to\delta\) (the operator \(\Box_E=\sum_a(D^a)^2\), which is not the Casimir difference above), the operator is a sub-Laplacian of step two with homogeneous dimension \(4+2\cdot2=8\), not an elliptic operator on a product; the ranks, the step and the dimension \(8\) depend only on the fields \(D^a\), \(J^{ab}\). The ghost and vector operators below are then built from \(\Box_E\). Under equiregularity the diagonal of its heat kernel is \(K(s)=c_0s^{-4}(1+O(s))\) with \(c_0\) given by the nilpotent approximation~\cite{BenArous1989}, here a group with four horizontal generators and a two-dimensional centre. After Fourier transform along the centre (variable \(\zeta\in\mathbb R^2\)) the nilpotent operator is a four-dimensional charged Laplacian in a constant field \(F_{ab}(\zeta)=\zeta_kS^{ab}{}_k\) of the two-plane strengths \(f_1,f_2\propto|\zeta|\), with Landau levels \(f_1(2n_1+1)+f_2(2n_2+1)\) and the diagonal \(\prod_j f_j/(4\pi\sinh sf_j)\) per unit horizontal volume and per \(\mathrm d^2\zeta/(2\pi)^2\). With \(\nu\) the Liouville volume of the frame \((D^0,\dots,D^3,Z_1,Z_2)\) of the nilpotent algebra, \(c_0=\nu^{-1}(64\pi^4)^{-1}\!\int_0^{2\pi}\!\mathrm d\varphi\,\phi_1\phi_2\!\int_0^\infty\!\mathrm du\,u^3/(\sinh u\phi_1\sinh u\phi_2)\), \(f_j=|\zeta|\phi_j\); for \(f_1=f_2\) this is \(3\zeta(3)/(64\pi^3\nu)\). With the Euclidean contraction \(\eta\to\delta\) of the four matrix indices (the Poisson algebra kept) the ratio \(f_1/f_2\) at \(12\) sampled points and \(9\) directions of \(\zeta\) each lies between \(0.006\) and \(0.84\), never at \(1\); no claim is made over the whole leaf. For the Lorentzian contraction \(F\eta\) has one real and one imaginary eigenvalue pair, the electric-type case, to which the Euclidean heat-kernel construction does not apply.



\emph{Consequences.} (i) The statement that a fuzzy factor replicates the four-dimensional coefficient of Theorem~\ref{thm:Geff} over its modes, with threshold weights, is the limit \(sf_j\ll1\) of the diagonal above, \(\prod_j sf_j/\sinh sf_j=1-s^2(f_1^2+f_2^2)/6+\dots\), valid when \(\LUV^2\) exceeds the largest populated \(f_j\). In the Poisson limit \(|\zeta|\) is unbounded and the regime is the hypoelliptic one; the replication formula does not apply there, and the weighted mode count is replaced by \(\nu^{-1}\int\mathrm d^2\zeta\,(2\pi)^{-2}\prod_j sf_j/\sinh sf_j\). (ii) In the same nilpotent limit \(H_V\) acquires the Zeeman term \(\pm2f_1,\pm2f_2\), so its lowest eigenvalue at fixed \(\zeta\) is \(-|f_1-f_2|\) (checked by diagonalisation in a truncated Fock space): zero modes of unbounded number for \(f_1=f_2\), and eigenvalues unbounded below for \(f_1\ne f_2\), neither regulated by the proper-time cutoff. The vector sector therefore needs a fibre truncation (a finite fuzzy fibre) or fermionic partners, and in the Euclidean continuation used here \(e^{-sH_V}\) is unbounded for every \(s>0\) at the leading symbol, so the proper-time trace of this sector diverges. (iii) For \(sf_j\ll1\), the vector, ghost and \(n_{\mathrm{ex}}\) scalar sectors together give \(2+n_{\mathrm{ex}}+\tfrac{22-n_{\mathrm{ex}}}{6}s^2(f_1^2+f_2^2)\); the \(s^2\) term has the structure of the quartic term \(\{X^a,X^b\}^2\) and renormalises its coefficient logarithmically, not the Einstein--Hilbert coefficient. (iv) An Einstein--Hilbert coefficient would be the weight-two term \(\propto s^{-3}\) of the diagonal and its second variation under \(A\); no such computation is made here, and neither the sign nor the size of \(w\), hence of \(\varrho=\Lambda/w\), follows, with \(w\) and \(\varrho\) the induced and tree-level coefficients of the preceding paragraph.



This is a statement for the Poisson limit, the proper-time cutoff of \(\Box_X\) as the regulator, one loop, the leading (nilpotent) symbol, equiregular points, the Euclidean contraction \(\eta\to\delta\) with the Lie--Poisson algebra kept, background gauge and the bosonic content listed. It does not treat Moyal terms, fermions, the subleading symbol, the dependence on \(A\), or the identification of \(G\) with the emergent metric of Chapter~\ref{ch:vierbein-emergence}, and it does not bear on Assumption~\ref{ass:EE}. This changes no ledger label; Problem~F remains open.





\paragraph{Conjugate-coordinate class: the contraction by the ambient metric and the five-matrix presentation (conditional; bounded no-go and conditional route).} \label{detail:sec-F-paragraphconjugatecoordinatecl-34}

Keep the Poisson limit, the leaf \(\mathbb{CP}^{1,2}\) of \(\mathfrak{so}(4,2)\) with \(\eta_6=\mathrm{diag}(-1,1,1,1,1,-1)\), the notation \(y^a=M^{a5}\), \(y^4=M^{45}\), \(Q^{ab}=(y^4)^2\eta^{ab}+y^ay^b\) of the preceding paragraph, and \(T^a=M^{a4}\), \(D^a=\{T^a,\cdot\}\). On the orbit \((\eta_6M)^2=-\ell^2\mathbf 1\) (checked at \(200\) random points to \(10^{-12}\)), whence \(y\cdot y\equiv\eta_{AB}y^Ay^B=-\ell^2\) (\(A,B=0,\dots,4\)) and \(T^ay_a=0\). The base metric induced from the ambient space is \(g=\eta+y_\mu y_\nu/(y^4)^2\) with inverse \(g^{-1}=\eta+yy/\ell^2\); it is positive definite, and the contraction of the fields \(D^a\) by \(Q^{-1}\) has base symbol exactly \(g^{-1}\) (checked numerically). The kinetic tensor of the four matrices \(T^a\) is \(K=(y^4)^2\eta=(y^4)^2\bigl(g^{-1}-yy/\ell^2\bigr)\). The contraction by \(Q^{-1}\) equals \((y^4)^{-2}\bigl[\eta_{ab}D^aD^b+(y_aD^a)^2/\ell^2\bigr]\) and therefore replaces the Lorentzian tensor \(K\) by the Euclidean metric \(g\) of the hyperboloid; it is not a Wick rotation of \(\eta\).



\emph{No supply in the four-matrix class.} Hypothesis class: the Poisson limit; the four matrices \(T^a=M^{a4}\); the quartic term plus one bilinear term, with Hessian \(L(A)_a=-\Box A_a+D_a(D_bA^b)-2s\mathcal J_{ab}A^b+\mu^2A_a\); the background gauge function \(D_aA^a\), the ghost operator \(-\Box\), and a regulator that is a function of \(\Box_X\); the matrix-index form constant, \(\eta_{ab}\). In this class every operator is a polynomial with constant coefficients in the \(D^a\) and \(J^{ab}=\{M^{ab},\cdot\}\), that is, an element of the enveloping algebra of \(\mathfrak{so}(4,1)'\), the Lie algebra generated by the \(D^a\) and \(J^{ab}\) (the subalgebra of \(\mathfrak{so}(4,2)\) on the indices \(0,\dots,4\); the generators \(M^{a5}\) and \(M^{45}\) are not in it). The operator \((y_aD^a)^2\) is not an element of that enveloping algebra: its coefficients are the non-constant functions \(y_a\), which would have to be supplied by the missing generators. Hence in this class none of the action, the bilinear term, the background gauge, the ghosts and the \(\Box_X\) regulator produces the term \((y_aD^a)^2\). The full Casimir \(\tfrac12\eta_{AC}\eta_{BD}J^{AB}J^{CD}\) of \(\mathfrak{so}(4,2)\) has base symbol \(-2\ell^2g^{-1}\) (checked numerically), but it lies outside the enveloping algebra of \(\mathfrak{so}(4,1)'\) and enters neither the Hessian nor the ghost operator. Extensions of the matrix set by \(y^a\) with a positive-norm contraction (eight matrices \(T^a,y^a\); seven matrices \(T^a,y^i\)) have rank-six horizontal span, but the equation of motion then needs different bilinear coefficients for different matrices (checked by exact Poisson algebra). The eight-matrix set \(T^a,y^a\) with the \(y\)-block contracted by \(-\eta\) is stationary for one coefficient \(\mu^2=-4\), with horizontal rank six, and is not examined further.



\emph{Five-matrix presentation.} Take \(X^A=M^{Ac_0}\), \(A\ne c_0\), with the restriction of \(\eta_6\) as contraction. Definition~\ref{def:matrix-coords-ext} (\(D=5\)) covers \(c_0\in\{0,5\}\) up to relabelling; the forms for \(c_0\in\{1,\dots,4\}\), of signature \((3,2)\), lie outside it and are included as a comparison. The orbit relation gives \(y\cdot y=\eta_{c_0c_0}\ell^2\) for \(y^A=M^{Ac_0}\). At \(25\) random points for each \(c_0\) the horizontal rank is \(4\) and the centre rank \(2\); the induced contraction \(h=c^{T}\eta c\) on the horizontal space is positive definite for \(c_0\in\{0,5\}\) and has signature \((2,2)\) for \(c_0\in\{1,\dots,4\}\); a Lorentzian signature does not occur. For \(X^A=M^{A5}\), \(A=0,\dots,4\) (stabiliser \(\mathfrak{so}(4,1)\)): \(y_AD^A=0\) on the orbit, the base symbol of \(\eta_{AB}D^AD^B\) is \(\ell^2g^{-1}\) (checked to \(10^{-12}\)), and at \(300\) random points and \(91\) directions of \(\zeta\), for both signs of the bracket, the vertical flux in an \(h\)-orthonormal frame is self-dual or anti-self-dual to \(10^{-11}\), so \(f_1=f_2\) identically. This constant contraction of the five matrices is consistent with the statement that self-duality holds identically only for the contraction by the position-dependent \(Q^{-1}\) of the four matrices \(T^a\): its base symbol \(\ell^2g^{-1}\) is that of \(Q^{-1}\) up to the factor \((y^4)^2\). The bilinear coefficient required by stationarity is \(\mu^2=+4\) (exact, both signs of \(s\)), the Zeeman term of the gauge-fixed vector operator has the same modulus as in the four-matrix class, and in the polynomial Poisson algebra, exactly,

\[

H_V=C_{\mathfrak{so}(4,2)}-C_{G}+C_{\Sigma}+\mu^2,\qquad \Box=C_{\mathfrak{so}(4,1)}-C_{\mathfrak{so}(4,2)},

\]

where \(C_G\) is the Casimir of the total \(\mathfrak{so}(4,1)\) action on vector-valued functions and \(C_\Sigma\) that of the vector index. By the Weyl formulae the \(\mathfrak{so}(4,1)\) component \((k+1,k)\) of the degree-\(k\) vector polynomials has \(C_{\mathfrak{so}(4,2)}-C_G+C_\Sigma=0\) for every \(k\), so it is an eigenspace of \(H_V\) with eigenvalue \(\mu^2\) (verified directly for \(k=1\), dimension \(35\)). On the \(75\)-dimensional space of degree-one vector polynomials the formal eigenvalues of \(H_V\) are \(-8,-4,-2,+2,+4\) with multiplicities \(1,5,20,14,35\) (checked); polynomials are not normalisable, so the negative values have no spectral meaning. At the leading symbol the vector operator then has the lowest Landau level at zero for the four horizontal components, a zero-mode family with density \(f^2/(2\pi^2)\) in \(\zeta\), independent of the proper time; the normal component has no Zeeman term, and the bosonic density per \(\mathrm d^2\zeta/(2\pi)^2\) is \(f^2(16\pi^2)^{-1}\bigl[8+(3+n_{\mathrm{ex}})/\sinh^2(sf)\bigr]\). The \(\zeta\)-integral of the second term converges; that of the first diverges and is not regulated by the proper-time cutoff. The fibre of the leaf is compact, so \(\zeta\) is expected to be bounded by the fibre harmonic number of the quantum space; the divergence is a feature of the Poisson limit.



\emph{Compact check at finite size.} For the compact form \(\mathfrak{so}(6)\) with \(\delta\) throughout, on the \(N\)-th symmetric power of the fundamental representation, the exact Hessian of the quartic term plus \(\tfrac12m^2\operatorname{Tr}X^2\) (finite matrices, \(N=1,\dots,4\)) satisfies \(H+G^\dagger G=C_6-C_5^{\mathrm{tot}}+C_5^{\mathrm{vec}}+m^2\) for \(X^A=M^{A6}\), \(m^2=-4\). Its lowest eigenvalue is \(-4=m^2\) for \(N\le4\), with the family \((k+1,k)\), \(1\le k\le N\), of total multiplicity \(\sum_{k=1}^{N}(2k+1)(2k+3)(2k+5)/3\) and \(14\) further modes at \(-2\), so that the number of negative modes is \(49,154,385,814\) for \(N=1,\dots,4\) and this compact background is unstable; for the four-matrix background \(X^a=M^{a5}\), \(m^2=-3\), it is \(-(N+3)\), growing with \(N\). The sign of \(m^2\) is \(-4\) on the sphere and \(+4\) on the hyperboloid.



\emph{Consequences.} (i) In the matrix sets examined (\(T^a\); \(M^{Ac_0}\); the extended sets listed above), the term \((y_aD^a)^2\) is supplied by the matrix data only if the matrices are the position functions \(X^A=M^{A5}\); the price is that the horizontal metric induced by the contraction is then the Euclidean metric \(g\), and the conformally flat Lorentzian \(K\) of the four-matrix class is not obtained, so the level-one and stress results of the four-matrix class are not obtained in this presentation. (ii) For the hyperboloid, the vector zero modes of the leading symbol are lifted by \(+\mu^2\) only if the vanishing of the Casimir combination, which is exact on polynomial vector fields and on the compact form, continues to the normalisable sector; this continuation is not shown, and the negative formal eigenvalues \(-8,-4,-2,+2\) of the same operator on degree-one polynomials, like the negative modes of the compact background, may persist on normalisable states. (iii) Neither the sign nor the size of \(w\) follows: the weight-two coefficient and its variation under \(A\) are not computed, and the passage from a Euclidean coefficient to the Lorentzian one is a Wick-rotation hypothesis tied to Problem~E.



This is a statement for the Poisson limit, the leading symbol, equiregular points, the bosonic content and the single bilinear term; the self-duality of the flux is verified numerically and not proved, and the finite-size spectrum is for the compact form. It does not treat fermions, Moyal terms for the hyperboloid, or the Lorentzian emergent metric. It does not identify \(g\) with the emergent metric of Chapter~\ref{ch:vierbein-emergence}, does not derive Lorentzian signature, and leaves Assumption~\ref{ass:lorentzian} in place (Problem~E). This is a scope statement and not a retraction of any earlier result. This changes no ledger label; Problem~F remains open.





\paragraph{Conjugate-coordinate class: fibre admissibility and the Lorentzian \(i\epsilon\) window (conditional; bounded no-go).} \label{detail:sec-F-paragraphconjugatecoordinatecl-35}

Keep the setting and notation of the preceding paragraph (Poisson limit, leaf \(\mathbb{CP}^{1,2}\), \(T^a=M^{a4}\), \(\Box=\eta_{ab}D^aD^b\), \(H_V\) with \(\mu^2=-3\), \(C_{42}\), \(C_{41}\), \(C_{31}\), \(\sigma(\cdot)\), \(\Delta_\kappa\)); write \(S^2=2C_{41}-C_{42}\), whose symbol \(\beta q_f\) is purely vertical.



\emph{Fibre cutoffs.} In the Poisson limit the fibre of the leaf is a compact two-sphere and \(|\zeta|\) is unbounded; it is expected to be bounded by the fibre harmonic number of the quantum space. Two bounds on the harmonic number are in the literature. (a) \(S^2=2s(s+1)\) on the sectors \(\mathcal C^s\) with \(s\le n\), from the \((n+1)\)-dimensional representation of the stabiliser \(SU(2)_L\) on the lowest-weight state~\cite{Steinacker2019NoGhost}; the classical counterpart is the radius of the fibre as a coadjoint orbit of \(SU(2)_L\), the self-dual part of \(M^{ab}\) on the \(\eta\)-orthogonal complement of \(y\): at \(10\) orbit points exactly one of the two radii vanishes and the other equals \(\ell\), independent of \(y\). (b) A local bound \(s\le m=r^{-1}x^0\) from the uncertainty of the coherent states, which the same group states to correct earlier statements on the higher-spin cutoff~\cite{MantaSteinacker2025}; it grows with the point and is not the area bound of (a), and the two are not reconciled here. Neither is a sharp cutoff of commuting matrix-built operators.



\emph{Admissible cone.} For admissible modes, the principal series of \(SO(4,1)\), one has \(r^{-2}\Box_H>s^2+s+9/4\) with \(\Box_H=-C_{42}+C_{41}\) in the convention of~\cite{Steinacker2019NoGhost}, equivalent, by (2.7) and (3.43) of that reference, to \(-C_{42}>9/2\) in that convention (its footnote~8 prints \(C_2[\mathfrak{so}(4,2)]>9/2\) and states the identification with the \(SO(4,2)\) content of \(\operatorname{End}(\mathcal H_n)\) as plausible, not investigated). With \(\Box_H\) horizontal, \(S^2\) vertical and \(\sigma(C_{42})=-\alpha q_b+\beta q_f\) of the preceding paragraph, the semiclassical content of this bound is the cone \(-\sigma(C_{42})\ge0\), that is \(\alpha q_b\ge\beta q_f\); the Casimirs here and there differ by a sign at the symbol level. Conditional on this identification, a cutoff of \(C_{42}\) restricted to the cone together with a bound \(\beta q_f\le Q\) is a compact region, so a function of the central operator \(C_{42}\) regularises once the fibre is bounded. The cone restricts to the admissible sector and is not what bounds the window: the bound \(\beta q_f\le Q\) is a cutoff of \(S^2=2C_{41}-C_{42}\), so this relaxes the hypothesis of the preceding paragraph and does not contradict it, whose consequence~(i) for functions of \(C_{42}\) alone stands. The Lorentzian moments, the flux and the exact-sector results below are for the Gaussian window of \(\Delta_\kappa\), not for this fibre-bounded window. For fixed \(Q\) and \(\LUV^2\to\infty\) the window \(\{-\sigma(\Delta_\kappa)\le\LUV^2\}\) has its vertical term bounded by \(\beta(1-\kappa)Q\), so \(\kappa\) only rescales \(\LUV^2\) by \(2\kappa-1\) at leading order and the divergence at \(\kappa\to1\) is absent. The cutoff function and the base scale remain external: the quantisation map is \(SO(4,2)\)-equivariant and so acts on each irreducible component by a scalar, but this multiplier is not computed, and neither \(\chi\) nor the base scale is fixed. The window uses \(M^{AB}\), which are background data and not dynamical matrices, so the regulated trace is not shown to be a diffeomorphism-covariant functional of \(X^a\).



\emph{Lorentzian trace in a window.} Define the one-loop sum as \(\sum\ln(\lambda-i\epsilon)=-\int_0^\infty\frac{\mathrm ds}{s}\operatorname{Tr}\bigl[\chi\,e^{-is(H_V-i\epsilon)}\bigr]\) up to a scale-dependent constant, with \(\chi\) a function of the commuting operator \(\Delta_\kappa\); the contour is the proper time of \(H_V\) with the factor \(e^{-\epsilon s}\), and uses no metric. In a window of compact support (or a Gaussian) the number of modes per unit volume is finite, so the trace is finite for every \(s\), and \(\sum\ln(\lambda-i\epsilon)\) is a finite product over the window whether or not the Schwinger integral converges. For \(\chi=e^{-\sigma_W/\LUV^2}\), \(\sigma_W=-\sigma(\Delta_\kappa)\), the leading-symbol density is the Gaussian \((4\pi)^{-3}\LUV^6\det(\sigma_W)^{-1/2}\prod_{k=1}^{4}(1+is\LUV^2\lambda_k)^{-1/2}\), with \(\lambda_k\) the nonzero generalised eigenvalues of \((\sigma(\Box),\sigma_W)\); the kernel of \(\sigma(\Box)\) is not the vertical plane (checked at eight points). The mass-type moment is \(\langle(\sigma-i0)^{-1}\rangle=i\int_0^\infty\mathrm ds\prod_k(1+is\lambda_k)^{-1/2}\), whose real part is the principal value. For \(\lambda=(-c,a,a,a)\) the real part is \(2/(a+c)\) (closed form, checked), and for \(4000\) samples with \(c\in[0.1,10]\), \(a_k\in[10^{-3},10^{3}]\) it is positive. At eight orbit points (\(\kappa=\tfrac34\), \((y^4)^2\) from \(0.08\) to \(144\)) it is \(1.1\) to \(2.3\) times the value for the Euclidean contraction \(\eta\to\delta\), with the same sign, and the imaginary part is positive.



\emph{Flux.} At the same points the Lorentzian centre flux \(F_{ab}(\zeta)\) has \(\eta F\) with one real pair \(\pm a\) and one imaginary pair \(\pm ib\), \(a,b>0\) at all sampled \(\zeta\), and \(F_{ab}F^{ab}\) takes both signs (negative for \((y^4)^2\lesssim0.4\), positive above). The vector block therefore has Zeeman eigenvalues \(\pm2b\) (phases in \(e^{-is(H_V-i\epsilon)}\)) and \(\pm2ia\) (a factor \(e^{\pm2sa}\)). The Schwinger integrand of the vector sector then grows for the electric pair, bounded by \(e^{2sa}\) with \(a\) linear in the cutoff, while \(\ln\det(\sigma+Z)=\ln(\sigma^2-4b^2)+\ln(\sigma^2+4a^2)\) has integrable logarithmic singularities at \(\sigma=\pm2b\). The exponential growth of the Euclidean contraction is absent order by order in the flux (each order is a polynomial in \(s\) times \(e^{-is\sigma}\)) and for the magnetic pair, at the leading symbol and at the sampled points, and is present for the electric pair as a property of the integral representation and not of the determinant. The order-\(F^2\) term is proportional to \(F_{ab}F^{ab}\), so its sign is not definite on the leaf.



\emph{Invariant sector, exactly.} On fibre-constant functions \(\Box=-(\Delta_{H^4}-\Delta_{H^3})\) with \(H^4=\mathbb R_\tau\times H^3\), metric \(\mathrm d\tau^2+\cosh^2\tau\,\mathrm dH_3^2\), \(y^4=\ell\sinh\tau\), and the diagonal of \(e^{z\Delta_{H^4}+w\Delta_{H^3}}\) is the \(H^3\) integral of the product of the two heat kernels. Its expansion gives \(D_0=(4\pi)^{-2}z^{-1/2}(z+w\cosh^2\tau)^{-3/2}\), the Gaussian symbol formula above, and, for the pure \(\Box\) term (window removed), the Seeley--DeWitt coefficient \((5m^2-3\ell^2)/(4m^2)\), \(m=y^4\), whose product with \(\sqrt{|G|}=m^{-4}\) is the density \((5m^2-3\ell^2)/(4m^6)\) of the preceding paragraphs, obtained by an independent route. With the window the weight-two part of \(\sum\ln(\lambda-i\epsilon)\) is \(-(4\pi)^{-2}\LUV^2\,Q(a)\), \(a=\sinh\tau\), plus \(-2\) times a scale-dependent logarithm proportional to the number of modes, which is the \(R/6\) curvature term of the auxiliary \(H^4\) window geometry; \(\operatorname{Re}Q\) is indefinite, and with the subtraction at \(x=1\) it changes sign at \(a\simeq0.62\) (not at \(a^2=\tfrac35\), the zero of the coefficient above), a location that moves with the subtraction point and has no invariant meaning; \(\operatorname{Im}Q\to\pi\). This is the Lorentzian density of one scalar sector at the background \(A=0\), with a scale-dependent additive constant; it is not the Einstein--Hilbert coefficient.



\emph{Consequences.} (i) The coefficient carrying \(w\) is the weight-two term, a local Poisson-algebraic coefficient multiplied by window moments; the signature enters through the moments and the window, and in the potential channel the Lorentzian moment is finite and positive in the class sampled. (ii) The mode-count term is proportional to the difference between the numbers of bosonic and fermionic modes in the window, so its logarithmic scale dependence cancels only if that difference vanishes at this order. (iii) The weight-two coefficient as a quadratic form in \(A\) is not computed, the vector, ghost and fibre-harmonic sectors are not computed, and neither the sign nor the size of \(w\), hence no sign of an attractive coupling, follows.



This is a statement for the Poisson limit, one loop, the leading symbol and the invariant sector, the four-matrix class with \(\mu^2=-3\), Gaussian windows and the sampled points; the admissible cone is conditional on the identification stated, the Lorentzian moments are not computed in the fibre-bounded window, and the positivity of the moment is numerical. It does not treat fermions, Moyal terms, the dependence on \(A\), or a window built from dynamical matrices only; it identifies no induced metric with the emergent metric of Chapter~\ref{ch:vierbein-emergence} and, since a Euclidean coefficient is not a Lorentzian one, leaves Assumption~\ref{ass:lorentzian} in place (Problem~E). This changes no ledger label; Problem~F remains open.





\paragraph{Conjugate-coordinate class: the supertrace at an arbitrary matrix fluctuation, and the species constants that can multiply an Einstein--Hilbert-type term (conditional; bounded no-go).} \label{detail:sec-F-paragraphconjugatecoordinatecl-36}

Keep the setting of the preceding paragraph: Poisson limit, leaf \(\mathbb{CP}^{1,2}\), \(T^a=M^{a4}\), background gauge function, ghost, vector with \(\mu^2=-3\), \(n_{\mathrm{ex}}\) further scalars of mass \(\mu_i^2\), the supplied sixteen real fermion components of weight \(-\tfrac12\) with no bilinear fermion mass, and one window function \(\chi\) of the commuting operator \(\Delta_\kappa\) for all species. Let \(X^a\) be arbitrary functions on the leaf, \(D_a=\{X_a,\cdot\}\), \(F_{ab}=\{X_a,X_b\}\), \(\mathcal Z_{ab}=\{F_{ab},\cdot\}\) and \(\Box_X=\eta^{ab}D_aD_b\).



\emph{Operators at arbitrary \(X\).} The Jacobi identity gives \([D_a,D_b]=\mathcal Z_{ab}\). The second variation of the quartic term plus the bilinear term about \(X\) is \(\tfrac12\int\varepsilon^aL_a\) with \(L_a=-\Box_X\varepsilon_a+D_aD_b\varepsilon^b-2\{F_{ab},\varepsilon^b\}+\mu^2\varepsilon_a\); in the background gauge \(D_a\varepsilon^a\) the ghost operator is \(-\Box_X\), the further scalars at vanishing transverse matrices have \(-\Box_X+\mu_i^2\), and the square of the Dirac operator is \(\mathcal D^2=\Box_X+\tfrac12\Gamma^{ab}\mathcal Z_{ab}\) (these three identities were checked symbolically for random polynomial \(X\) on a four-dimensional phase space). Every species is therefore \(H_i=(-\Box_X)\otimes1+Z_i+m_i\) with \(Z_i=\Sigma_i^{ab}\otimes\mathcal Z_{ab}\), the same operator \(-\Box_X\) and the same \(\mathcal Z_{ab}\) for all species, constant masses \(m_i\) (\(\mu^2\), \(\mu_i^2\), and zero for the ghost and the fermions), and constant matrices \(\Sigma_i^{ab}\) acting on the species index (for the vector, minus the generator of the vector index, which gives \(-2\{F_{ab},\varepsilon^b\}\); for the fermions, minus \(\tfrac14[\Gamma^a,\Gamma^b]\) on four copies of the real four-component spinor, because the operator compared with \(-\Box_X\) is \(-\mathcal D^2\); the traces satisfy \(\operatorname{tr}\Sigma_D\Sigma_D=\tfrac12\operatorname{tr}\Sigma_V\Sigma_V\) for all index pairs).



\emph{Species sums.} With weights \(w_i\) and dimensions \(d_i\) (\(1\cdot4\), \(1\cdot n_{\mathrm{ex}}\), \(-2\cdot1\), \(-\tfrac12\cdot16\)) one has, as exact tensor identities checked for all index tuples, \(\sum_iw_id_i=0\) and \(\sum_iw_i\operatorname{tr}(Z_i^{a_1b_1}\cdots Z_i^{a_nb_n})=0\) for \(n=1,2,3\) (not only for the totally symmetrised products); for \(n=4\) the symmetrised sum is the \(s^4\) coefficient \((f_e^2+f_m^2)^2\) of the character of the preceding paragraph and does not vanish. Because \(\Sigma_i\) acts on the species index only, the term of any Duhamel expansion of \(\sum_iw_i\operatorname{Tr}[\chi\,\Phi(H_i)]\) in powers of \(m_i\) and \(Z_i\) is a species constant times a species-independent operator trace; no commutation of \(\chi\), \(\mathcal Z_{ab}\) and \(\Box_X\) is needed. Hence, for every \(X\) and every \(\Phi\), the terms without masses and with at most three insertions of \(\mathcal Z_{ab}\) cancel identically (a check with random non-commuting matrices confirms this), and the term linear in the masses and without insertions of \(\mathcal Z_{ab}\) is \(+M_{\mathrm{str}}\operatorname{Tr}[\chi\,\Phi'(H_X)]\) (for \(\Phi=e^{-(\cdot)}\) this is \(-M_{\mathrm{str}}\operatorname{Tr}[\chi\,e^{-H_X}]\)) with \(M_{\mathrm{str}}=\sum_iw_id_im_i=4\mu^2+n_{\mathrm{ex}}\mu_i^2\). The other constants that survive at the next order are \(\sum_iw_id_im_i^2\), \(\sum_iw_im_i\operatorname{tr}\Sigma_i^{ab}\Sigma_i^{cd}\) (the vector mass alone) and the quartic constant.



\emph{Consequences.} (i) The weight-two coefficient of the preceding paragraph is, for every \(X\) and not only at \(X=T\), the mass-supertrace term: no count term \(\propto R\), no spin term quadratic in \(\mathcal Z_{ab}\) and no term cubic in it appears, whatever the dependence of \(\Box_X\) and \(\mathcal Z_{ab}\) on \(A\). (ii) At the leading symbol that term is proportional to \(M_{\mathrm{str}}\int\chi\,(\sigma_X-\mathrm i0)^{-1}\) with \(\sigma_X=\eta_{ab}\,\xi(D^a)\xi(D^b)\), which depends on \(X\) through \(\mathrm dX^a\) only. For \(A^a=\alpha^a\cos\phi\), \(\phi=q\cdot y\), with constant \(\alpha^a\), the Hamiltonian field of \(X^a\) is \(W^a-\epsilon\sin\phi\,\alpha^aY_q\), where \(W^a\) is that of \(T^a\) and \(Y_q=q_\mu\{y^\mu,\cdot\}\) (chain rule; the first-order change of the field of \(F_{ab}\) was checked through the Poisson algebra for the linear ansatz \(A^a=\alpha^a\,q\cdot y\)). This leading-symbol term depends on \(A\) only through its first derivatives, so its second variation contains no second derivative of the fluctuation and no term of Einstein--Hilbert form; the next-to-leading symbol, which involves the second jet of \(X\), is not computed. With the Gaussian window of \(\Delta_{3/4}\) the Hessian of the real (principal-value) part of this moment along \(\alpha\), at frozen coefficients and \(\phi=\pi/2\), was evaluated numerically in \(59\) random draws of orbit point and \(q\): it is indefinite in \(40\), positive definite in \(5\) and negative definite in \(14\), and along \(\alpha^a\propto q^a\) it is numerically nonzero in all draws with non-null \(q\). (iii) A term of Einstein--Hilbert form can therefore enter only through the constants listed above. The one proportional to \(M_{\mathrm{str}}\) is the mass-insertion term of the weighted \(a_2\) in the Euclidean class (there a shift of the operator of species \(i\) by \(m_i\) changes the weighted \(a_2\) by \(-m_ia_1^{(i)}\), with curvature part \(-\tfrac16M_{\mathrm{str}}\,\rho R[g_\rho]\) and connection part \(-(\sum_iw_i\kappa_im_i)\rho^2X\), \(\kappa_i\) being the weight of \(\rho^2X\) of the species; the second constant above is the vector mass times \(\kappa_V\)); its value in the six-dimensional leaf requires the next-to-leading symbol and is not computed. The pure flux constant is the order-\(F^4\) term: at the phase \(\phi=0\), where the first-jet terms vanish, and at the leading order in \(q\), the flux in the leading symbol is shifted as \(F_{ab}(\xi)\to F^0_{ab}(\xi)+\epsilon\,c\,(\xi\cdot Y_q)(q_a\alpha_b-q_b\alpha_a)\), \(c=-s\,y^4\cos\phi\), and at fixed proper time and \(\kappa=\tfrac34\) the Gaussian averages in the integrand depend on \(q\cdot y\) at fixed \(q\cdot q\), through the vector \(y^\mu\) and the window (at an orbit point with \(y^i=0\), \(\langle(\xi\cdot Y_q)^2\rangle\) changes by a factor above \(2\) over a boost of rapidity \(1.2\) of a unit \(q\)); its integrand is therefore not Lorentz invariant (the window and the leaf both single out \(y^\mu\), and this diagnostic does not separate them). (iv) No finite coefficient of this term follows in the present window: at four orbit points the integrand of the order-\(F^4\) term at \(A=0\) grows as \(s^2\) at large proper time (log-slope \(2.00\)), because the flux does not vanish on the kernel plane of \(\sigma(\Box)\); the expansion in \(\mathcal Z_{ab}\) is not the expansion in a small parameter there, and a sign of an Einstein--Hilbert-type coefficient does not follow. The non-expanded leading-symbol logarithm is \(\sum_iw_i\ln\det(\sigma+Z_i)=\ln(\sigma^2-4f_m^2)+\ln(\sigma^2+4f_e^2)+4\ln\sigma-2\ln[(\sigma+f_m)^2+f_e^2]-2\ln[(\sigma-f_m)^2+f_e^2]\), whose large-\(\sigma\) series starts with \(-6(f_e^2+f_m^2)^2/\sigma^4\) (checked); its second variation under \(A\) is not computed.



This is a statement for the Poisson limit, one loop, the commuting-symbol (leading) level, the four-matrix class with \(\mu^2=-3\), the supplied fermion content with no bilinear fermion mass and the same window for all species, applied term by term in an expansion that is not shown to converge. It does not compute the next-to-leading symbol, the dependence on \(A\) of the terms of Einstein--Hilbert form, the case of transverse matrices with flux or of species-dependent windows, Moyal terms, or a window built from dynamical matrices; it identifies no induced metric with the emergent metric of Chapter~\ref{ch:vierbein-emergence} and gives no value or sign of \(G_{\mathrm{eff}}\). This changes no ledger label; Problem~F remains open.





\paragraph{Conjugate-coordinate class: exact Landau (nilpotent) resummation of the leading symbol in the background flux, and its effect on the mass-supertrace term and on the massless resummed character term (conditional; bounded no-go).} \label{detail:sec-F-paragraphconjugatecoordinatecl-37}

Keep the setting of the preceding paragraphs: Poisson limit, leaf \(\mathbb{CP}^{1,2}\), \(T^a=M^{a4}\), \(H_X=-\Box_X\), the supplied content with no bilinear fermion mass, \(M_{\mathrm{str}}=4\mu^2+n_{\mathrm{ex}}\mu_i^2\), and the nilpotent (leading) symbol, in which the horizontal momenta \(\omega_a=\xi(W_a)\) obey \([\omega_a,\omega_b]=\mathrm iF_{ab}(k)\) with \(F_{ab}\) the constant flux at the centre variable \(k\) of the kernel plane of \(\sigma(\Box)\); \(\eta F\) has eigenvalues \(\pm f_e\) and \(\pm\mathrm if_m\). At four sampled orbit points the ratio \(f_e/f_m\) on the kernel plane equals \(\ell/|y^4|\) in every direction of \(k\) to five digits, which is the closed form of the kernel-plane flux above. Write \(a=f_e\), \(b=f_m\), \(x=sa\), \(y=sb\).



\emph{Orbital factor.} For \(H_0=\eta^{ab}\omega_a\omega_b\) the Gaussian (KMS) two-point function \(G=-\mathrm iFM^{T}(1-M^{T})^{-1}\), \(M=\mathrm e^{2sF\eta}\), gives, per unit horizontal volume, \(\Theta(s)=\operatorname{Tr}\mathrm e^{-\mathrm isH_0}=-\mathrm i\,ab\,\bigl(16\pi^2\sinh x\,\sin y\bigr)^{-1}\), with \(s\to s-\mathrm i0\) (the poles of \(1/\sin y\) lie above the contour, as for the Landau-level sum \(\tfrac{b}{2\pi}\sum_n\mathrm e^{-\mathrm isb(2n+1)}\)). The identity \(\mathrm d\ln\Theta/\mathrm ds=-\mathrm i\eta^{ab}G_{ab}\) holds to \(10^{-10}\) at complex \(s\), and \(\Theta\to-\mathrm i/(16\pi^2s^2)\) for \(s\to0\). The electric pair gives \(1/\sinh\) and the magnetic pair \(1/\sin\).



\emph{Mass-supertrace term.} The term \(+M_{\mathrm{str}}\operatorname{Tr}[(H_X-\mathrm i0)^{-1}]=M_{\mathrm{str}}\,\mathrm i\!\int_0^\infty\!\mathrm ds\,\operatorname{Tr}\mathrm e^{-\mathrm isH_X}\) becomes, per unit horizontal volume and per centre point, \(M_{\mathrm{str}}(16\pi^2)^{-1}(s_0^{-1}+g)\) with

\[

g(a,b)=\int_0^\infty\!\mathrm ds\Bigl[\frac{ab}{\sinh(as)\,\sin(bs-\mathrm i0)}-\frac1{s^2}\Bigr]=\pi\bigl[\,b\,S(b/a)+\mathrm i\,a\,S(a/b)\bigr],\qquad S(r)=\sum_{n\ge1}\frac{(-1)^n}{\sinh(n\pi r)}<0 .

\]

The integrand is analytic in the open fourth quadrant and bounded by \(\exp[-|s|(a\cos\vartheta+b\sin\vartheta)]\) on the ray \(s=|s|\mathrm e^{-\mathrm i\vartheta}\), so \(g\) is finite for \(a,b>0\) and independent of \(\vartheta\); the closed form follows from the two residue sums (poles of \(1/\sinh\) at \(s=-\mathrm in\pi/a\), poles of \(1/\sin\) at \(s=n\pi/b\)) and was confirmed against the ray integral to \(10^{-24}\) at ten pairs \((a,b)\). \(g\) is homogeneous of degree one, \(\operatorname{Re}g<0\) and \(\operatorname{Im}g<0\) for all \(a,b>0\), \(\operatorname{Re}g\to-a\ln2\) for \(b\ll a\) and \(\operatorname{Im}g\to-b\ln2\) for \(a\ll b\). Thus the linear divergence in \(s\) of the Maxwell tail \(c_0\) of the preceding paragraph is cut off at \(s\sim1/f\), and the flux supplies an exponentially convergent infrared end of the proper-time integral.



\emph{Completeness of the \(M_{\mathrm{str}}\) term.} With the flux kept to all orders, the mass-linear term is not \(M_{\mathrm{str}}\operatorname{Tr}[(H_X-\mathrm i0)^{-1}]\) alone: the vector mass multiplies the Zeeman-dressed resolvent, and the complete term is \(\mathrm i\!\int\!\mathrm ds\,\Theta(s)\bigl[M_{\mathrm{str}}+\mu^2(\chi_V-4)\bigr]\), \(\chi_V=2\cosh2x+2\cos2y\). The second part is the all-order version of the constant \(\sum_iw_im_i\operatorname{tr}\Sigma_i\Sigma_i\) listed above and has the growth of the next paragraph.



\emph{Massless resummed character term.} With \(\chi_{\mathrm{str}}=4(\cos y-\cosh x)^2\) (the species sum \(2\cosh2x+2\cos2y+6-2-8\cosh x\cos y\), identity checked to \(10^{-40}\)) the resummed integrand is \(\Theta\chi_{\mathrm{str}}/s\). On the ray of angle \(\vartheta\) one has \(\bigl|\Theta\chi_{\mathrm{str}}\bigr|\sim\exp\bigl[|s|\,|a\cos\vartheta-b\sin\vartheta|\bigr]\) (measured rates agree to \(4\cdot10^{-3}\) at four angles away from \(\vartheta_*\)), and \(\Theta\chi_{\mathrm{str}}\to(ab/\pi^2)\sinh^2[s(a-\mathrm ib)/2]\) at large \(|s|\) (ratio \(1\) to \(10^{-8}\)). The exponential growth therefore occurs on the real axis (the \(\mathrm i\epsilon\) contour) for every \(a>0\), the same exponent holds for \(\Theta(\chi_V-4)\), and the massless exponent vanishes only on the single ray \(\vartheta_*=\arctan(a/b)\), which depends on the flux; on the leaf \(\vartheta_*=\arctan(\ell/|y^4|)\). The exponent is unchanged by the phases \(\mathrm e^{-\mathrm ism_i}\) of real masses on the real axis; on a ray of angle \(\vartheta\) a mass \(m\) shifts it by \(-m\sin\vartheta\), so for the vector with \(\mu^2=-3\) the rate is \(|a\cos\vartheta-b\sin\vartheta|+3\sin\vartheta>0\) on every ray, whereas a positive mass-squared would open a window of convergent rays about \(\vartheta_*\) for the massless term treated here. On that ray the integrand is bounded, but it tends to the constant \(-ab/(2\pi^2)\) plus oscillations, so \(\int\mathrm ds/s\) grows logarithmically in the real part (numerically \(-0.032\) per doubling of the upper limit at \(a=0.7\), \(b=1.3\); the imaginary part converges). For the massless term no ray of the closed fourth quadrant gives a convergent integral (rays of the first quadrant correspond to the opposite \(\mathrm i\epsilon\) prescription and are neither claimed nor excluded), and with \(\mu^2=-3\) the same holds for the mass-dressed term, so the massless resummed character term has no finite value from the proper-time representation of the class; its non-analytic dependence on \(f_e\) is that of a spectral sum over the complex eigenvalues \(\lambda\pm2\mathrm if_e\) of the vector sector and is not computed.



\emph{Dependence on \(A\).} The leading nilpotent coefficient of the \(M_{\mathrm{str}}\) term depends on \(A\) only through the centre-valued flux, so its second variation is \(\tfrac12D^2g\,[Q,Q]\) with \(Q_{ab}=\zeta_Y(k)(\kappa_a\alpha_b-\kappa_b\alpha_a)\); it contains no second derivative of the fluctuation and no term of Einstein--Hilbert form; this concerns the \(M_{\mathrm{str}}\) term only, since the Zeeman-dressed part has no contour. It rests on the frozen nilpotent frame in which \(Y\) is central, and on the leaf \([W_a,Y]\ne0\) at order \(q\) by the Jacobi identity; the same omission is made in the preceding frozen-frame paragraph. At four orbit points, averaging over the direction of \(k\) with the window of \(\Delta_{3/4}\), this form is gauge invariant (\(\alpha\propto\kappa\) in its kernel) but is not proportional to \((\kappa\cdot\kappa)\eta^{ab}-\kappa^a\kappa^b\) (relative residual \(0.4\)--\(1.0\) in real and imaginary parts at the points with \(0.3\le f_e/f_m\le0.8\)), and its eigenvalues on the quotient by \(\kappa\) have both signs. The infrared-finite replacement of the Maxwell tail of the \(M_{\mathrm{str}}\) term is therefore not of the covariant Maxwell form and is not sign-definite (leaf average, four points).



\emph{Einstein--Hilbert sector in a four-dimensional Landau model.} Let \(\Omega_a=\omega_a+\epsilon\sin\phi\,(m_a{}^c\omega_c+n_a)\), \(\phi=\kappa\cdot x\), \(m_a{}^c\kappa_c=0\), on the algebra above with \(F\) constant; the choice \(m_a{}^c=-\alpha_ay^c\), \(n_a=-\alpha_a c\) reproduces the sector of the preceding paragraph with the commutators \([\omega_a,\omega_b]=\mathrm iF_{ab}\) kept, but without the commutator of \(Y\) with \(W_a\) that the Jacobi identity fixes on the leaf, so it is a model and not the leaf (an omission shared with the frozen-frame paragraph above). The second order in \(\epsilon\) of \(\operatorname{Tr}\mathrm e^{-\mathrm isH}\) was computed exactly by Wick's theorem in the state \(\mathrm e^{-\mathrm isH_0}\). It reproduces the Maxwell coefficient at \(F=0\) (ratio \(1\) to \(10^{-10}\)), the form \(A\) of the preceding paragraph at \(F=0\) (\(c_1/A=\mathrm i\) to \(10^{-8}\) in five random draws, including the dilaton term), gauge invariance in \(n\) at \(F\ne0\) (\(2\cdot10^{-15}\)), and the constant-field limits (Hessian of \(\Theta\) in \(F\), and \(|\det M|^{-1}\Theta(MFM^T)\) for a constant vielbein). At \(F\ne0\) the coefficient \(c=A/16\pi^2\) of \(\mathrm ds/s\) is unchanged, the integrand decays like \(\mathrm e^{-|s|(a\cos\vartheta+b\sin\vartheta)}\), and the finite part \(\mathrm{Fin}=\lim[\int_{s_0}^\infty T_2\,\mathrm ds-c\ln(1/s_0)]\), in units of \(c\), is complex; its real part is about \(7.6\), \(8.1\), \(6.8\) in three Lorentz frames of \(F\) at \((a,b)=(0.7,1.3)\) (data \(\kappa,\alpha,y\) rotated, \(A\) fixed; imaginary parts between \(-0.32\) and \(0.20\)), and \(7.6\), \(5.9\), \(7.3\) in three frames at \((a,b)=(1.92,1.0)\) (imaginary parts \(1.05\)--\(1.28\)). The flux therefore turns the infrared end of \(\int\mathrm ds/s\) into a finite remainder that depends on the orientation of \(F\) relative to \(\kappa,\alpha,y\): at fixed \((a,b)\) and fixed \(A\) (the meaningful comparison, since \(\mathrm{Fin}/c\) shifts by a constant under a rescaling of \(s\)) the real part varies between frames by roughly \(20\%\), and the Maxwell part of the same sector changes between frames by about \(30\)--\(45\%\) at \((a,b)=(0.7,1.3)\) and by a factor of order one at \((1.92,1.0)\) (complex values). The coefficient of the logarithm is the flux-independent \(A\) of the preceding paragraph.



This is a statement for the Poisson limit, one loop, the nilpotent (leading) symbol resummed in a constant flux at the centre variable, the Lorentzian \(\mathrm i\epsilon\) proper time, the supplied content, and numerical evaluation at sampled points; the closed form of \(g\) is checked numerically, not proved, and the Landau model for the Einstein--Hilbert sector is not the leaf. It does not include the next-to-leading symbol of the hypoelliptic operator with the commutators of \(Y\), the dependence of the window on the flux, non-perturbative masses, or the spectral definition of the vector sector, and it identifies no induced form with the emergent metric of Chapter~\ref{ch:vierbein-emergence}, makes no Wick-rotation claim, does not consider the first-quadrant rays, and gives no value or sign of \(G_{\mathrm{eff}}\). This changes no ledger label; Problem~F remains open.







\section{Problem G: HBEV}

\label{sec:detailed-Problem G: HBEV}



\paragraph{Ledger label.} \label{detail:sec-G-paragraphledgerlabel-38}

Status: OPEN. Priority: CRITICAL PATH. Completion:

PROGRAMME-CONDITIONAL. Reduced-scope results: G-model and

G-external support at most a scope-limited model; G remains OPEN

and neither satisfies the matrix-derived criterion.



The boundary functional of Chapter~12 regulates \(\Lambda_{\mathrm{eff}}\)

once an infrared entropy density is supplied. The working value

\(\Lambda_{\mathrm{eff}}=3/L_{\mathrm{IR}}^2\) is a choice consistent with

that functional, not a derived necessity.







\section{Problem H: modular origin of the interpolating function}

\label{sec:detailed-Problem H: modular origin of the interpolating function}



\paragraph{Ledger label.} \label{detail:sec-H-paragraphledgerlabel-39}

Status: OPEN. Priority: OPTIONAL within the phenomenological scope.

Completion: PROGRAMME-CONDITIONAL. The modular derivation of the

interpolating function is not established.



Appendix~\ref{app:empirical} tests the function

\[

\nu_{\sqrt{}}(g_{\mathrm N};g_s)

=

\frac{1}{1-\exp\!\left(-\sqrt{g_{\mathrm N}/g_s}\right)}

\]

against SPARC with \(g_s\) fixed to the empirical value

\(g_\dagger\). This is a coefficient-level consistency check, not

a microscopic derivation. Appendix~\ref{app:matrix_entropy} gives

the standard heat-kernel area contribution but does not compute a

source-dependent modular spectrum. Appendix~\ref{app:gaia} applies

an external-field prescription to wide binaries; it does not derive

that prescription from the modular data.





\paragraph{Missing step.} \label{detail:sec-H-paragraphmissingstep-40}

Construct the regulated modular operator, or an equivalent

well-defined spectral response, for a weak-field causal diamond

containing a baryonic source. Derive the relevant excitation gap

\(\Delta k(g_{\mathrm N},\kappa_{\mathrm{IR}})\) and its spectral

weight from the microscopic theory. Fix the normalization relating

\(g_s\) to \(\kappa_{\mathrm{IR}}\) independently of galaxy data.

In particular, reconcile the convention

\(g_s\sim\kappa_{\mathrm{IR}}/(2\pi)\) in

Appendix~\ref{app:rar-bridge} with the undetermined factor

\(g_s=c_\star\kappa_{\mathrm{IR}}\) in Chapter~17 before claiming

that the Bose occupation yields the target.



Until this calculation is supplied, Appendices~\ref{app:empirical},

\ref{app:matrix_entropy}, and~\ref{app:gaia} provide an empirical fit

check, a standard area-term calculation, and an external-field

consistency test, respectively; together they do not derive or

uniquely select \(\nu_{\sqrt{}}\). Problem~H remains open.





\paragraph{Ledger label.} \label{detail:sec-H-paragraphledgerlabel-41}

Status: OPEN. Priority: OPTIONAL falsifiable phenomenological

extension; prerequisite only if drag is claimed as derived or predicted.

Completion: PROGRAMME-CONDITIONAL.



Chapter~\ref{ch:dissipation} adopts a vector drag proportional to

\(-\mathbf v/d^4\) as a working assumption. The microscopic target

is to derive the low-frequency force-response tensor, including its

separation and finite-size dependence, from specified dynamics and

a specified stationary reference state.





\paragraph{Static tail, distinct from the response.} \label{detail:sec-H-paragraphstatictaildistinctfro-42}

Chapters~\ref{ch:entanglement-equilibrium} and \ref{ch:dissipation}

and Appendix~\ref{app:rar-bridge} use the same large-separation

scaling,

\[

I(A:B)\sim c_{p_*}

\left(\frac{\ell_A\ell_B}{d^2}\right)^{p_*},

\qquad d\gg\ell_A,\ell_B,

\]

which is dimensionless. It is established in the literature for

R\'enyi mutual information of regions in a conformal vacuum, with the

von Neumann limit shown for free-scalar spheres

\cite{Cardy2013MI}; its leading power for NAMCG cores is an

assumption. Neither this static tail nor a coincidence of its power

with the assumed drag exponent determines the retarded force

response, its tensor structure, or its coefficient.





\paragraph{Ledger label.} \label{detail:sec-H-paragraphledgerlabel-43}

Status: OPEN. Priority: OPTIONAL within the stated ECG model.

Completion: PROGRAMME-CONDITIONAL. The reduced results do not derive

singularity resolution from the matrix theory.



For the explicit invariant \(\mathcal P\) in~\eqref{eq:ECG}, consider

the reduced action per unit time and angular volume

\[

L[N,f]=N r^2\bigl(R-2\Lambda_{\mathrm{eff}}+\kappa\mathcal P\bigr),

\qquad

\kappa=16\pi G_{\mathrm{eff}}\beta_{\mathrm{ECG}}.

\]

This reduction assumes that the static, spherically symmetric

variations capture the corresponding field equations; an appropriate

symmetric-criticality argument is not established here. The reduced

Euler--Lagrange operators, including all higher-derivative terms, are

\[

\mathcal E_q

=\frac{\partial L}{\partial q}

-\frac{\mathrm d}{\mathrm dr}\frac{\partial L}{\partial q'}

+\frac{\mathrm d^2}{\mathrm dr^2}\frac{\partial L}{\partial q''}

-\cdots,\qquad q=N,f.

\]

Variations are taken with the fields and the derivatives required by

the order of \(L\) fixed at the endpoints (or with compact support);

otherwise the associated boundary terms must be retained. For the

specified \(\mathcal P\), these reduced equations have been checked

symbolically:

\[

\left.\mathcal E_f\right|_{N=1}\equiv0,\qquad

\left.\mathcal E_N\right|_{N=1}

=\frac{\mathrm d}{\mathrm dr}

\left[2r(1-f)+3\kappa\mathcal B[f]

-\frac23\Lambda_{\mathrm{eff}}r^3\right].

\]

Thus \(N=1\) is a consistent vacuum branch of the reduced equations,

and, with \(\alpha=-\kappa\), the lapse equation has first integral

\[

\mathcal F[f]-\frac23\Lambda_{\mathrm{eff}}r^3=4M,\qquad

\mathcal F[f]=2r(1-f)-3\alpha\,\mathcal B[f],

\]

where \(\alpha\) is the reduced ECG coupling and

\[

\mathcal B[f]

=

\frac43(f')^3

-\frac8r\,ff''

 \left(\frac{rf'}2+1-f\right)

+\frac4r(f')^2

-\frac8{r^2}ff'(f-1).

\]

For the rational profile with \(a=r_0^3\) and

\(\Lambda_{\mathrm{eff}}=0\),

\[

f_{\mathrm{rat}}(r)

=

1-\frac{2Mr^2}{r^3+a}

=

1-\frac{2M}{r}+\frac{2Ma}{r^4}+O(r^{-7}),

\]

substitution gives

\[

\mathcal F[f_{\mathrm{rat}}]-4M

=-\frac{4Ma}{r^3}-\frac{432\alpha M^2}{r^5}+O(r^{-6}).

\]

The \(r^{-3}\) term cannot be cancelled by the cubic contribution,

which begins at \(r^{-5}\). Hence, for the displayed equation, this

profile is not an exact vacuum solution when \(Mr_0\ne0\).



Moreover, a smooth regular center has \(f(r)=1+O(r^2)\), which

implies \(\mathcal B[f]\to0\) and

\(\mathcal F(0)-\tfrac23\Lambda_{\mathrm{eff}}0^3=0\).

Consequently the integration constant vanishes: this displayed

equation has no smooth regular-centre vacuum solution with a nonzero

integration constant, for any \(\Lambda_{\mathrm{eff}}\). For

\(\Lambda_{\mathrm{eff}}=0\), the asymptotic falloff

\(f=1-2M/r+o(r^{-1})\) identifies that constant as \(4M\) and \(M\)

as the ADM mass. For nonzero \(\Lambda_{\mathrm{eff}}\), its

interpretation is not automatically the ADM mass.





\paragraph{Large-radius tail.} \label{detail:sec-H-paragraphlargeradiustail-44}

For the displayed equation, write \(f=1-2M/r+\delta(r)\). Then

\(2r(1-f)=4M-2r\delta\) and

\(\mathcal B[1-2M/r]=144M^2/r^5+O(r^{-6})\). The leading

correction and the next asymptotic order give

\[

f(r)=1-\frac{2M}{r}-\frac{216\,\alpha M^2}{r^6}

+\frac{368\,\alpha M^3}{r^7}+\cdots .

\]

The rational profile instead has a \(+2Ma/r^4\) term. The tail

expansion is an asymptotic check of the displayed first integral; it

does not determine the model parameter \(r_0\) or establish a match to

an independently published expansion. In the stated action

normalization \(\alpha=-16\pi G_{\mathrm{eff}}\beta_{\mathrm{ECG}}\).

If \(G_{\mathrm{eff}}>0\), then \(\alpha>0\) iff

\(\beta_{\mathrm{ECG}}<0\), which fixes the sign of the leading

correction. For \(M>0\), its magnitude defines the characteristic

scale \(r_*\sim|\alpha M|^{1/5}\); it is not a core scale.







\section{Problem N: control of the effective expansion}

\label{sec:detailed-Problem N: control of the effective expansion}



\paragraph{Ledger label.} \label{detail:sec-N-paragraphledgerlabel-45}

Status: OPEN. Priority: CRITICAL PATH for the matrix-derived label.

Completion: PROGRAMME-CONDITIONAL.



Assumption~\ref{ass:EFT-a2} treats the four-derivative terms as

expansion vertices at order \((p/\LUV)^2\) and not as propagating

modes. It is stated in Chapter~\ref{ch:ultraviolet} and justified

there only by the composite character of the metric. The ghost mass

of Stelle gravity is of order \(1/\sqrt{\alpha G}\), and its relation

to \(\LUV\) depends on the physical \(G_{\mathrm{eff}}\) and on

\(\alpha\), which are part of Problem~F.





\paragraph{Induced Weyl term and the extra spin-two pole (conditional; not closure).} \label{detail:sec-N-paragraphinducedweyltermandthe-46}

Work in the truncated identity of Theorem~\ref{thm:Geff}, expanded about

a flat background with the vacuum term cancelled or neglected so that

flat space is a solution, with every included sector assigned the common

limits \(\LUV,\mu\), \(\mu\ll\LUV\), with \(S_E=-iS_L\) as in

Appendix~\ref{app:seeley-dewitt}, and with curvature invariants

continued without change of sign (an added convention). Rewrite the

traced \(a_{2,s}\) of~\eqref{eq:appA-a2} in the basis

\(\{C^2,R^2,E_{\mathrm{GB}}\}\), and let

\(\mathcal W_C:=\sum_sw_sc_{C,s}\) and \(\mathcal W_R:=\sum_sw_sc_{R,s}\),

where \(c_{C,s}\) and \(c_{R,s}\) are the coefficients of

\(C_{\mu\nu\rho\sigma}C^{\mu\nu\rho\sigma}\) and \(R^2\). For a real

scalar sector with \(E=E_0+QR\), \(c_{C,s}=\operatorname{tr}\mathbb I_s/120\)

and \(c_{R,s}=\tfrac12\operatorname{tr}(\tfrac16\mathbb I_s-Q)^2\); for

other sectors \(c_{C,s}\) also depends on \(E\) and on the bundle

curvature and is not computed here. The logarithmic term

of~\eqref{eq:W-expand} gives, in~\eqref{eq:W4},

\[

  c_W=\frac{\mathcal W_C L}{32\pi^2},\qquad

  c_R=\frac{\mathcal W_R L}{32\pi^2},\qquad

  L:=\ln\frac{\LUV^2}{\mu^2},

\]

that is, \(\alpha=-2c_W\) and \(\beta=2c_R\) in~\eqref{eq:stelle}; the

Gauss--Bonnet term does not contribute at quadratic order on a flat

background. With \(g_{\mu\nu}=\eta_{\mu\nu}+h_{\mu\nu}\) and

\(k^2=\eta^{\mu\nu}k_\mu k_\nu\), the spin-two part of the quadratic form

of \(R/(16\pi G_{\mathrm{eff}})+c_WC^2+c_RR^2\) is proportional to

\(k^2(1-32\pi G_{\mathrm{eff}}c_Wk^2)\) and its spin-zero part to

\(k^2(1+96\pi G_{\mathrm{eff}}c_Rk^2)\) (a symbolic check, not a proof).

There is therefore a spin-two pole at \(k^2=-m_2^2\) with residue

opposite in sign to that of the massless graviton, and a spin-zero pole

at \(k^2=-m_0^2\) with residue of the graviton's sign, where

\[

  m_2^2=-\frac{1}{32\pi G_{\mathrm{eff}}c_W}

       =-\frac{\Upsilon\LUV^2}{L},

  \qquad

  m_0^2=\frac{1}{96\pi G_{\mathrm{eff}}c_R},

  \qquad

  \Upsilon:=\frac{\pi}{G_{\mathrm{eff}}\mathcal W_C\LUV^2}.

\]

If \(G_{\mathrm{eff}}\mathcal W_C>0\), the extra spin-two pole of this

fixed-\(L\) form is a tachyonic ghost, with \(|m_2|<\LUV\) exactly when

\(L>\Upsilon\). If \(1/G_{\mathrm{eff}}\) is taken from the

universal \(a_1\) term alone, \(1/G_{\mathrm{eff}}

=\mathcal A_1(\LUV^2-\mu^2)/(2\pi)\) with

\(\mathcal A_1:=\sum_sw_s\mathcal A_{1,R,s}\), and

\(\Upsilon=\mathcal A_1/(2\mathcal W_C)\) up to \(O(\mu^2/\LUV^2)\),

independent of \(\LUV\). For sectors of minimally coupled real scalars,

\(\mathcal A_1=\mathcal N_0/6\) and \(\mathcal W_C=\mathcal N_0/120\), so

\(\Upsilon=10\) for every \(\mathcal N_0>0\),

\(|m_2|<\LUV\) in this fixed-\(L\) form once \(\LUV/\mu>e^{5}\), and

\(m_0^2=2\LUV^2/L\). Since \(|m_2|^2/\mu^2=\Upsilon e^{L}/L\), which is

much larger than one for \(\mu\ll\LUV\) unless \(\Upsilon\) is

exponentially small, this pole lies at momenta where the fixed-\(L\)

form does not apply (see below); it is a property of that form, not a

one-loop prediction. For \(|k|\lesssim\mu\), where the fixed-\(L\) form

applies, the Weyl correction relative to the Einstein term is

\(32\pi G_{\mathrm{eff}}c_Wk^2=(L/\Upsilon)\,k^2/\LUV^2\).



The fixed-\(L\) form is not the one-loop result at momenta

\(\mu\ll|k|\ll\LUV\). At second order in curvature the one-loop action

is nonlocal, with momentum-dependent form factors~\cite{BarvinskyVilkovisky1990}; assuming that

the Weyl form factor decays when \(sk^2\gg1\), its proper-time integral

is cut off at \(s\sim k^{-2}\) rather than at \(\mu^{-2}\), and for

massless sectors \(L\) is replaced by \(\ln(\LUV^2/k^2)+\delta_W\), with

a constant \(\delta_W\) fixed by the non-logarithmic part of the form

factor and by the finite Weyl counterterm. For the analogous

second-order form factor \(\int_0^1\mathrm d\tau\,

e^{-\tau(1-\tau)sk^2}\) of a potential insertion, a direct

numerical evaluation gives \(\ln(\LUV^2/k^2)+2-\gamma_{\mathrm E}\),

with \(\gamma_{\mathrm E}\) Euler's constant,

independent of \(\mu\) for \(k\gg\mu\); this illustrates the mechanism

and does not compute \(\delta_W\). With \(x:=k^2/\LUV^2\) and the

principal logarithm, the spin-two factor becomes

\[

  F(x)=1-\frac{x}{\Upsilon}\Bigl(\ln\frac1x+\delta_W\Bigr),

\]

whose zeros on the principal sheet are

\(x=-\Upsilon/W_j(-\Upsilon e^{-\delta_W})\), \(j\in\{0,-1\}\), with

\(W_j\) the branches of the Lambert function; \(\mu\) no longer enters.

At \(\delta_W=0\) the zeros cross \(|x|=1\) at \(x=\pm i\) when

\(\Upsilon=\pi/2\). A numerical scan with an argument-principle count

(a check, not a proof) finds that \(F\) has no zero in the disc

\(|k^2|<\LUV^2\) exactly when \(\Upsilon>\Upsilon_*(\delta_W)\), with

\(\Upsilon_*(0)=\pi/2\) and \(\Upsilon_*(\delta_W)=\delta_W\) for

\(\delta_W\ge1\), where the crossing is at \(x=1\). For \(\Upsilon=10\)

the condition is \(-9.58<\delta_W<10\); at \(\delta_W=0\) the zeros

are a complex-conjugate pair with \(|k^2|\approx3.9\LUV^2\).



These statements are conditional on the truncated identity, the

continuation and common limits stated, massless sectors (massive

sectors require the threshold functions of Problem~F), and the use of

the leading-logarithmic form factor up to \(|k^2|\sim\LUV^2\), where

the derivative expansion is not controlled. They do not compute

\(\mathcal W_C\), \(\mathcal A_1\), \(G_{\mathrm{eff}}\) or \(\delta_W\)

for the locked block, do not derive the induced coefficients from the

matrix path integral, and do not decide whether complex-conjugate zeros

beyond \(\LUV\) are propagating ghosts. Within the leading-logarithmic

model \(F\), absence of a zero in \(|k^2|<\LUV^2\) is a numerical

condition on \(\Upsilon\) and \(\delta_W\), which are outputs of

Problem~F, and is not implied by the metric being composite.

Problems~F and~N remain open.



